% Options for packages loaded elsewhere
\PassOptionsToPackage{unicode}{hyperref}
\PassOptionsToPackage{hyphens}{url}
\PassOptionsToPackage{dvipsnames,svgnames,x11names}{xcolor}
%
\documentclass[
  11pt,
  a4paper,
]{article}
\usepackage{amsmath,amssymb}
\usepackage{iftex}
\ifPDFTeX
  \usepackage[T1]{fontenc}
  \usepackage[utf8]{inputenc}
  \usepackage{textcomp} % provide euro and other symbols
\else % if luatex or xetex
  \usepackage{unicode-math} % this also loads fontspec
  \defaultfontfeatures{Scale=MatchLowercase}
  \defaultfontfeatures[\rmfamily]{Ligatures=TeX,Scale=1}
\fi
\usepackage{lmodern}
\ifPDFTeX\else
  % xetex/luatex font selection
\fi
% Use upquote if available, for straight quotes in verbatim environments
\IfFileExists{upquote.sty}{\usepackage{upquote}}{}
\IfFileExists{microtype.sty}{% use microtype if available
  \usepackage[]{microtype}
  \UseMicrotypeSet[protrusion]{basicmath} % disable protrusion for tt fonts
}{}
\makeatletter
\@ifundefined{KOMAClassName}{% if non-KOMA class
  \IfFileExists{parskip.sty}{%
    \usepackage{parskip}
  }{% else
    \setlength{\parindent}{0pt}
    \setlength{\parskip}{6pt plus 2pt minus 1pt}}
}{% if KOMA class
  \KOMAoptions{parskip=half}}
\makeatother
\usepackage{xcolor}
\usepackage[margin=27mm]{geometry}
\usepackage{longtable,booktabs,array}
\usepackage{calc} % for calculating minipage widths
% Correct order of tables after \paragraph or \subparagraph
\usepackage{etoolbox}
\makeatletter
\patchcmd\longtable{\par}{\if@noskipsec\mbox{}\fi\par}{}{}
\makeatother
% Allow footnotes in longtable head/foot
\IfFileExists{footnotehyper.sty}{\usepackage{footnotehyper}}{\usepackage{footnote}}
\makesavenoteenv{longtable}
\setlength{\emergencystretch}{3em} % prevent overfull lines
\providecommand{\tightlist}{%
  \setlength{\itemsep}{0pt}\setlength{\parskip}{0pt}}
\setcounter{secnumdepth}{-\maxdimen} % remove section numbering
\usepackage{amsmath,amssymb,amsthm,mathtools}
\usepackage{microtype,booktabs,longtable,array,xurl}
\usepackage{fancyhdr}
\pagestyle{fancy}\fancyhf{}
\fancyhead[L]{\small Calibration along gradient descent}
\fancyhead[R]{\small Y. Shmalo}
\fancyfoot[C]{\thepage}
\setlength{\headheight}{14pt}
\setlength{\emergencystretch}{3em}
\allowdisplaybreaks[2]
\providecommand{\sech}{\operatorname{sech}}
\providecommand{\argmin}{\operatorname*{arg\,min}}
\providecommand{\argmax}{\operatorname*{arg\,max}}
\providecommand{\logit}{\operatorname{logit}}
\providecommand{\diag}{\operatorname{diag}}
\ifLuaTeX
  \usepackage{selnolig}  % disable illegal ligatures
\fi
\IfFileExists{bookmark.sty}{\usepackage{bookmark}}{\usepackage{hyperref}}
\IfFileExists{xurl.sty}{\usepackage{xurl}}{} % add URL line breaks if available
\urlstyle{same}
\hypersetup{
  pdftitle={Calibration along gradient descent in Gaussian tanh random-feature classifiers},
  pdfauthor={Yitzchak Shmalo},
  colorlinks=true,
  linkcolor={blue},
  filecolor={Maroon},
  citecolor={Blue},
  urlcolor={blue},
  pdfcreator={LaTeX via pandoc}}

\title{Calibration along gradient descent in Gaussian tanh
random-feature classifiers}
\author{Yitzchak Shmalo}
\date{31 October 2026}

\begin{document}
\maketitle

\begin{center}
\small Einstein Institute of Mathematics, The Hebrew University of Jerusalem\\
Givat Ram, Jerusalem, Israel
\end{center}

\begin{abstract}
We study ordinary full-batch gradient descent on cross-entropy in a
two-layer Gaussian random-feature model with tanh activation and a
logistic single-index teacher. At every fixed positive iteration and
arbitrary positive finite sample/input and width/input ratios, an explicit
causal recursion determines the asymptotic population risk and exact
unbinned calibration conditional on the trained model. The proof verifies
matrix universality for the actual nonlinear features and upgrades joint
Gaussian approximation to total variation. In an explicit sample-rich
regime, the calibration approximation is uniform over all positive times
and the limiting estimator. In a separate strongly overparameterized
regime, a large-norm witness gives a uniform calibration floor at every
positive stopping time and a finite window for population-risk-optimal
stopping. A scalar single-crossing argument proves strict average
overconfidence at the risk-optimal first-step learning scale. All-time
claims use stated sufficient ratio and step conditions; hidden features
remain fixed throughout training.
\end{abstract}

\clearpage
\setcounter{tocdepth}{1}
\tableofcontents
\clearpage

\hypertarget{introduction-and-relation-to-prior-work}{%
\section{Introduction and relation to prior
work}\label{introduction-and-relation-to-prior-work}}

Calibration compares a predictor's stated probability with the label
frequency conditional on that prediction. It is distinct from
classification accuracy and from recovering the entire conditional
probability given the input. This paper concerns exact score-conditional
calibration for the realized trained network. The training matrix is
reused at each update, and the algorithm is ordinary gradient descent.

Bai, Mei, Wang and Xiong {[}1{]} established high-dimensional
overconfidence for logistic regression, including an underparameterized
setting. Clarté, Loureiro, Krzakala and Zdeborová {[}2{]} already
derived calibration and uncertainty descriptions for random-feature
estimators. Thus neither random-feature calibration nor the Gaussian
conditioning formula used below is a new subject or a new general
formula. The contribution here is the discrete training
trajectory of ordinary reused-data GD, exact model-conditional
calibration convergence, and stopping-time conclusions under explicit
assumptions.

Early-stopping benefits and calibration tradeoffs have also been studied
in other scalings and models {[}6--8{]}. Our results keep proportional
aspect ratios and a positive learning scale fixed. They do not rule out
calibrated early stopping under different sample scalings, vanishing
learning scales, feature learning, or post-training calibration.

The matrix spectral input is due to Louart, Liao and Couillet {[}3{]};
the rectangular memory and covariance formulas and their universality
are due to Fan {[}4{]} and Wang, Zhong and Fan {[}5{]}. The appendix
verifies their hypotheses for the present matrix and gives an exact
auxiliary representation whose corrections cancel in the physical GD
update. Theorem A covers every fixed finite horizon. Theorems B and D
use different nonempty ratio regimes to control all positive stopping
times. Theorem C concerns signed average overconfidence, which is
stronger than positive calibration error and requires its stated choice
of scale.

This is an internally reviewed research manuscript. Its precise
contributions are presented without a claim of external acceptance or
exclusive publication priority.

\hypertarget{model-algorithm-and-calibration-definition}{%
\section{1. Model, algorithm, and calibration
definition}\label{model-algorithm-and-calibration-definition}}

Let \(d,n,m\to\infty\), with \[
n/d\to\alpha\in(0,\infty),\qquad m/d\to\gamma\in(0,\infty).
\] Fix a deterministic unit vector \(\theta_d\), finite \(\beta\ge0\),
and independent \[
X_i\sim N(0,I_d),\qquad W_j\sim N(0,I_d/d),
\qquad Y_i\mid X_i\sim\operatorname{Bernoulli}(\sigma(\beta\theta_d^\top X_i)),
\] where the labels are conditionally independent and
\(\sigma(u)=(1+e^{-u})^{-1}\). The two-layer score and probability are
\[
S_b(x)=\sum_{j=1}^m b_j\tanh(W_j^\top x),\qquad p_b(x)=\sigma(S_b(x)).
\] Only the output coefficients are trained. Write
\(H_{ij}=\tanh(W_j^\top X_i)\), and define mean cross-entropy \[
\mathcal L_n(b)=\frac1n\sum_i[\log(1+e^{(Hb)_i})-Y_i(Hb)_i].
\] For fixed \(c>0\), the \textbf{actual algorithm} is \[
b^0=0,\qquad
b^{t+1}=b^t+\frac{cd}{nm}H^\top[Y-\sigma(Hb^t)]. \tag{1}
\] This is mean-loss GD with step \(cd/m\) in \(b\), equivalently step
\(cd\) in the normalized output coefficients \(w=\sqrt m\,b\). The same
sample and feature matrix are reused at every step.

Let \(\mathcal E_d\) contain all training data and features. For a fresh
independent input, put \(Z=\theta_d^\top X\) and \(S_{t,d}=S_{b^t}(X)\).
Define \[
A_{t,d}=\mathbb E[ZS_{t,d}\mid\mathcal E_d],\qquad
Q_{t,d}=\mathbb E[S_{t,d}^2\mid\mathcal E_d].
\] The exact, unbinned, model-conditional calibration error is \[
\mathrm{CE}_{t,d}=\mathbb E\big[\,|\mathbb E[Y\mid S_{t,d},\mathcal E_d]-\sigma(S_{t,d})|\mid\mathcal E_d\big]. \tag{2}
\] Conditioning on the score or its sigmoid is equivalent. This is not a
fixed-bin estimator or an average of calibration computed under
resampled models.

Throughout, \(\|\mu-\nu\|_{\rm TV}=\sup_{|f|\le1}|\int f\,d(\mu-\nu)|\),
the full total-variation norm, twice the probability-distance
convention.

For a centered Gaussian pair with \(\operatorname{Var}Z=1\), covariance
\(A\), and score variance \(Q>A^2\), set \[
\begin{aligned}
g_{A,Q}(s)&=\mathbb E_V\sigma\left(\beta\left[\frac AQs+\sqrt{1-A^2/Q}\,V\right]\right),\\
C_\beta(A,Q)&=\mathbb E_U\left|g_{A,Q}(\sqrt Q\,U)-\sigma(\sqrt Q\,U)\right|,\\
\mathcal R_\beta(A,Q)&=\mathbb E_U\log\left(2\cosh(\sqrt Q\,U/2)\right)-\mu_\beta A,\\
\mu_\beta&=\mathbb E[U(\sigma(\beta U)-1/2)],
\end{aligned} \tag{3}
\] where \(U,V\) are independent standard normals.

\hypertarget{finite-time-theorem}{%
\section{2. Finite-time theorem}\label{finite-time-theorem}}

\textbf{Theorem A.} Under the model above, for every fixed integer
\(T\ge1\) and fixed \(c>0\), the explicit recursion in Section 3 defines
deterministic \(A_t,Q_t\) and a positive-definite matrix \(\Sigma_T\),
with \[
Q_t-A_t^2=(\Sigma_T)_{tt}>0,\qquad 1\le t\le T.
\] For each such time, \[
A_{t,d}\xrightarrow{P}A_t,\quad Q_{t,d}\xrightarrow{P}Q_t,\quad
\left\|\mathcal L((Z,S_{t,d})\mid\mathcal E_d)-
N\left(0,\begin{pmatrix}1&A_t\\A_t&Q_t\end{pmatrix}\right)\right\|_{\rm TV}
\xrightarrow{P}0. \tag{4}
\] The entire finite test-score path converges conditionally weakly in
probability to the Gaussian path in (10) below. In particular, \[
\mathrm{CE}_{t,d}\xrightarrow{P}C_\beta(A_t,Q_t)>0,\qquad
\mathbb E[\ell(Y,S_{t,d})\mid\mathcal E_d]\xrightarrow{P}\mathcal R_\beta(A_t,Q_t). \tag{5}
\] The TV and calibration convergences also hold in every finite \(L^p\)
over the training environment, since their errors are uniformly bounded.
The calibration curves satisfy the stronger integrated statement \[
\mathbb E\left[\left|\mathbb E[Y\mid S_{t,d},\mathcal E_d]-g_{A_t,Q_t}(S_{t,d})\right|\mid\mathcal E_d\right]\xrightarrow{P}0. \tag{6}
\] All finite aspect ratios are allowed, including \(m>n\). The theorem
fixes time before taking the dimension limit; it alone does not justify
a growing-time conclusion.

\hypertarget{explicit-recursion-the-dependence-on-sample-ratio-width-and-time}{%
\section{3. Explicit recursion: the dependence on sample ratio, width,
and
time}\label{explicit-recursion-the-dependence-on-sample-ratio-width-and-time}}

Set \[
a=\mathbb E[U\tanh U],\quad v=\mathbb E[\tanh^2U],\quad
\nu=v-a^2>0,\quad \lambda=c/\alpha,\quad s=a\sqrt\gamma.
\] There are two scalar probability spaces. On the training-row space,
\(Z\sim N(0,1)\) and
\(Y\mid Z\sim\operatorname{Bernoulli}(\sigma(\beta Z))\). Gaussian row
fields \(G_1,G_2,\ldots\) are independent of \((Z,Y)\). On the column
space, standard normal \(H_\star\) is independent of Gaussian column
fields. Both Gaussian covariance sequences are constructed below.

Initialize the training residual function \(R_1=Y-1/2\), the column
coefficient \(q_0=0\), and empty matrices \(\Psi_0,\Gamma_0,\Sigma_0\).
At stage \(t\), the previous row covariance determines the joint law of
\((G_1,\ldots,G_{t-1})\) and hence the already defined functions
\(R_1,\ldots,R_t\). Compute \[
\widehat\Delta_t=\alpha(\mathbb E[R_iR_j])_{i,j\le t},\quad
(\widehat\Phi_t)_{ij}=\alpha\mathbb E[\partial_{G_j}R_i]\ (j<i),\quad
\widehat\ell_t=\alpha\mathbb E[ZR_t], \tag{7}
\] with all other entries of \(\widehat\Phi_t\) zero. Embed the
preceding \(\Psi,\Gamma\) in \(t\times t\) matrices with zero last row
and column. First calculate \[
\begin{aligned}
U_t&=(I-a^2\widehat\Phi_t\Psi_t)^{-1},\quad T_t=(I-a^2\Psi_t\widehat\Phi_t)^{-1},\\
\mathsf B_t&=\gamma^{-1}(\nu\widehat\Phi_t+a^2U_t\widehat\Phi_t),\\
\Omega_t&=\gamma^{-1}\left[\nu\widehat\Delta_t+
a^2U_t(\widehat\Delta_t+a^2\widehat\Phi_t\Gamma_t\widehat\Phi_t^\top)U_t^\top\right].
\end{aligned} \tag{8}
\] The last column of \(\widehat\Phi_t\) is zero, so these expressions
use only previous column information. If \(\psi_j\) denotes row \(j\) of
\(\Psi\), padded with zeros, complete the current row and coefficient by
\[
\psi_t=\psi_{t-1}+\lambda\left[e_t^\top+\sum_{j<t}(\mathsf B_t)_{tj}\psi_j\right],
\quad q_t=q_{t-1}+\lambda\left[\sum_{j<t}(\mathsf B_t)_{tj}q_j+
\frac{s}{\gamma}\widehat\ell_t\right].
\] Here \(\psi_0=0\). Put \(q_{1:t}=(q_1,\ldots,q_t)^\top\) and \[
\begin{aligned}
\Gamma_t&=\Psi_t\Omega_t\Psi_t^\top+q_{1:t}q_{1:t}^\top,\\
\mathsf A_t&=\nu\Psi_t+a^2T_t\Psi_t,\\
\Sigma_t&=\nu\Gamma_t+a^2T_t
(\Gamma_t+a^2\Psi_t\widehat\Delta_t\Psi_t^\top)T_t^\top,
\end{aligned} \tag{9}
\] where \(T_t\) is now evaluated with the completed \(\Psi_t\). Its
inverses exist without a small-step assumption: the products with
\(\widehat\Phi_t\) are strictly lower triangular, so their inverses are
finite geometric polynomials.

Let \(G_{1:t}\sim N(0,\Sigma_t)\), consistently with previous stages and
independently of \((Z,Y)\), and set \[
S_t^{\rm train}=G_t+\sum_{j\le t}(\mathsf A_t)_{tj}R_j+s q_tZ,
\qquad R_{t+1}=Y-\sigma(S_t^{\rm train}).
\] Derivatives required in (7) are computed recursively, with
deterministic coefficients held fixed: \[
\partial_{G_j}R_{t+1}=-\sigma'(S_t^{\rm train})
\left[\mathbf1_{j=t}+\sum_{k\le t}(\mathsf A_t)_{tk}\partial_{G_j}R_k\right].
\] This finishes one causal stage. Only finite matrices and Gaussian
expectations occur. The fresh-test law and required overlaps are \[
S_t^{\rm test}=s q_tZ+G_t,
\quad A_t=s q_t,
\quad Q_t=s^2q_t^2+(\Sigma_t)_{tt}. \tag{10}
\] Thus (3), (7)--(10) explicitly characterize calibration and risk as
functions of \(\alpha,\gamma,\beta,c,t\).

As an exact normalization check, at the first update \[
A_1=ca^2\mu_\beta,\qquad
Q_1-A_1^2=c^2\tau^2,\quad
\tau^2=\frac{\mu_\beta^2a^2v}{\gamma}+
\frac{a^4+v^2/\gamma}{4\alpha}. \tag{11}
\] This agrees with the direct one-step calculation in Appendix B. The
two variance terms both matter.

\hypertarget{uniform-training-time-and-early-stopping-theorem}{%
\section{4. Uniform training-time and early-stopping
theorem}\label{uniform-training-time-and-early-stopping-theorem}}

We give explicit sufficient constants, without claiming they are sharp.
Let \[
L=8\sqrt{1+\gamma},\quad s_0=\nu/2,\quad k_0=\sqrt{s_0/2},\quad
p_0=s_0^2/(20L^4),
\] \[
\begin{gathered}
B_0=\sqrt{2\log(4/p_0)},\quad p_1=\sigma(-\beta B_0)p_0/2,\quad
\mathfrak a=k_0p_1/4,\\
C_* = \max\{1,(32L/(k_0p_1))^2\},\qquad \Lambda=16L^2.
\end{gathered}
\] Assume additionally \[
\alpha/\gamma>C_*,\qquad 0<c<\gamma/(2\Lambda). \tag{12}
\] These are finite constants for every fixed finite \(\beta\) and
\(\gamma>0\), so the regime is nonempty. It has more samples than
trained output coefficients.

\textbf{Theorem B.} Under (12), with probability tending to one, the
empirical logistic loss has a unique finite MLE \(b^*\), and GD
converges to it geometrically with a dimension-independent contraction
factor. Use \(t=\infty\) for this MLE. The deterministic trajectory from
Section 3 has limits \(A_\infty,Q_\infty\), and \[
\sup_{t\in\{1,2,\ldots,\infty\}}
\{|A_{t,d}-A_t|+|Q_{t,d}-Q_t|\}\xrightarrow{P}0,
\] \[
\sup_{t\in\{1,2,\ldots,\infty\}}
|\mathrm{CE}_{t,d}-C_\beta(A_t,Q_t)|\xrightarrow{P}0. \tag{13}
\] There is a deterministic \(\varepsilon_*>0\), depending only on the
fixed parameters, such that \[
\mathbb P\left\{\inf_{t\in\{1,2,\ldots,\infty\}}\mathrm{CE}_{t,d}\ge\varepsilon_*/2\right\}\longrightarrow1. \tag{14}
\] Consequently no training-measurable stopping rule with a positive
integer value, or the value \(\infty\), attains asymptotic calibration
in this regime. A stopping rule may grow arbitrarily with dimension. If
it diverges in probability, its calibration tends to
\(C_\beta(A_\infty,Q_\infty)>0\). This conclusion keeps \(c\) fixed and
positive; it does not cover dimension-dependent vanishing steps or all
late-time overparameterized regimes.

At time zero, the constant prediction \(1/2\) is calibrated by teacher
symmetry. Equation (14) concerns positive training times and nontrivial
learning. It does not confuse calibration with Bayes accuracy.

\hypertarget{strict-overconfidence-with-an-improved-predictor}{%
\section{5. Strict overconfidence with an improved
predictor}\label{strict-overconfidence-with-an-improved-predictor}}

\textbf{Theorem C.} Let \(\beta>0\) and any fixed \(\alpha,\gamma>0\).
Put \(k=a^2\mu_\beta>0\), let \(\tau>0\) be given by (11), and set
\(U=kZ+\tau G\) with independent standard normals. Define
\(g(u)=\mathbb E[\sigma(\beta Z)\mid U=u]\) and \[
R(c)=\mathbb E[\log(1+e^{cU})-cU g(U)].
\] There is a unique population-risk-optimal first-step scale \(c_*>0\).
For GD (1) run with this fixed scale, its first-step population risk and
signed overconfidence satisfy \[
\mathcal R_{1,d}\xrightarrow{P}R(c_*)<\log2,
\] \[
\mathrm{OC}_{1,d}:=
\mathbb E[(\sigma(S_{1,d})-\sigma(\beta Z))\operatorname{sign}(S_{1,d})\mid\mathcal E_d]
\xrightarrow{P}\mathrm{OC}_*>0. \tag{15}
\] The signed quantity is \textbf{average predicted confidence minus
classification accuracy}. It does not assert pointwise overconfidence at
every score. In fact the optimal-scale Gaussian predictor is
overconfident close to its decision boundary and underconfident in the
far tails.

This theorem uses Theorem A with \(c=c_*\). The scale need not
automatically lie in (12); the following corollary gives explicit
sufficient conditions for that compatibility. Optimizing first-step risk
is not asserted to optimize over the entire GD path.

\textbf{Corollary: one fixed model can satisfy all three conclusions.}
The regimes are not disjoint. Retain \(L,k_0,p_0,B_0\) above, and choose
\[
\alpha/\gamma>C_0:=\max\{1,(256L/(k_0p_0))^2\},\qquad
0<\beta<\min\{\log(3)/B_0,\ v/(8aL\sqrt\alpha)\}.
\tag{16}
\] Then \(\sigma(-\beta B_0)>1/4\), so \(C_*(\beta)\le C_0\). Further,
\[
c_*<\frac{\beta k}{k^2+\tau^2}
\le\frac{a^2\alpha\gamma\beta^2}{v^2}
<\frac{\gamma}{64L^2}<\frac{\gamma}{2\Lambda}.
\] Here \(\mu_\beta\le\beta/4\) and \(\tau^2\ge v^2/(4\alpha\gamma)\).
Thus the \textbf{same} fixed-step network, with \(c=c_*\), has the
complete deterministic training-time calibration curve, an improved and
strictly overconfident first iterate, and the uniform positive
calibration floor for every positive stopping time. Equation (16) is an
explicit nonempty open regime with a weak but nonzero teacher. It does
not impose this restriction on Theorem A or on Theorem C separately.

\hypertarget{all-positive-stopping-times-in-a-strongly-overparameterized-regime}{%
\section{6. All positive stopping times in a strongly overparameterized
regime}\label{all-positive-stopping-times-in-a-strongly-overparameterized-regime}}

There is a second all-time regime, with more trained coefficients than
samples. Its proof does not require a late-time Gaussian approximation.
Define, with \(\nu\) as in Section 3, \[
s_X=\nu/2,\quad \kappa_X=\sqrt{s_X/2},\quad
L_X=8\sqrt{1+\alpha},\quad p_X=s_X^2/(20L_X^4),
\] \[
C_X=\max\{1,(32L_X/(\kappa_Xp_X))^2\},\qquad
k_X=\kappa_X^2p_X/16,\quad K_X=64L_X^2.
\] Assume \[
\gamma/\alpha>C_X,\qquad 0<c<2\alpha/K_X. \tag{17}
\] For every fixed \(\alpha>0\) and finite \(\beta\), this is a nonempty
proportional regime. The constants are conservative sufficient
conditions.

\textbf{Theorem D.} Under (17), there is \(\varepsilon>0\) such that \[
\mathbb P\{\inf_{t\in\{1,2,\ldots\}}\mathrm{CE}_{t,d}
\ge\varepsilon\}\longrightarrow1. \tag{18}
\] Thus every training-measurable stopping rule taking a finite positive
integer value, including one diverging arbitrarily with dimension, has a
positive exact-calibration floor. The empirical loss converges to zero,
the coefficient norms diverge, and population cross-entropy diverges
uniformly at late times: for every \(M<\infty\), there is a
deterministic \(T_M\) with \[
\mathbb P\{\inf_{t\ge T_M}\mathcal R_{t,d}>M\}\longrightarrow1.
\tag{19}
\] There is a deterministic finite \(T_0\) such that, with probability
tending to one, every population-risk minimizing time on the entire GD
path lies in \(\{0,\ldots,T_0\}\). Theorem A therefore determines the
limiting optimal risk by minimizing over this finite window. If the
limiting minimizer is unique, the minimizing time converges to it. For
\(\beta>0\), a sufficiently small fixed positive \(c\) gives a
first-step improvement over the constant predictor while retaining the
calibration obstruction.

\textbf{Proof.} Appendix F proves that \(HH^\top/m\) lies between
\(k_XI\) and \(K_XI\) with probability tending to one. A separating
comparator and the convex GD inequality give a dimension-uniform loss
upper bound tending to zero. Jensen's inequality then forces every late
iterate's norm above any specified constant. A uniform small-ball event
for fresh features and finite teacher noise gives a score-measurable
witness to calibration error for every sufficiently large coefficient
vector, whether delocalized or not. Combining that witness with Theorem
A for the finitely many earlier times proves (18). A related witness
lower bounds population cross-entropy by a positive multiple of
coefficient norm, proving (19) and the finite optimization window. The
full inequalities, constants, and quantifiers are in Appendix F.

Time \(\infty\) is not a finite MLE in this separable regime, and no
claim is made here about the calibration of a terminal hard classifier.
Theorem D supplies a uniform lower bound, not an exact late-time
calibration curve. The exact curve in Theorem A fixes time before
dimension. Theorem B gives a uniform curve under its separate
assumptions.

\hypertarget{proof-of-the-combined-theorem}{%
\section{7. Proof of the combined
theorem}\label{proof-of-the-combined-theorem}}

The detailed proofs are in Appendices A--F. The following proof
establishes their logical composition and identifies the imported
mathematical results.

\textbf{Step 1: actual random-feature matrix and independent teacher
signal.} Rotate \(\theta=e_1\), and write input and feature vectors as
their first coordinate plus perpendicular coordinates. Taylor expansion
of tanh gives, for training rows and a proportional independent test
block together, \[
H/\sqrt m=A_0+(a\sqrt\gamma/m)zh^\top+E_d,
\qquad \|E_d\|_{\rm op}\xrightarrow{P}0.
\] The perpendicular matrix \(A_0\) is independent of \(z,h\), labels,
and the training/test mask. The error follows from Gaussian sup-norm
bounds, a centered derivative-feature matrix bound, and a Frobenius
bound on the Taylor remainder; no entrywise-to-operator inference is
omitted.

\textbf{Step 2: the universality hypotheses are proved for this matrix.}
Clip each underlying Gaussian matrix in operator norm and then clip the
normalized feature matrix. These Frobenius-nonexpansive projections
preserve signed-permutation symmetry and make each fixed matrix-power
entry globally \(O(d^{-1/2})\)-Lipschitz in all Gaussian entries.
Gaussian concentration gives summable failure bounds for every fixed
power: off-diagonal entries vanish at the required rate, and diagonal
entries approach their normalized trace uniformly. The clipped and
original matrices agree eventually almost surely. Conditional Hermite
covariance analysis and the Louart--Liao--Couillet spectral theorem
identify the limiting singular law with the biorthogonally invariant
Gaussian comparison matrix. These facts verify Wang--Zhong--Fan
Proposition D.1(b2), including more than just an empirical spectral
distribution.

The row-normalized rectangular cumulants are \[
\kappa_2=v,\qquad \kappa_{2j}=a^{2j}\delta^{j-1}\quad(j\ge2),
\] where \(\delta\) is the total training-plus-query row ratio before
the query block is eliminated. Substitution into Fan's rectangular
covariance and memory formulas sums finite triangular series. After
conditioning on training rows, every \(\delta\) cancels exactly,
yielding (8)--(9).

\textbf{Step 3: ordinary GD is represented causally.} Write
\(w^t=\sqrt m\,b^t\) for the normalized coefficient vector. Introduce
auxiliary column fields by subtracting the deterministic memory terms
\(\mathsf B\) from \(A_0^\top\mathbf R_t\), and row fields by
subtracting \(\mathsf A\) from \(A_0w^t\), where \(\mathbf R_t\) is the
sample residual vector. Define the coordinatewise nonlinearities by
adding exactly these terms back. They cancel algebraically in the
physical update. Teacher overlaps are first fixed to their deterministic
recursive expectations. The resulting history-dependent maps meet the
continuity, growth, Lipschitz, and independent-side-information
assumptions of Wang--Zhong--Fan Theorem 2.22.

No degeneracy assumption is left open. The residual second-moment matrix
begins with \(\Delta_1>0\). From (8)--(9), \(\Omega_t\) dominates a
positive multiple of \(\Delta_t\), and \(\Sigma_t\succeq\nu\Gamma_t\).
The derivative matrix \(\Psi_t\) has positive diagonal \(\lambda\), so
\(\Gamma_t\succ0\). Conditional on the past and a training row, the next
sigmoid residual depends strictly on a fresh Gaussian innovation; hence
it cannot lie in the linear span of preceding residuals. This proves
positive definiteness inductively for every fixed time and every
\(c>0\).

The AMP theorem now gives joint empirical state convergence. Its overlap
convergence, a normalized Euclidean stability bound, and a finite
induction replace frozen teacher overlaps by the actual empirical ones.
The same stability argument removes \(E_d\). On query rows all residuals
are zero, so the limiting physical score is exactly (10). Conditional
laws of large numbers and conditional fourth-moment bounds identify the
realized-network covariances. This proves finite-time conditional weak
convergence and the deterministic overlaps for the original nonlinear
GD, rather than a different AMP training algorithm.

\textbf{Step 4: exact calibration requires, and receives, total
variation.} A leave-one-feature comparison proves \(\|b^t\|_2=O_P(1)\)
and \(\|b^t\|_\infty=O_P(\sqrt{\log d/d})\) at any fixed horizon. Thus
the fresh score has bounded gradient and Hessian Hilbert--Schmidt norm,
while its Hessian operator norm vanishes. An Ornstein--Uhlenbeck
covariance identity proves Gaussian weak approximation with matched
covariance. The teacher-orthogonal derivative Gram determinant is
concentrated away from zero whenever the residual score variance is
positive. Localized inverse-vector fields and Gaussian integration by
parts give a dimension-free total-variation smoothing bound.
Mollification then upgrades weak convergence to joint TV. This analytic
argument is proved in Appendix C.

For any joint law \(\mu\) of \((Z,S)\), define the signed measure \[
\zeta_\mu(B)=\int\mathbf1_{s\in B}[\sigma(\beta z)-\sigma(s)]\,\mu(dz,ds).
\] Its TV norm equals calibration error. Therefore
\(|C(\mu)-C(\widetilde\mu)|\le\|\mu-\widetilde\mu\|_{\rm TV}\). The same
measure argument controls the integrated conditional regression curve.
Logistic loss is at most \(\log2+|S|\), so bounded conditional second
moments justify truncation for risk. These observations prove (4)--(6).

Strict positivity in (5) is elementary. For \(A\le0\) or \(\beta=0\),
the conditional teacher curve cannot equal a strictly increasing
nonconstant sigmoid. For \(A>0,\beta>0\), write it as
\(g(s)=\mathbb E\sigma(ks+rG)\) with \(r>0\). Its upper tail satisfies
\(e^{ks}(1-g(s))\to e^{r^2/2}\). Equality with \(\sigma(s)\) would first
require \(k=1\), then \(r=0\), a contradiction.

\textbf{Step 5: all times are controlled on one event.} Under (12), the
conditional feature covariance and its teacher-orthogonal residual both
dominate \((\nu/2)I\). A bounded-cutoff small-ball argument and the
noisy labels yield empirical coercivity uniformly in every coefficient
direction. Uniform unsigned small-ball bounds give a sample covariance
lower bound even after deleting an arbitrary small fraction of rows. On
the bounded initial loss sublevel set, this yields uniform positive
logistic Hessian curvature. GD and every feature-deleted GD path stay in
that same ball and contract by a common dimension-independent factor
\(q<1\).

Feature-deleted minimizers are defined on a fixed compact ball even on
exceptional data, preserving independence from the omitted feature.
Their missing-coordinate gradients have centered Gaussian concentration.
A union bound only through a deterministic \(O(\log d)\) horizon,
followed by the common geometric tail, proves \[
\sup_{t\ge0}\|b^t\|_\infty=O_P(\sqrt{\log d/d}).
\] The first update makes a nonvanishing loss decrease, so convexity
bounds every positive iterate's norm away from zero. This gives uniform
lower residual variance and upper score variance. The TV theorem is
sequentially uniform on this bounded class. Its calibration functional
is continuous and strictly positive on the resulting compact covariance
set, proving (14).

Finally, the geometric coefficient tail gives a geometric tail for
\(A_{t,d},Q_{t,d}\). Fixed-time deterministic limits from Theorem A are
therefore exponentially Cauchy. Splitting a supremum into finitely many
times and the common tail proves (13). This supplies the needed exchange
of training-time and dimension limits rather than assuming it.

\textbf{Step 6: signed overconfidence at the best first-step scale.} For
\(U=kZ+\tau G\), Gaussian conditioning gives
\(g(u)=\mathbb E\sigma(\lambda_0u+rG)\), where
\(\lambda_0=\beta k/(k^2+\tau^2)>0\) and \(r>0\). Appendix E proves that
\(\operatorname{logit}g(u)/u\) is strictly increasing on \(u>0\), from
\(4g'(0)\) to \(\lambda_0\). Its proof derives a normal variance-mixture
representation of logistic noise from its characteristic function, then
proves the necessary strict single-crossing property by an explicit
kernel-ratio inequality.

The cross-entropy has strictly positive second derivative and a unique
minimizer \(c_*\in(4g'(0),\lambda_0)\). The error
\(e(u)=\sigma(c_*u)-g(u)\) is positive below one threshold \(u_*\) and
negative above it. Risk stationarity gives \(\mathbb E[Ue(U)]=0\).
Consequently \[
\mathrm{OC}_*=2\int_0^\infty e(u)(1-u/u_*)\varphi_{k^2+\tau^2}(u)\,du>0.
\] Both factors have the same strict sign. Also \(R'(0)<0\), so
\(R(c_*)<\log2\). Theorem A's TV and risk convergence transfer both
strict conclusions to the actual trained network, proving (15). This
proves Theorems A--C. Appendix F proves Theorem D by a separate
large-norm argument, as described in Section 6.

\clearpage

\hypertarget{appendix-a.-feature-matrices-and-the-finite-time-gd-representation}{%
\section{Appendix A. Feature matrices and the finite-time GD
representation}\label{appendix-a.-feature-matrices-and-the-finite-time-gd-representation}}

Theorem, lemma, and proposition numbers in this appendix are local;
equation and section references carry the appendix letter.

\hypertarget{a.1-model-and-normalization}{%
\subsection{A.1 Model and
normalization}\label{a.1-model-and-normalization}}

Use the Gaussian model of Section 1: inputs \(X_i\sim N(0,I_d)\), a
deterministic unit teacher \(\theta\), conditionally independent labels
\(Y_i\mid X_i\sim\mathrm{Bernoulli}(\sigma(\beta\theta^\top X_i))\),
fixed finite \(\beta\ge0\), and independent feature rows
\(W_j\sim N(0,I_d/d)\). Let \(n/d\to\alpha>0,m/d\to\gamma>0\),
\(\phi=\tanh\), \(H_{ij}=\phi(W_j^\top X_i)\), and \(K=HH^\top/m\).

Write the score as \(S_t(x)=(b^t)^\top h(x)\), where \(b^t=a^t/\sqrt m\)
in the old network convention. Mean-loss GD with output-weight step
\(cd\) is exactly

\[
b^0=0,\quad u^t=Hb^t,\quad r^t=Y-\sigma(u^t),\quad
b^{t+1}=b^t+\frac{\lambda_d}{m}H^\top r^t,
\qquad\lambda_d=\frac{cd}{n}\longrightarrow\frac c\alpha. \tag{A.1.1}
\]

Consequently \(u^{t+1}=u^t+\lambda_dKr^t\). All statements below allow
any fixed \(c>0\), except the explicit small-step conclusion in
Proposition 3. The horizon \(T\) is a fixed integer independent of
dimension. No claim concerns \(T_d\to\infty\).

Define \[
a=\mathbb E[G\phi(G)],\quad v=\mathbb E\phi(G)^2,\quad
\nu=v-a^2>0,
\quad A(q)=\mathbb E\phi'(\sqrt qG).
\tag{A.1.2}
\] Strict positivity of \(\nu\) follows because tanh is not a scalar
multiple of its input in Gaussian \(L^2\).

\hypertarget{a.2-uniform-feature-matrix-bounds}{%
\subsection{A.2 Uniform feature-matrix
bounds}\label{a.2-uniform-feature-matrix-bounds}}

Conditional on \(F\), whose rows are \(W_j^\top\), the rows of \(H\) are
independent centered random vectors. For any unit \(v_0\in\mathbb R^m\),
the function \(x\mapsto v_0^\top\phi(Fx)\) is
\(\|F\|_{\rm op}\)-Lipschitz and odd. Gaussian concentration therefore
gives \[
\mathbb E[\exp(sv_0^\top h(X))\mid F]
\le \exp(s^2\|F\|_{\rm op}^2/2).
\tag{A.2.1}
\] For a unit \(u_0\in\mathbb R^n\), conditional independence gives the
same bound for \(u_0^\top Hv_0\). Taking \(1/4\)-nets of the two unit
spheres and using the standard net inequality proves \[
\|H\|_{\rm op}=O_P(\sqrt d),\qquad
\|K\|_{\rm op}=O_P(1). \tag{A.2.2}
\] The failure probabilities at sufficiently large fixed constants decay
exponentially in \(d\). The same net argument for the Gaussian matrices
proves \(\|F\|_{\rm op}=O_P(1)\) and \(\|X\|_{\rm op}=O_P(\sqrt d)\),
again with exponential tails. Integrating these bounds shows that all
fixed moments of the corresponding normalized norms are uniformly
bounded.

The elementary proof above does not infer a bounded spectral norm merely
from entrywise boundedness.

\hypertarget{a.3-delocalization-along-every-fixed-finite-gd-path}{%
\subsection{A.3 Delocalization along every fixed finite GD
path}\label{a.3-delocalization-along-every-fixed-finite-gd-path}}

\textbf{Theorem 1.} For each fixed \(T<\infty\), \[
\max_{0\le t\le T}\|b^t\|_2=O_P(1),\qquad
\max_{0\le t\le T}\|b^t\|_\infty
=O_P\!\left(\sqrt{\frac{\log d}{d}}\right). \tag{A.3.1}
\] The first bound has uniformly bounded fixed moments. The proof uses
the same reused training sample at every step.

\textbf{Proof.} Since \(\|r^t\|_2\le\sqrt n\), (A.1.1) gives directly \[
\|b^t\|_2\le \frac{T\lambda_d}{m}\|H\|_{\rm op}\sqrt n=O_P(1). \tag{A.3.2}
\]

For the coordinate bound, denote feature column \(j\) by \(h_j\),
replace that column by zero, and retain the denominator \(m\). Write
\(K^{-j}=K-h_jh_j^\top/m\), and let \(u^{-j,t},r^{-j,t}\) be the
corresponding GD scores and residuals. The matrices satisfy
\(0\preceq K^{-j}\preceq K\).

Conditional on \(X,Y,F_{-j}\), the vector \(r^{-j,t}\) is independent of
\(W_j\), bounded in Euclidean norm by \(\sqrt n\), and \[
G\mapsto (r^{-j,t})^\top\phi(XG/\sqrt d)
\] is odd and \(\|X\|_{\rm op}\sqrt{n/d}\)-Lipschitz for
\(G\sim N(0,I_d)\). Its conditional mean is zero. Gaussian concentration
followed by a union bound over \(j\le m,t<T\), on
\(\|X\|_{\rm op}\le C\sqrt d\), therefore yields \[
\max_{j,t<T}|h_j^\top r^{-j,t}|=O_P(\sqrt{d\log d}). \tag{A.3.3}
\] The conditioning here is essential: later residuals in the original
model are not independent of \(W_j\).

Set \(e^{j,t}=u^t-u^{-j,t}\). Since sigmoid is \(1/4\)-Lipschitz, \[
\|e^{j,t+1}\|_2
\le\left(1+\frac{\lambda_d\|K\|_{\rm op}}4\right)\|e^{j,t}\|_2
+\frac{\lambda_d\|h_j\|_2}{m}|h_j^\top r^{-j,t}|. \tag{A.3.4}
\] We have \(e^{j,0}=0\) and \(\|h_j\|_2\le\sqrt n\). On the
simultaneous high-probability bounds (A.2.2) and (A.3.3), iteration for
fixed \(T\) gives \[
\max_{j,t\le T}\|e^{j,t}\|_2=O_P(\sqrt{\log d}). \tag{A.3.5}
\] Thus \[
|h_j^\top r^t|
\le|h_j^\top r^{-j,t}|+\tfrac14\sqrt n\,\|e^{j,t}\|_2
=O_P(\sqrt{d\log d})
\] uniformly over the finitely many times and all features. Summing
(A.1.1) over time and dividing by \(m\) proves the second assertion in
(A.3.1). \(\square\)

\hypertarget{a.4-joint-fresh-test-gaussian-approximation-for-the-entire-finite-path}{%
\subsection{A.4 Joint fresh-test Gaussian approximation for the entire
finite
path}\label{a.4-joint-fresh-test-gaussian-approximation-for-the-entire-finite-path}}

Let \(\mathcal E_d\) contain the features and training sample, and
define the exact conditional covariances \[
A_{t,d}=\mathbb E_X[(\theta^\top X)S_t(X)\mid\mathcal E_d],\qquad
Q_{st,d}=\mathbb E_X[S_s(X)S_t(X)\mid\mathcal E_d]. \tag{A.4.1}
\] Write \(C_d\) for the covariance matrix of
\((\theta^\top X,S_1(X),\ldots,S_T(X))\). All means are zero by oddness.

\textbf{Theorem 2.} For fixed \(T\), in bounded-Lipschitz distance, \[
d_{\rm BL}\!\left(
\mathcal L((\theta^\top X,S_1(X),\ldots,S_T(X))\mid\mathcal E_d),
N(0,C_d)\right)\longrightarrow0\quad\text{in probability}. \tag{A.4.2}
\] This is a Gaussian approximation whose covariance is allowed to
depend on the realized network. It does not assert that \(C_d\) has a
deterministic limit.

\textbf{Proof.} Fix real \(u,v_1,\ldots,v_T\), put \(b=\sum_tv_tb^t\),
and let \[
f(x)=u\theta^\top x+b^\top\phi(Fx).
\] Then \(\mathbb Ef(X)=0\), and globally \[
\|\nabla f\|_2\le L:=|u|+\|F\|_{\rm op}\|b\|_2=O_P(1),
\quad
\|\nabla^2 f\|_{\rm op}\le
\varepsilon:=\|\phi''\|_\infty\|F\|_{\rm op}^2\|b\|_\infty=o_P(1).
\tag{A.4.3}
\]

Here is a self-contained normal-approximation argument. For the Gaussian
Ornstein--Uhlenbeck semigroup \(P_s\), set \[
U_f(x)=\int_0^\infty e^{-s}P_s\nabla f(x)\,ds,
\qquad T_f(x)=\nabla f(x)\cdot U_f(x).
\] Gaussian integration by parts, or integration of \(\frac d{ds}P_sf\),
gives for smooth bounded \(g\) \[
\mathbb E[f(X)g(f(X))]=\mathbb E[T_f(X)g'(f(X))],
\quad\mathbb ET_f=\mathbb Ef^2=:q.
\tag{A.4.4}
\] The semigroup derivative identity gives \(\|U_f\|_\infty\le L\) and
\(\|DU_f\|_{\rm op,\infty}\le\varepsilon/2\), whence
\(\|\nabla T_f\|_2\le3L\varepsilon/2\). Gaussian Poincaré yields \[
\mathbb E|T_f-q|\le\tfrac32L\varepsilon. \tag{A.4.5}
\] For \(\psi(t)=\mathbb Ee^{itf(X)}\), (A.4.4) gives
\(\psi'(t)=-tq\psi(t)-t\mathbb E[(T_f-q)e^{itf(X)}]\). Solving this
scalar differential equation and using \(q\ge0\) gives \[
|\psi(t)-e^{-qt^2/2}|\le\tfrac34t^2L\varepsilon. \tag{A.4.6}
\] No lower bound on \(q\) is needed.

The conditional covariance matrices are tight because (A.4.3) and
Poincaré bound all diagonal entries. Every subsequence admits a further
subsequence on which the bounds in (A.4.3) hold almost surely and
\(C_d\) converges along further deterministic realizations in a compact
set. Applying (A.4.6) to a countable dense collection of coefficient
vectors, tightness and the characteristic-function criterion identifies
every conditional subsequential limit with the corresponding Gaussian
covariance limit. This proves (A.4.2). \(\square\)

The old covariance-operator lemma applies uniformly to these
data-dependent weights because their Euclidean norms are tight: \[
A_{t,d}=a\theta^\top F^\top b^t+o_P(1),
\quad
Q_{st,d}=a^2(F^\top b^s)^\top(F^\top b^t)+\nu(b^s)^\top b^t+o_P(1). \tag{A.4.7}
\] Indeed the exact first formula is
\(A_{t,d}=\theta^\top F^\top D_A b^t\), where
\(D_A=\operatorname{diag}(A(\|W_j\|^2))\), and
\(\|D_A-aI\|_{\rm op}\to0\). The exact second formula is
\(Q_{st,d}=(b^s)^\top\Sigma_Fb^t\).

For any subsequence on which \((A_{t,d},Q_{st,d})\) converges in
probability to deterministic limits, Theorem 2 supplies the full joint
finite-time Gaussian test-law limit on that subsequence. With fresh
labels, apply the common Markov kernel
\(Y\mid Z\sim\mathrm{Bernoulli}(\sigma(\beta Z))\).

The total-variation and exact-calibration strengthening is proved
separately in Appendix C; its conditions are furnished by Theorem 1
together with a nonvanishing coefficient norm, for example Proposition 3
below.

\hypertarget{a.5-a-rigorous-nondegeneracy-condition}{%
\subsection{A.5 A rigorous nondegeneracy
condition}\label{a.5-a-rigorous-nondegeneracy-condition}}

Put \(\rho_\beta=\mathbb E\sigma'(\beta G)>0\). This is positive for
every fixed finite \(\beta\).

\textbf{Proposition 3.} \[
\|b^1\|_2^2\ge\frac{c^2v\rho_\beta}{\alpha\gamma}+o_P(1). \tag{A.5.1}
\] If \(\lambda_d\|K\|_{\rm op}\le4\), then pathwise \[
\|b^t\|_2\ge\tfrac12\|b^1\|_2\quad\text{for every integer }t\ge1. \tag{A.5.2}
\] In particular a sufficiently small, fixed positive \(c\) makes
(A.5.2) hold with probability tending to one. At every fixed positive
time in this regime, the teacher-orthogonal fresh-score variance is
bounded away from zero in probability.

\textbf{Proof.} Conditional on \(X,F\), write \(R=Y-1/2=\mu+\epsilon\),
where \(\mu_i=\sigma(\beta Z_i)-1/2\), independent centered
\(\epsilon_i\) have variance \(D_{ii}=\sigma'(\beta Z_i)\), and \(D\) is
diagonal. Since \[
\|b^1\|_2^2=\frac{\lambda_d^2}{m}R^\top KR,
\] expanding the quadratic form and using independence and bounded
fourth moments gives \[
\operatorname{Var}(R^\top KR\mid X,F)
\le C\{\operatorname{tr}(K^2)+\|K\mu\|_2^2\}
\le Cn\|K\|_{\rm op}^2. \tag{A.5.3}
\] Thus
\(\|b^1\|_2^2=\lambda_d^2m^{-1}(\mu^\top K\mu+\operatorname{tr}KD)+o_P(1)\).
The first term is nonnegative. Conditional on \(F\), the row quantities
\(D_{ii}K_{ii}\) are iid, bounded, and have mean \[
\frac1m\sum_j\mathbb E_X[\sigma'(\beta\theta^\top X)\phi(W_j^\top X)^2].
\] Uniformly over \(j\), \(\|W_j\|^2\to1\) and \(\theta^\top W_j\to0\).
Continuity of bounded Gaussian expectations therefore makes this mean
tend to \(v\rho_\beta\). The conditional LLN proves
\(n^{-1}\operatorname{tr}KD\to v\rho_\beta\), yielding (A.5.1).

For (A.5.2), let \(L(b)=\sum_i\ell(Y_i,h_i^\top b)\),
\(g_0=\nabla L(0)=-H^\top R\), and \(\eta=\lambda_d/m\). Its gradient is
\(\|H\|_{\rm op}^2/4\)-Lipschitz. Under the displayed step condition,
descent gives \(L(b^t)\le L(b^1)\le L(0)-\eta\|g_0\|^2/2\). Convexity
gives \(L(b^t)\ge L(0)+g_0^\top b^t\). Combining these and
Cauchy--Schwarz yields \(\|b^t\|\ge\eta\|g_0\|/2=\|b^1\|/2\). Finally,
the exact Hermite decomposition implies \[
Q_{tt,d}-A_{t,d}^2\ge(\nu-o_P(1))\|b^t\|_2^2, \tag{A.5.4}
\] uniformly for the fixed-time path. \(\square\)

\hypertarget{a.6-teacher-coordinate-removal-is-operator-small}{%
\subsection{A.6 Teacher-coordinate removal is
operator-small}\label{a.6-teacher-coordinate-removal-is-operator-small}}

Rotate coordinates so that \(\theta=e_1\). Write \(X_i=(z_i,x_i)\) and
\(W_j=(t_j,f_j)\), with \(x_i\in\mathbb R^{d-1}\),
\(f_j\in\mathbb R^{d-1}\). Let \(V\) have rows \(f_j^\top\),
\(U_{ij}=x_i^\top f_j\), \(H^0=\phi(U)\). The vectors \(z,t\), the noise
matrix \(H^0\), and label randomness conditional on \(z\) have the
required independence: \(H^0\) is independent of \((z,t,Y)\).

\textbf{Proposition 4.} For any fixed proportional number of rows,
including training rows together with an additional proportional block
of test rows, \[
\left\|\frac{H-H^0-a zt^\top}{\sqrt m}\right\|_{\rm op}\longrightarrow0
\quad\text{in probability}. \tag{A.6.1}
\]

\textbf{Proof.} Taylor's theorem gives \[
H-H^0-a zt^\top
=\operatorname{diag}(z)[\phi'(U)-a\mathbf1\mathbf1^\top]\operatorname{diag}(t)+R,
\quad |R_{ij}|\le C z_i^2t_j^2. \tag{A.6.2}
\] The remainder satisfies \[
\|R\|_F^2\le C\Big(\sum_i z_i^4\Big)\Big(\sum_jt_j^4\Big)=O_P(1),
\] so \(\|R\|_{\rm op}/\sqrt m=o_P(1)\).

Conditional on \(V\), the centered rows of the derivative matrix
\(\phi'(U)\) are independent. The same concentration-and-net proof as in
Section A.2 applies with Lipschitz constant
\(\|\phi''\|_\infty\|V\|_{\rm op}\); hence the centered matrix has
operator norm \(O_P(\sqrt d)\). Its row mean minus \(a\mathbf1\) is the
vector \(\delta_j=A(\|f_j\|^2)-a\). Bounded \(A'\) and
\(\mathbb E(\|f_j\|^2-1)^2=O(d^{-1})\) give \(\|\delta\|_2=O_P(1)\).
Thus \[
\|\phi'(U)-a\mathbf1\mathbf1^\top\|_{\rm op}=O_P(\sqrt d).
\] Finally \(\|z\|_\infty=O_P(\sqrt{\log d})\) and
\(\|t\|_\infty=O_P(\sqrt{\log d/d})\), so the first matrix in (A.6.2),
after division by \(\sqrt m\), is \(O_P(\log d/\sqrt d)=o_P(1)\).
\(\square\)

This is an operator-norm statement for the normalized \emph{feature
matrix}, not just for its population covariance. If \(A,\widetilde A\)
have bounded operator norm and \(\|A-\widetilde A\|_{\rm op}=o_P(1)\),
finite-time logistic GD in the parameter \(w=\sqrt m b\), with
normalized scores \(Aw\), differs by \(o_P(\sqrt m)\) in coefficient
norm and \(o_P(\sqrt n)\) in score norm. To verify, subtract the two
recursions, use \(1/4\)-Lipschitzness of sigmoid, \(\|r\|\le\sqrt n\),
and \(\|w^t\|=O_P(\sqrt m)\), and iterate the resulting linear
inequality for fixed \(T\).

\hypertarget{a.7-the-perpendicular-rf-matrix-satisfies-generalized-invariance}{%
\subsection{A.7 The perpendicular RF matrix satisfies generalized
invariance}\label{a.7-the-perpendicular-rf-matrix-satisfies-generalized-invariance}}

We verify the matrix assumptions of {[}5, Proposition D.1(b2){]}. The
following section then embeds the physical GD iteration in the
corresponding causal recursion.

Let \(N/d\to\delta>0\) be the total number of rows, \(p=d-1\),
\(X\in\mathbb R^{N\times p}\) have independent standard Gaussian
entries, and \(G\in\mathbb R^{m\times p}\) be an independent standard
Gaussian matrix. Set \[
V=G/\sqrt d,\qquad
A_0=\frac1{\sqrt m}\phi(XV^\top),\qquad r=N/m\to\delta/\gamma.
\tag{A.7.1}
\] For the almost-sure source-theorem application below, choose the
standard product construction in which the full data-array sequence is
independent of the full feature-array sequence, and different dimensions
may be taken independent. This choice is harmless for the final
in-probability conclusions, which depend only on the per-dimension laws.
It avoids conditioning on an entire feature sequence that could encode
data from another dimension under an adversarial cross-dimension
coupling.

\textbf{Theorem 5 (matrix hypothesis verification).} The matrix \(A_0\)
has the rectangular generalized-invariance property used by
Wang--Zhong--Fan, with the same limiting diagonal distribution as \[
A_G=\frac{aXV^\top+\sqrt\nu\,E}{\sqrt m},
\tag{A.7.2}
\] where \(E\) is an independent standard Gaussian \(N\times m\) matrix.
In particular:

\begin{enumerate}
\def\labelenumi{\arabic{enumi}.}
\tightlist
\item
  The operator norm is eventually bounded by a deterministic constant.
\item
  The singular-value distribution has a deterministic compactly
  supported limit, identical to that of (A.7.2), and its second moment
  is positive.
\item
  For every fixed integer \(k\ge0\) and every \(\epsilon>0\), the
  following inequalities hold eventually almost surely (with positive
  powers used for the first two lines): \[
  \max_i\left|((A_0A_0^\top)^k)_{ii}-\frac1N\operatorname{tr}(A_0A_0^\top)^k\right|
  +\max_{i\ne j}|((A_0A_0^\top)^k)_{ij}|
  \le m^{-1/2+\epsilon},
  \tag{A.7.3}
  \] \[
  \max_i\left|((A_0^\top A_0)^k)_{ii}-\frac1m\operatorname{tr}(A_0^\top A_0)^k\right|
  +\max_{i\ne j}|((A_0^\top A_0)^k)_{ij}|
  \le m^{-1/2+\epsilon},
  \tag{A.7.4}
  \] \[
  \max_{i,j}|((A_0A_0^\top)^kA_0)_{ij}|\le m^{-1/2+\epsilon}.
  \tag{A.7.5}
  \]
\item
  It is invariant under independent signed permutations of its rows and
  columns.
\end{enumerate}

\textbf{Proof of the uniform entry estimates.} For a matrix \(B\), let
\(\mathcal P_R(B)\) denote its Frobenius-metric projection onto the
closed convex set \(\{C:\|C\|_{\rm op}\le R\}\). This map clips singular
values at \(R\) and is Frobenius-nonexpansive. It commutes with left and
right orthogonal transformations.

Fix sufficiently large deterministic constants \(C_X,C_G,C_A\), and
define \[
\bar X=\mathcal P_{C_X\sqrt d}(X),\quad
\bar G=\mathcal P_{C_G\sqrt d}(G),\quad
B=\mathcal P_{C_A}\!\left(\frac1{\sqrt m}\phi(\bar X\bar G^\top/\sqrt d)\right).
\tag{A.7.6}
\] The exponential norm bounds in Section A.2 show \(B=A_0\) outside an
event of probability at most \(C e^{-cd}\). Thus \(B=A_0\) eventually
almost surely. This explicit clipping step avoids the invalid inference
that a derivative bound only on a nonconvex high-probability set is
automatically a global Lipschitz bound.

For any two base inputs, the product inequality and the scalar Lipschitz
property of tanh give \[
\|B(X,G)-B(X',G')\|_F
\le\frac C{\sqrt m}\big(\|X-X'\|_F+\|G-G'\|_F\big)
\le\frac {C'}{\sqrt d}\|(X,G)-(X',G')\|_F.
\tag{A.7.7}
\] Because \(\|B\|_{\rm op}\le C_A\), each scalar entry of
\((BB^\top)^k\), \((B^\top B)^k\), or \((BB^\top)^kB\) is a globally
\(C_k\)-Lipschitz function of \(B\) in Frobenius norm. This follows by
telescoping products; for example \[
\|(BB^\top)^k-(B'B'^\top)^k\|_F
\le 2k C_A^{2k-1}\|B-B'\|_F.
\] Consequently each such scalar entry is \(C_k'/\sqrt d\)-Lipschitz in
the independent standard Gaussian base entries.

Oddness of tanh and equivariance of all clipping maps imply that \(B\)
is invariant under row and column signed permutations. Therefore every
offdiagonal entry of \((BB^\top)^k\) or \((B^\top B)^k\) has expectation
zero; all diagonal entries in a given block have the same expectation;
and all entries of \((BB^\top)^kB\) have expectation zero. Gaussian
concentration and a union bound over \(O(d^2)\) entries give, for any
fixed \(\eta>0\), failure probability at most \[
C_kd^2\exp(-c_k d^{2\eta})
\] for the assertion that every entry differs from its expectation by at
most \(d^{-1/2+\eta}\). This sequence is summable. The difference
between a diagonal expectation and the realized average trace is bounded
by the same maximum diagonal error. Choose \(\eta<\epsilon\) and absorb
fixed constants and dimension ratios. Borel--Cantelli and \(B=A_0\)
eventually prove (A.7.3)--(A.7.5), simultaneously for the countable
collection of integer \(k\) and rational positive \(\epsilon\).
Arbitrary positive \(\epsilon\) follows by decreasing it to a rational
number.

\textbf{Proof of spectral matching.} Let \[
\Sigma_V=\mathbb E_x[\phi(Vx)\phi(Vx)^\top],\qquad
C_V=a^2VV^\top+\nu I_m.
\] The Hermite covariance calculation strengthens to almost sure
operator convergence here: \[
\|\Sigma_V-C_V\|_{\rm op}\to0. \tag{A.7.8}
\] Indeed Gaussian norm and spherical-inner-product tails imply,
eventually almost surely, \[
\max_j|\|f_j\|^2-1|=O(\sqrt{\log d/d}),\qquad
\max_{j\ne k}\frac{|f_j^\top f_k|}{\|f_j\|\|f_k\|}
=O(\sqrt{\log d/d}).
\] The offdiagonal higher-odd-Hermite remainder has operator norm at
most its maximal absolute row sum, hence at most
\(C m(\log d/d)^{3/2}=o(1)\). The diagonal remainder and linear Hermite
coefficient converge uniformly, proving (A.7.8).

Condition on any realization of the sequence \(V\) satisfying these
bounds. The Louart--Liao--Couillet spectral theorem applies to
\(\phi(XV^\top)\), since \(V^\top\) is bounded in operator norm, tanh is
Lipschitz, and all dimensions are proportional. It gives the same
spectral deterministic equivalent as a Gaussian sample matrix with row
covariance \(\Sigma_V\). A Gaussian coupling and \[
\|\Sigma_V^{1/2}-C_V^{1/2}\|_{\rm op}
\le\|\Sigma_V-C_V\|_{\rm op}^{1/2}
\] replace \(\Sigma_V\) by \(C_V\) without changing the limiting
spectral law. Conditional on \(V\), (A.7.2) has exactly Gaussian
independent rows with covariance \(C_V/m\). Fubini's theorem transfers
this conditional almost-sure conclusion to the joint model.

The empirical eigenvalue law of \(VV^\top\) tends to the
Marchenko--Pastur law with aspect \(\gamma\), so the empirical law of
\(C_V\) tends to its image under \(x\mapsto\nu+a^2x\). The fixed-point
description in the spectral theorem consequently has a deterministic
limit. More explicitly, if \(H\) is that image law, the feature-side
squared-singular law is characterized for \(\zeta>0\) by \[
\delta_\zeta=\int\frac{t}{\zeta+rt/(1+\delta_\zeta)}\,dH(t),\qquad
m(-\zeta)=\int\frac1{\zeta+rt/(1+\delta_\zeta)}\,dH(t).
\tag{A.7.9}
\] The positive solution is unique: divide its defining equation by
\(1+\delta_\zeta\) to obtain
\(\delta_\zeta/(1+\delta_\zeta)=\int t/[\zeta(1+\delta_\zeta)+rt]\,dH(t)\),
whose left side is strictly increasing and right side decreasing. The
bound \(0\le\delta_\zeta\le(\sup\operatorname{supp}H)/\zeta\) makes
continuity under weak convergence of uniformly compactly supported
covariance laws immediate by subsequences. The deterministic operator
bounds turn weak spectral convergence into convergence of every fixed
spectral moment. Also
\(m^{-1}\operatorname{tr}(A_0^\top A_0)\to r v>0\).
\textbf{Normalization for the AMP source:} its singular-value law \(D\)
is padded to the row dimension \(N\), whereas (A.7.9) is normalized to
the feature dimension \(m\). For \(k\ge1\),
\(\mathbb E[D^{2k}]=r^{-1}\int\lambda^k\,d\mu_{\rm feature}(\lambda)\);
in particular \(\mathbb E D^2=v\), not \(rv\). The aspect ratio is
\(r=N/m\). As an additional normalization check,
\(\mathbb E D^4=v^2+r(v^2+a^4\gamma)\), so the singular-value law is
nonconstant. This also supplies the positive-variance singular-law
hypothesis used by the AMP degeneracy-removal argument.

Finally (A.7.2) is bi-orthogonally invariant: Gaussian left
multiplication of \(X,E\) and Gaussian right multiplication of \(V,E\)
preserve their joint law. The independently signed-permutation
representation required by the generalized-invariance definition follows
by taking an independent copy of \(A_0\) and multiplying it by
independent uniform signed permutations. The law is unchanged, and these
permutations can be independent of any side information already
independent of \(A_0\). Items 1--4 are precisely the sufficient matrix
conditions of the Proposition D.1(b), which identifies the limiting
diagonal distribution. \(\square\)

Theorem 5 does not infer dynamics from an empirical spectral law. It
verifies the additional entrywise-polynomial and independence hypotheses
needed by a dynamical universality theorem. Combined with Proposition 4,
the residual issue is the finite-time algorithmic embedding, rather than
an unchecked Gaussian-equivalence conjecture for the RF matrix.

\hypertarget{a.8-causal-embedding-of-the-rank-one-logistic-recursion-in-rectangular-amp}{%
\subsection{A.8 Causal embedding of the rank-one logistic recursion in
rectangular
AMP}\label{a.8-causal-embedding-of-the-rank-one-logistic-recursion-in-rectangular-amp}}

This section supplies the exact algorithmic embedding and the
deterministic-overlap comparison. The Gaussian query/state recursion and
positive-definiteness proof are in Appendix B. The present formulas
specify their normalization and show that no extra universality for
arbitrary empirical scalar coefficients is being assumed.

Use a proportional block of additional test inputs, for example
\(N=2n\), and the deterministic row mask \(M_i=1\) for the \(n\)
training rows and \(0\) for the test rows. Set \[
\pi=n/N=1/2,\quad r=N/m\to2\alpha/\gamma,\quad
\lambda= c/\alpha,\quad s=a\sqrt\gamma,\quad v=\sqrt d\,t.
\tag{A.8.1}
\] Here the letter \(v\) denotes the feature-side teacher vector only
within this section; the scalar feature moment in earlier sections
remains \(\mathbb E\tanh(G)^2\). Write \(w_t=\sqrt m\,b^t\). Proposition
4 reduces the normalized feature matrix to \[
D=A_0+\frac{s}{m}zv^\top.
\] The small differences between \((\lambda_d,s_d,r_d)\) and their
limits are harmless by the same finite-time stability estimate. The
exact resulting physical iteration is \[
q_t^{\rm emp}=m^{-1}v^\top w_t,\quad
S_t=A_0w_t+s q_t^{\rm emp}z,\quad
r_t=M\odot[Y-\sigma(S_t)],
\tag{A.8.2}
\] \[
\rho_t^{\rm emp}=N^{-1}z^\top r_t,\qquad
w_{t+1}=w_t+\lambda[A_0^\top r_t+s r\rho_t^{\rm emp}v],\qquad w_0=0.
\tag{A.8.3}
\] The side vectors \((M,z,Y)\) and \(v\) are independent of \(A_0\).
Their empirical laws converge to \((M,Z,Y)\) and \(V\), where
\(M\sim\mathrm{Bernoulli}(\pi)\) is independent of \((Z,Y)\), \(Z,V\)
are standard Gaussian, and
\(Y\mid Z\sim\mathrm{Bernoulli}(\sigma(\beta Z))\). Gaussian and bounded
side variables have all moments. Polynomial density follows, for
example, by splitting the two finite coordinates \(M,Y\) and using
polynomial density under the Gaussian measure; their joint law is
dominated by a constant times Gaussian measure times counting measure.

Let \(a_{th},b_{th}\) be the deterministic rectangular-AMP correction
coefficients prescribed by the row-normalized singular law in Theorem 5.
The AMP has column fields \(Z_t^{\rm A}\), row fields \(Y_t^{\rm A}\),
feature iterates \(Q_t\), and residual iterates \(R_t\): \[
Z_t^{\rm A}=A_0^\top R_t-\sum_{h<t}b_{th}Q_h,\qquad
Y_t^{\rm A}=A_0Q_t-\sum_{h\le t}a_{th}R_h.
\tag{A.8.4}
\] Initialize \(Q_0=0\) and \(R_1=M\odot(Y-1/2)\). Define its
nonlinearities recursively by \[
Q_t=Q_{t-1}+\lambda\left[Z_t^{\rm A}+\sum_{h<t}b_{th}Q_h+s r\rho_{t-1}V\right],
\tag{A.8.5}
\] \[
R_{t+1}=M\left[Y-\sigma\!\left(Y_t^{\rm A}+\sum_{h\le t}a_{th}R_h+s q_tZ\right)\right],
\tag{A.8.6}
\] where \[
\rho_{t-1}=\mathbb E[Z R_t],\qquad q_t=\mathbb E[VQ_t]
\tag{A.8.7}
\] are taken in the limiting AMP law. Each \(Q_h\) in (A.8.5) is already
a deterministic function of the past column fields and feature-side
information, and each \(R_h\) in (A.8.6) is already a deterministic
function of past row fields and sample-side information. Thus these are
legitimate memory-dependent coordinatewise nonlinearities.

\textbf{Causality.} Before the \(t\)th column update, the law of \(R_t\)
and hence \(\rho_{t-1}\) are already known. The column field law and
(A.8.5) then determine \(Q_t\) and \(q_t\). Only after this does (A.8.6)
define \(R_{t+1}\). Consequently the overlap definitions introduce no
circular fixed point at a time step. The source's AMP covariance and
correction-coefficient prescriptions follow the same alternating order.

\textbf{Regularity.} At a fixed horizon, (A.8.5) is affine in its
Gaussian-field arguments. Equation (A.8.6) is recursively globally
Lipschitz in its Gaussian-field arguments because \(0\le M\le1\) and
\(\|\sigma'\|_\infty\le1/4\). To meet a theorem requiring definitions on
the entire Euclidean side-information space, replace the formal
arguments \(M,Y\) by their projections onto \([0,1]\); this changes no
observed value and keeps the Lipschitz constants uniform in the side
arguments. All functions have at most polynomial growth. The Bernoulli
sampling threshold is not an update nonlinearity: the already generated
label is supplied as side information.

\textbf{Exact cancellation.} Substituting (A.8.4) into (A.8.5) cancels
all \(b_{th}\) corrections. Substituting (A.8.4) into (A.8.6) cancels
all \(a_{th}\) corrections. Thus the finite-dimensional AMP trajectory
is exactly logistic GD with the scalar overlaps frozen to (A.8.7). These
identities are algebraic and do not use a limit.

\textbf{Deterministic-overlap freezing.} Suppose the source theorem is
applied to this AMP, using the positive-definiteness result of Appendix
B. Let \((\bar w_t,\bar r_t)\) denote its frozen-overlap physical
iterates. Joint empirical \(W_2\) convergence gives \[
\epsilon^q_t:=|m^{-1}v^\top\bar w_t-q_t|\to0,
\qquad
\epsilon^\rho_t:=|N^{-1}z^\top\bar r_t-\rho_t|\to0,
\tag{A.8.8}
\] since both products are continuous with quadratic growth. Define \[
e_t=\|w_t-\bar w_t\|_2/\sqrt m,\quad
h_t=\|r_t-\bar r_t\|_2/\sqrt N,\quad
V_m=\|v\|_2/\sqrt m,\quad Z_N=\|z\|_2/\sqrt N.
\] On \(\|A_0\|_{\rm op}\le C\), Cauchy--Schwarz and sigmoid's
derivative bound give \[
h_t\le\frac14\left[\left(C/\sqrt r+sZ_NV_m\right)e_t+sZ_N\epsilon^q_t\right],
\tag{A.8.9}
\] \[
e_{t+1}\le e_t+\lambda\left[\left(C\sqrt r+s rV_mZ_N\right)h_t+s rV_m\epsilon^\rho_t\right].
\tag{A.8.10}
\] We have \(e_0=0\), \(V_m,Z_N\to1\), and finitely many times.
Induction using (A.8.8) proves \(\max_{t\le T}(e_t+h_t)\to0\) in
probability. Differences in dimension-dependent constants add only
\(o_P(1)\) terms. Therefore the adaptive-overlap physical iteration has
exactly the same finite-time empirical limit as the frozen AMP.

This comparison is independent of whether the source's remark on
empirical Onsager coefficients includes the rank-one overlaps: it proves
the required replacement directly.

\hypertarget{a.9-assembled-finite-horizon-theorem-and-its-boundary}{%
\subsection{A.9 Assembled finite-horizon theorem and its
boundary}\label{a.9-assembled-finite-horizon-theorem-and-its-boundary}}

The preceding matrix inputs, Appendix B, and Appendix C give the
following assembly.

\textbf{Theorem 6 (original Gaussian tanh GD, each fixed finite
horizon).} Fix \(\alpha,\gamma,c\in(0,\infty)\), \(0\le\beta<\infty\),
and a fixed integer \(T\). Run the original model (A.1.1). Compute
\(q_{1:T},\Sigma_T\) by the explicit recursion (B.16) in Appendix B,
with \(s=a\sqrt\gamma\), and put \[
A_t=sq_t,\qquad Q_{st}=s^2q_sq_t+(\Sigma_T)_{st}.
\] Then \(\Sigma_T\) is positive definite. Conditional on the realized
features and training sample, the joint fresh-test law converges weakly
in probability to \[
\begin{gathered}
\bigl(Z,S_1,\ldots,S_T,Y\bigr),\qquad S_t=A_tZ+G_t,\\
G_{1:T}\sim N(0,\Sigma_T),\quad G\perp Z,\qquad
Y\mid Z\sim\mathrm{Bernoulli}(\sigma(\beta Z)).
\end{gathered}
\tag{A.9.1}
\] For every \(1\le t\le T\), the conditional law of \((Z,S_t,Y)\) in
fact converges in total variation in probability. The exact, unbinned,
model-conditional calibration error converges in probability and every
finite \(L^p\) to \[
C_t=\mathbb E_{S\sim N(0,Q_{tt})}|g_t(S)-\sigma(S)|,
\quad
 g_t(x)=\mathbb E_H\sigma\!\left(\beta\left[\frac{A_t}{Q_{tt}}x+
 \sqrt{\frac{(\Sigma_T)_{tt}}{Q_{tt}}}\,H\right]\right).
\tag{A.9.2}
\] Its population logistic risk converges in \(L^1\) to \[
R_t=\mathbb E_G\log\left(2\cosh\left(\frac{\sqrt{Q_{tt}}G}{2}\right)\right)-\kappa A_t,
\qquad
\kappa=\mathbb E[Z(\sigma(\beta Z)-1/2)].
\tag{A.9.3}
\] All these convergences are uniform over the finite index set
\(\{1,\ldots,T\}\). At time zero, \(C_0=0\) and \(R_0=\log2\) exactly.

\textbf{Assembly proof.} Theorem 5 and the side-information/regularity
checks in Section A.8 verify WZF's matrix and algorithm assumptions. the
positive-definiteness induction in Appendix B verifies all of its
nonsingularity assumptions without perturbation, for every \(c>0\). The
AMP embedding cancels the memory corrections exactly; (A.8.9)--(A.8.10)
remove frozen-overlap dependence; Proposition 4 and finite-time operator
stability restore the actual tanh features. On a mask-zero test row all
residual iterates vanish, so the limiting score is exactly
\(sq_tZ+G_t\), proving (A.9.1) for appended test rows.

Conditional on the original training environment those appended rows are
independent fresh inputs. Conditional LLNs for bounded observables
therefore yield the claimed conditional test-law convergence. For the
covariance identification, Gaussian concentration gives
\(\mathbb E_X|S_t(X)|^4\le C\|F\|_{\rm op}^4\|b^t\|_2^4=O_P(1)\); hence
conditional variances of the appended empirical quadratic averages are
\(O_P(N^{-1})\). This identifies the realized-model covariance limit
with \(Q\), including teacher covariances \(A\).

Theorem 1 supplies coefficient delocalization. Since Appendix B proves
\((\Sigma_T)_{tt}>0\), the deterministic smoothing theorem in Appendix C
upgrades each conditional pair to total variation. Adding the fresh
label is a common Markov kernel and contracts total variation. the
signed-measure inequality in Appendix C then proves (A.9.2). Gaussian
integration gives (A.9.3); uniform second-moment bounds for the scores
imply uniform integrability of logistic losses and conditional risks, so
the convergence is in \(L^1\). Finite unions give the stated uniformity
over \(t\le T\). \(\square\)

For clarity, \(C_t>0\) at every positive finite time. If \(\beta=0\),
the regression curve is \(1/2\) and the score is nondegenerate. If
\(\beta>0\) and \(A_t\le0\), monotonicity prevents \(g_t\) from equaling
the increasing sigmoid. If \(A_t>0\), dominated convergence gives \[
1-g_t(x)\sim
\exp\!\left[-\frac{\beta A_t}{Q_{tt}}x+
\frac{\beta^2(\Sigma_T)_{tt}}{2Q_{tt}}\right]
\quad(x\to\infty).
\] Equality with \(1-\sigma(x)\sim e^{-x}\) would require both
\(\beta A_t/Q_{tt}=1\) and \((\Sigma_T)_{tt}=0\), contradicting positive
definiteness. Continuity and the everywhere-positive score density turn
nonidentity into strictly positive calibration error.

For any training-dependent stopping time \(\tau_d\in\{0,\ldots,T\}\),
the realized population risk and calibration are approximated, in
probability, by \(R_{\tau_d}\) and \(C_{\tau_d}\). If the deterministic
finite-grid risk has a unique minimizer, every minimizer of the realized
population risk on that grid converges to it in probability. Minimizing
calibration alone includes the exact zero-score solution at time zero.
No monotonicity of the risk or calibration curve is asserted.

This theorem concerns every fixed finite horizon, at arbitrary positive
finite aspect ratios and step size. It does not, by itself, identify an
infinite-time limit or optimize over an unbounded stopping horizon.
Appendix D gives an additional all-time result under its explicitly
stated sample-ratio and small-step assumptions; those hypotheses must
remain visible whenever that extension is used.

\clearpage

\hypertarget{appendix-b.-explicit-causal-state-evolution}{%
\section{Appendix B. Explicit causal state
evolution}\label{appendix-b.-explicit-causal-state-evolution}}

Theorem, lemma, and proposition numbers in this appendix are local;
equation and section references carry the appendix letter.

\hypertarget{b.1-model-dimensions-and-the-physical-update}{%
\subsection{B.1 Model, dimensions, and the physical
update}\label{b.1-model-dimensions-and-the-physical-update}}

Let \(X_i\sim N(0,I_d)\), let \(\theta\) be a deterministic unit vector,
and let \[
Y_i\mid X_i\sim\operatorname{Bernoulli}(\sigma(\beta\theta^\top X_i)),
\qquad \sigma(x)=(1+e^{-x})^{-1},\quad 0\le\beta<\infty.
\] There are \(n\) training examples and \(m\) independent random
features \(W_j\sim N(0,I_d/d)\), with \[
n/d\to\alpha>0,\qquad m/d\to\gamma>0,\qquad
H_{ij}=\tanh(W_j^\top X_i).
\] Write \[
a=\mathbb E[G\tanh G],\qquad v=\mathbb E[\tanh^2G],\qquad
\nu=v-a^2>0,\quad G\sim N(0,1).
\] Strict positivity follows because \(\tanh G\) is not a scalar
multiple of \(G\) in Gaussian \(L^2\).

In Appendix A's normalization, \(w^t=\sqrt m\,b^t\), \(A=H/\sqrt m\),
and a mean-loss GD step of size \(cd\) gives \[
w^0=0,\qquad w^{t+1}=w^t+\lambda_d A^\top r^t,
\quad r^t=M\odot(Y-\sigma(Aw^t)),\quad \lambda_d=cd/n\to\lambda=c/\alpha.
\tag{B.1}
\] Here a proportional independent query block is appended to the rows,
and \(M_i=1\) on training rows and \(M_i=0\) on query rows. If the total
number of rows is \(N\), set \[
N/d\to\delta>\alpha,\qquad r=N/m\to\delta/\gamma,
\qquad p=n/N\to\alpha/\delta.
\tag{B.2}
\] The common choice \(N=2n\) gives \(p=1/2\), \(\delta=2\alpha\). No
result depends on that choice.

Rotate \(\theta\) to the first coordinate. Let \(z_i=\theta^\top X_i\)
and \(h_j=\sqrt d\,\theta^\top W_j\). Then \(z_i,h_j\) are standard
normal, and Appendix A's decomposition is \[
A=A_0+\frac{s_d}{m}zh^\top+E_d,
\quad s_d=a\sqrt{m/d}\to s=a\sqrt\gamma,
\quad \|E_d\|_{\rm op}=o_{\mathbb P}(1).
\tag{B.3}
\] The perpendicular matrix \(A_0\) is independent of \((z,Y,M,h)\). For
the exact spike matrix \(A_0+s_dzh^\top/m\), the physical GD update is
\[
Q_t=Q_{t-1}+\lambda_d\left[A_0^\top R_t+s_d r_d\ell_t^{(d)}h\right],
\quad R_{t+1}=M\odot\left[Y-\sigma(A_0Q_t+s_dq_t^{(d)}z)\right],
\tag{B.4}
\] where \(Q_t=w^t\), \(R_t=r^{t-1}\), \[
q_t^{(d)}=m^{-1}h^\top Q_t,\qquad
\ell_t^{(d)}=N^{-1}z^\top R_t,\qquad r_d=N/m.
\] In particular the column spike coefficient is \(s r\ell_t\), not
\(s\ell_t\).

\hypertarget{b.2-the-rf-rectangular-cumulants-have-a-closed-form}{%
\subsection{B.2 The RF rectangular cumulants have a closed
form}\label{b.2-the-rf-rectangular-cumulants-have-a-closed-form}}

Let \(\mu_k\) be the limiting \textbf{row-normalized} moments \[
\mu_k=\lim N^{-1}\operatorname{tr}(A_0A_0^\top)^k.
\] The squared singular-value law is padded with zeros to \(N\) values
when necessary. the Gaussian comparison matrix in Appendix A has
independent rows conditional on its feature covariance, whose limiting
covariance law is \[
\mathcal H=\mathcal L(\nu+a^2 L),\qquad L\sim\operatorname{MP}_{\gamma},
\quad \mathbb EL=1.
\] Let \(M(z)=\sum_{k\ge1}\mu_kz^k\) and
\(M_{\mathcal H}(z)=\sum_{k\ge1}(\int t^k\,d\mathcal H(t))z^k\). The
spectral fixed point in Appendix A, expanded at infinity, gives the
formal-power-series identity \[
rM(z)=M_{\mathcal H}\big(rz(1+M(z))\big).
\tag{B.5}
\] For completeness, writing \(m_F(-\zeta)\) for the feature-side
Stieltjes transform and \(e_\zeta\) for the auxiliary scalar in Appendix
A, its fixed point implies
\(r e_\zeta/(1+e_\zeta)=1-\zeta m_F(-\zeta)\). Since the feature moment
series is \(rM\), substitution \(z=-1/\zeta\) gives \(e_\zeta=-M/(1+M)\)
and then (B.5).

The ordinary free \(R\)-transform of \(\mathcal H\) is \[
R_{\mathcal H}^{\rm ord}(u)=\nu+\frac{a^2}{1-a^2\gamma u}.
\] Equivalently, \[
M_{\mathcal H}(w)=w(1+M_{\mathcal H}(w))
\left[\nu+\frac{a^2}{1-a^2\gamma w(1+M_{\mathcal H}(w))}\right].
\] Combine this identity with (B.5), and set \(t=z(1+M)(1+rM)\). Since
\(r\gamma=\delta\), this yields \[
M(z)=\nu t+\frac{a^2t}{1-a^2\delta t}.
\tag{B.6}
\] Fan's rectangular moment/cumulant identity is
\(M(z)=R_{\rm rect}(z(1+M(z))(1+rM(z)))\). Therefore \[
\boxed{R_{\rm rect}(t)=\nu t+\frac{a^2t}{1-a^2\delta t},\qquad
\kappa_2=v,\quad\kappa_{2k}=a^{2k}\delta^{k-1}\ (k\ge2).}
\tag{B.7}
\] This is a formal-series identity; no convergence radius or small-step
condition is needed. In particular \[
\mu_1=v,\qquad
\mu_2=(1+r)v^2+\delta a^4
=v^2+r(v^2+a^4\gamma),\qquad \kappa_4=\delta a^4.
\tag{B.8}
\]

\hypertarget{b.3-exact-finite-dimensional-memory-and-covariance-formulas}{%
\subsection{B.3 Exact finite-dimensional memory and covariance
formulas}\label{b.3-exact-finite-dimensional-memory-and-covariance-formulas}}

Use \(R_t\) for the scalar row residual and \(Q_t\) for the scalar
column iterate. The centered row and column Gaussian fields will be
\(G_t\) and \(Z_t\), respectively. Their covariance matrices are
\(\Sigma_t\) and \(\Omega_t\). The teacher scalar is denoted
\(Z_\star\), and the independent column teacher coordinate is
\(H_\star\).

Define the \(t\times t\) matrices \[
\Delta_t=(\mathbb E[R_iR_j])_{i,j\le t},\quad
\Gamma_t=(\mathbb E[Q_iQ_j])_{i,j\le t},
\] \[
(\Phi_t)_{ij}=\mathbb E\left[\frac{\partial R_i}{\partial G_j}\right]
\quad(j<i),\qquad
(\Psi_t)_{ij}=\mathbb E\left[\frac{\partial Q_i}{\partial Z_j}\right]
\quad(j\le i),
\tag{B.9}
\] and put zero in the remaining entries. Thus \(\Phi_t\) is strictly
lower triangular and \(\Psi_t\) is lower triangular. All derivatives
hold the deterministic recursion coefficients fixed.

Our coefficient matrices \(\mathsf A_t=(a_{ij})\) and
\(\mathsf B_t=(b_{ij})\) are \textbf{lower} triangular. They are the
transposes of Fan's displayed upper-triangular matrices \(A_t,B_t\). Set
\(h_0=a^2\delta\). Then the explicit formulas are \[
\boxed{\begin{aligned}
\mathsf A_t&=\nu\Psi_t+a^2\Psi_t(I-h_0\Phi_t\Psi_t)^{-1},\\
\mathsf B_t&=r\nu\Phi_t+ra^2\Phi_t(I-h_0\Psi_t\Phi_t)^{-1},\\
\Sigma_t&=\nu\Gamma_t+a^2(I-h_0\Psi_t\Phi_t)^{-1}
\left[\Gamma_t+h_0\Psi_t\Delta_t\Psi_t^\top\right]
(I-h_0\Psi_t\Phi_t)^{-\top},\\
\Omega_t&=r\nu\Delta_t+ra^2(I-h_0\Phi_t\Psi_t)^{-1}
\left[\Delta_t+h_0\Phi_t\Gamma_t\Phi_t^\top\right]
(I-h_0\Phi_t\Psi_t)^{-\top}.
\end{aligned}}
\tag{B.10}
\] Every inverse is a finite polynomial: \(\Phi_t\Psi_t\) and
\(\Psi_t\Phi_t\) are strictly lower triangular, so
\((I-h_0L)^{-1}=\sum_{j=0}^{t-1}h_0^jL^j\).

Here is an explicit check against the primary-source formulas. With
\(L=\Phi\Psi\) and \(K=\Psi\Phi\), Fan's covariance coefficients, in our
notation, are \[
\Theta^{(j)}=\sum_{i=0}^jL^i\Delta(L^\top)^{j-i}
+\sum_{i=0}^{j-1}L^i\Phi\Gamma\Phi^\top(L^\top)^{j-1-i},
\] \[
\Xi^{(j)}=\sum_{i=0}^jK^i\Gamma(K^\top)^{j-i}
+\sum_{i=0}^{j-1}K^i\Psi\Delta\Psi^\top(K^\top)^{j-1-i}.
\] The second sums are empty for \(j=0\). His prescriptions give
\(\Sigma=\sum_{j\ge0}\kappa_{2j+2}\Xi^{(j)}\),
\(\Omega=r\sum_{j\ge0}\kappa_{2j+2}\Theta^{(j)}\),
\(\mathsf A=\sum_{j\ge0}\kappa_{2j+2}\Psi(\Phi\Psi)^j\), and
\(\mathsf B=r\sum_{j\ge0}\kappa_{2j+2}\Phi(\Psi\Phi)^j\). Inserting
(B.7) and summing the two finite geometric series proves (B.10). In
particular no additional factor of \(r\) belongs inside either
covariance bracket.

\hypertarget{b.4-causal-recursion-and-its-implementation-without-circular-definitions}{%
\subsection{B.4 Causal recursion and its implementation without circular
definitions}\label{b.4-causal-recursion-and-its-implementation-without-circular-definitions}}

Let \[
Z_\star,H_\star\sim N(0,1),\quad
Y\mid Z_\star\sim\operatorname{Bernoulli}(\sigma(\beta Z_\star)),
\quad M\sim\operatorname{Bernoulli}(p),
\] with \(M\) independent of \((Z_\star,Y)\) and \(H_\star\) on the
separate column probability space. Initialize \[
R_1=M(Y-1/2),\qquad Q_0=0,\qquad q_0=0.
\] For \(t=1,2,\ldots\), carry out the following operations.

\begin{enumerate}
\def\labelenumi{\arabic{enumi}.}
\tightlist
\item
  From the already defined row functions
  \(R_{1:t}(G_{1:t-1},Z_\star,Y,M)\) and \(\Sigma_{t-1}\), compute
  \(\Delta_t\), \(\Phi_t\), and \(\ell_t=\mathbb E[Z_\star R_t]\).
\item
  Embed the previous \(\Psi_{t-1},\Gamma_{t-1}\) in \(t\times t\)
  matrices with zero last row/column. Compute \(\mathsf B_t,\Omega_t\)
  from (B.10). Their values do not depend on the missing current column
  function: the last column of \(\Phi_t\) is zero, so \(\Phi_t\Psi_t\)
  and \(\Phi_t\Gamma_t\Phi_t^\top\) only use previous column
  information.
\item
  Let \(Z_{1:t}\sim N(0,\Omega_t)\) independently of \(H_\star\),
  consistently with previous stages. Define \[
  Q_t=Q_{t-1}+\lambda\left[Z_t+\sum_{j<t}b_{tj}Q_j+s r\ell_tH_\star\right],
  \qquad q_t=\mathbb E[H_\star Q_t].
  \tag{B.11}
  \]
\item
  Compute the current derivative row of \(\Psi_t\) and \(\Gamma_t\),
  then compute \(\mathsf A_t,\Sigma_t\) using (B.10). Let
  \(G_{1:t}\sim N(0,\Sigma_t)\) independently of \((Z_\star,Y,M)\).
  Define \[
  S_t=G_t+\sum_{j\le t}a_{tj}R_j+s q_tZ_\star,
  \qquad R_{t+1}=M(Y-\sigma(S_t)).
  \tag{B.12}
  \]
\end{enumerate}

All column functions are linear. More precisely, \[
Q_{1:t}=\Psi_tZ_{1:t}+q_{1:t}H_\star,\qquad
\Gamma_t=\Psi_t\Omega_t\Psi_t^\top+q_{1:t}q_{1:t}^\top,
\tag{B.13}
\] and their coefficients can be calculated without integration: \[
(\Psi_t)_{tj}=(\Psi_{t-1})_{t-1,j}
+\lambda\mathbf1_{\{j=t\}}+\lambda\sum_{k<t}b_{tk}(\Psi_{t-1})_{kj},
\] \[
q_t=q_{t-1}+\lambda\sum_{j<t}b_{tj}q_j+\lambda s r\ell_t,
\tag{B.14}
\] where nonexistent entries at \(t=1\) are zero. In particular every
diagonal entry of \(\Psi_t\) equals \(\lambda\).

For the row derivatives, differentiate the already explicit scalar
recursion: \[
\frac{\partial R_{t+1}}{\partial G_j}
=-M\sigma'(S_t)\left[\mathbf1_{\{j=t\}}
+\sum_{k\le t}a_{tk}\frac{\partial R_k}{\partial G_j}\right],\quad j\le t.
\tag{B.15}
\] Thus the remaining computations are finite-dimensional Gaussian
integrals and a Bernoulli average; there is no unknown limiting quantity
outside the recursion.

\hypertarget{b.5-query-block-size-cancels-exactly}{%
\subsection{B.5 Query-block size cancels
exactly}\label{b.5-query-block-size-cancels-exactly}}

The apparent dependence on \(\delta\) in (B.7)--(B.12) is only a
bookkeeping choice. Let \(R^{\rm tr}_t\) be the row recursion
conditioned on \(M=1\). Define \[
\widehat\Delta_t=\alpha(\mathbb E[R_i^{\rm tr}R_j^{\rm tr}])_{i,j\le t},\qquad
\widehat\Phi_t=\alpha(\mathbb E[\partial_{G_j}R_i^{\rm tr}])_{i,j\le t},
\qquad \widehat\ell_t=\alpha\mathbb E[Z_\star R_t^{\rm tr}].
\] Since \(R_t=0\) when \(M=0\), \[
\widehat\Delta_t=\delta\Delta_t,\qquad
\widehat\Phi_t=\delta\Phi_t,\qquad \widehat\ell_t=\delta\ell_t.
\] Define \(T_t=(I-a^2\Psi_t\widehat\Phi_t)^{-1}\) and
\(U_t=(I-a^2\widehat\Phi_t\Psi_t)^{-1}\). Then (B.10) becomes \[
\boxed{\begin{aligned}
\mathsf A_t&=\nu\Psi_t+a^2\Psi_tU_t,\\
\mathsf B_t&=\gamma^{-1}\left[\nu\widehat\Phi_t+a^2\widehat\Phi_tT_t\right],\\
\Sigma_t&=\nu\Gamma_t+a^2T_t\left[\Gamma_t+a^2\Psi_t\widehat\Delta_t\Psi_t^\top\right]T_t^\top,\\
\Omega_t&=\gamma^{-1}\left\{\nu\widehat\Delta_t+a^2U_t
\left[\widehat\Delta_t+a^2\widehat\Phi_t\Gamma_t\widehat\Phi_t^\top\right]U_t^\top\right\}.
\end{aligned}}
\tag{B.16}
\] Replace \(s r\ell_t\) in (B.11) by \((s/\gamma)\widehat\ell_t\) and
use \(M=1\) in the training scalar recursion. This is an entirely
\(\delta\)-free prescription, depending only on
\((\alpha,\gamma,\beta,c,a,\nu)\).

On a query row \(M=0\), every \(R_j\) is zero. The physical test path is
therefore \[
\boxed{S_t^{\rm test}=s q_tZ_\star+G_t,\qquad
G_{1:T}\sim N(0,\Sigma_T)\text{ independently of }Z_\star.}
\tag{B.17}
\] Its deterministic teacher/score overlaps are \[
\mathcal A_t=s q_t,\qquad
\mathcal Q_{st}=s^2q_sq_t+(\Sigma_T)_{st}.
\tag{B.18}
\] Equation (B.16) proves algebraically that these quantities are
independent of the number of appended queries.

\hypertarget{b.6-positive-definiteness-no-regularization-or-algorithm-change}{%
\subsection{B.6 Positive definiteness: no regularization or algorithm
change}\label{b.6-positive-definiteness-no-regularization-or-algorithm-change}}

Suppose \(c>0\). Equations (B.10) imply \[
\Omega_t\succeq r\nu\Delta_t,\qquad \Sigma_t\succeq\nu\Gamma_t.
\tag{B.19}
\] Initially \(\Delta_1=p/4>0\). If \(\Delta_t\succ0\), then
\(\Omega_t\succ0\). Since \(\Psi_t\) is triangular with diagonal
\(\lambda>0\), (B.13) gives \(\Gamma_t\succ0\), hence
\(\Sigma_t\succ0\).

To continue the induction, suppose \(R_{t+1}\) were in the linear span
of \(R_1,\ldots,R_t\) almost surely. Condition on
\((Z_\star,Y,M,G_1,\ldots,G_{t-1})\). On the event \(M=1\) of positive
probability, all previous residuals are fixed. But the conditional
variance of \(G_t\) is positive because \(\Sigma_t\succ0\), and \[
R_{t+1}=Y-\sigma\left(G_t+\sum_{j\le t}a_{tj}R_j+s q_tZ_\star\right)
\] is a strictly monotone nonconstant function of that conditional
Gaussian. This is impossible. Therefore \(\Delta_{t+1}\succ0\).
Induction establishes positive definiteness at every finite stage,
including the uninformative-teacher case \(\beta=0\).

The case \(c=0\) is immediate from (B.1): every score is zero. It need
not be passed through a theorem requiring nonsingular covariances.

\hypertarget{b.7-why-this-is-ordinary-gd-including-adaptive-overlaps}{%
\subsection{B.7 Why this is ordinary GD, including adaptive
overlaps}\label{b.7-why-this-is-ordinary-gd-including-adaptive-overlaps}}

For any fixed deterministic arrays \(a_{tj},b_{tj}\), introduce fields
\[
\mathbf Z_t=A_0^\top\mathbf R_t-\sum_{j<t}b_{tj}\mathbf Q_j,
\qquad
\mathbf G_t=A_0\mathbf Q_t-\sum_{j\le t}a_{tj}\mathbf R_j.
\] Using (B.11) and (B.12) coordinatewise adds these same terms back.
With empirical teacher overlaps, the resulting algorithm is identically
(B.4), at every finite dimension. This algebra does not impose an
Onsager correction on GD.

For the theorem application, first use the deterministic scalar overlaps
\(q_t,\ell_t\) of (B.11)--(B.12). The resulting auxiliary AMP has
coordinatewise nonlinearities, so WZF applies. Its empirical overlaps
converge to \(q_t,\ell_t\) by its joint \(W_2\) convergence with the
teacher side information. Then compare it with the empirical-overlap
recursion (B.4).

Here are explicit bounds closing that comparison. Normalize errors by
\(e^Q_t=\|Q_t^{\rm GD}-Q_t^{\rm aux}\|/\sqrt m\) and
\(e^R_t=\|R_t^{\rm GD}-R_t^{\rm aux}\|/\sqrt N\). On events where
\(\|A_0\|_{\rm op}\), \(\|h\|/\sqrt m\), \(\|z\|/\sqrt N\), and the
fixed dimension ratios are bounded, Cauchy--Schwarz gives \[
|q_t^{\rm GD}-q_t|\le C e^Q_t+
|m^{-1}h^\top Q_t^{\rm aux}-q_t|,
\] \[
|\ell_t^{\rm GD}-\ell_t|\le C e^R_t+
|N^{-1}z^\top R_t^{\rm aux}-\ell_t|.
\] The two trailing terms tend to zero. The physical updates and
\(\operatorname{Lip}(\sigma)=1/4\) then imply \[
e^Q_t\le e^Q_{t-1}+C e^R_t+o_{\mathbb P}(1),\qquad
e^R_{t+1}\le C e^Q_t+o_{\mathbb P}(1).
\tag{B.20}
\] The initial error is zero, so both errors vanish at every fixed time.
Replacing \(s_d,\lambda_d,r_d\) by their limits only adds vanishing
terms of the same kind.

Finally, \(\|E_d\|_{\rm op}=o_{\mathbb P}(1)\) in (B.3) transfers this
conclusion to the \textbf{actual tanh RF matrix}. At fixed time
\(\|Q_t\|/\sqrt m=O_{\mathbb P}(1)\) and \(\|R_t\|/\sqrt N\le1\). Thus
every occurrence of \(E_dQ_t\) or \(E_d^\top R_t\) is
\(o_{\mathbb P}(1)\) in the relevant normalized norm, and the same
finite induction proves convergence of actual-GD iterates and score
vectors.

The appended queries never enter the update, since their mask is zero.
Conditional on the training data and features, their score paths are iid
fresh-test paths. A conditional law of large numbers for bounded
Lipschitz functions therefore converts empirical query convergence into
convergence in probability of the conditional fresh-test law to (B.17).
Quadratic observables also converge: the fixed-time operator bounds and
Gaussian Lipschitz concentration give bounded conditional fourth moments
in probability, so the conditional variance of their query averages is
\(O_{\mathbb P}(1/N)\). Consequently the realized fresh-test covariance
in Appendix A converges to (B.18).

\hypertarget{b.8-source-hypotheses-and-normalization}{%
\subsection{B.8 Source hypotheses and
normalization}\label{b.8-source-hypotheses-and-normalization}}

The primary sources are:

\begin{itemize}
\tightlist
\item
  Zhou Fan, \emph{Approximate Message Passing algorithms for
  rotationally invariant matrices}, Annals of Statistics 50(1), 197--224
  (2022), \href{https://arxiv.org/pdf/2008.11892}{arXiv:2008.11892v5}.
  In that paper, Section 2.4 gives the rectangular cumulant convention,
  and Section 5.1, equations (5.7)--(5.13), gives the memory/covariance
  prescription. Corollary 5.4 addresses degeneracy for the invariant
  comparison ensemble.
\item
  Tianhao Wang, Xinyi Zhong, and Zhou Fan, \emph{Universality of
  Approximate Message Passing algorithms and tensor networks}, Annals of
  Applied Probability 34(4), 3943--3994 (2024),
  \href{https://arxiv.org/html/2206.13037v5}{arXiv:2206.13037v5}. The
  applicable result is Theorem 2.22 with Assumption 2.17, algorithm
  (2.14), and prescriptions (D.4)--(D.5). Proposition D.1(b) supplies
  the matrix criterion checked in Appendix A.
\end{itemize}

The application checks are as follows.

\begin{longtable}[]{@{}
  >{\raggedright\arraybackslash}p{(\columnwidth - 2\tabcolsep) * \real{0.5000}}
  >{\raggedright\arraybackslash}p{(\columnwidth - 2\tabcolsep) * \real{0.5000}}@{}}
\toprule\noalign{}
\begin{minipage}[b]{\linewidth}\raggedright
Requirement
\end{minipage} & \begin{minipage}[b]{\linewidth}\raggedright
Verification
\end{minipage} \\
\midrule\noalign{}
\endhead
\bottomrule\noalign{}
\endlastfoot
Generalized-invariant matrix with invariant comparison law & Appendix A
Theorem 5, applied to the perpendicular RF matrix with all
training/query rows present \\
Independent initialization and side information & \(A_0\) is independent
of \(z,Y,M,h\) after teacher-coordinate separation \\
Nonzero initial AMP second moment & \(R_1=M(Y-1/2)\) and
\(\mathbb ER_1^2=p/4\); zero GD weights do not mean zero AMP
initialization \\
All side-information moments and polynomial density & Gaussian teacher
coordinates and bounded labels/masks have exponential moments; their
polynomial algebra is dense in \(L^2\) \\
Coordinatewise continuous, polynomial-growth, Lipschitz functions &
Column functions are linear; row functions are finite compositions of
sigmoid and affine functions with deterministic coefficients \\
Uniform Lipschitz constants over side information & Extend the mask off
\(\{0,1\}\) by the smooth bounded function \(\sin^2(\pi x/2)\); this
leaves actual iterates unchanged and keeps row derivatives uniformly
bounded \\
Eventual operator bound and positive spectral second moment & Appendix A
Theorem 5; \(\mathbb ED^2=v>0\) \\
Nonsingular \(\Omega_t,\Sigma_t\) & The direct induction in Section
B.6 \\
Empirical teacher-overlap dependence & Apply the theorem to
deterministic overlaps first, then use (B.20) \\
Fixed time and consistent finite aspect ratios & Every statement fixes
\(T\) before \(d\to\infty\); deterministic parameter replacements are
handled by (B.20) \\
Fresh-test extraction & Mask-zero appended rows, (B.17), and the
conditional law of large numbers \\
\end{longtable}

The source notation uses a matrix with ``\(m\) rows and \(n\) columns.''
Here these correspond to \(N\) rows and \(m\) features, so its aspect
parameter is \(r=\delta/\gamma\), \textbf{not} the RF ratio
\(\gamma=m/d\). The zero padding in Section B.2 is essential when
\(N>m\).

\hypertarget{b.9-exact-first-step-validation}{%
\subsection{B.9 Exact first-step
validation}\label{b.9-exact-first-step-validation}}

Put \[
\mu=\mathbb E[Z_\star(\sigma(\beta Z_\star)-1/2)].
\] Using the query-independent recursion (B.16), \[
\widehat\Delta_{11}=\alpha/4,\quad \widehat\Phi_{11}=0,
\quad \widehat\ell_1=\alpha\mu,\quad
\Omega_{11}=\frac{\alpha v}{4\gamma},\quad \Psi_{11}=\lambda.
\] Consequently \[
Q_1=\lambda Z_1+\frac{ca\mu}{\sqrt\gamma}H_\star,
\quad q_1=\frac{ca\mu}{\sqrt\gamma},\quad
\Gamma_{11}=\frac{c^2v}{4\alpha\gamma}+\frac{c^2a^2\mu^2}{\gamma}.
\] The forward memory coefficient is \(a_{11}=\lambda v\). It vanishes
from a test score because \(R_1=0\) on the test block. The Gaussian
residual variance is \[
\begin{aligned}
\Sigma_{11}
&=v\Gamma_{11}+a^4\lambda^2\widehat\Delta_{11}\\
&=c^2\left[\frac{\mu^2a^2v}{\gamma}
+\frac{a^4+v^2/\gamma}{4\alpha}\right].
\end{aligned}
\] Since \(s q_1=ca^2\mu\), the fresh-test law is exactly \[
\boxed{S_1^{\rm test}=ca^2\mu\,Z_\star+ au_1G,\qquad
\tau_1^2=c^2\left[\frac{\mu^2a^2v}{\gamma}
+\frac{a^4+v^2/\gamma}{4\alpha}\right],}
\tag{B.21}
\] with \(G\) independent standard normal. This reproduces both terms of
the first-step variance. Omitting \(\kappa_4=\delta a^4\) would
incorrectly remove \(c^2a^4/(4\alpha)\); forgetting the mask-zero query
block would incorrectly retain the training self term
\(\lambda v(Y-1/2)\).

\clearpage

\hypertarget{appendix-c.-total-variation-and-exact-calibration-transfer}{%
\section{Appendix C. Total variation and exact calibration
transfer}\label{appendix-c.-total-variation-and-exact-calibration-transfer}}

Theorem, lemma, and proposition numbers in this appendix are local;
equation and section references carry the appendix letter.

\hypertarget{c.1-notation-and-the-exact-calibration-functional}{%
\subsection{C.1 Notation and the exact calibration
functional}\label{c.1-notation-and-the-exact-calibration-functional}}

Write

\[
\begin{gathered}
X\sim N(0,I_d),\quad Z=\theta_d^\top X,\quad \|\theta_d\|=1,\quad S=f_d(X),\\
Y\mid X,\mathcal E_d\sim\operatorname{Bernoulli}(p(Z)),\qquad p(z)=\sigma(\beta z).
\end{gathered}
\]

The environment \(\mathcal E_d\) contains the training data and random
features, and the fresh \(X\) is independent of this environment. A
realized score is a deterministic function of this fresh input. We use
the \textbf{total-variation norm} convention

\[
\|\mu-\nu\|_{\rm TV}
 =\sup_{|h|\le1}\left|\int h\,d\mu-\int h\,d\nu\right|,
\tag{C.1.1}
\]

which is twice \(\sup_A|\mu(A)-\nu(A)|\) for probability measures.

Let \(\mu\) be the conditional law of \((Z,S)\). Its exact calibration
functional is

\[
\mathcal C_\beta(\mu)
 =\mathbb E_\mu\left|
       \mathbb E_\mu[p(Z)\mid S]-\sigma(S)\right|.
\tag{C.1.2}
\]

This equals the original calibration error conditioned on the predicted
probability and the environment, since sigmoid is one-to-one and
\(\mathbb E[Y\mid X,\mathcal E_d]=p(Z)\).

For a finite signed measure on the score line define

\[
\zeta_\mu(B)=\int\mathbf1_{\{s\in B\}}
                    [p(z)-\sigma(s)]\,\mu(dz,ds).
\]

Its Radon--Nikodym derivative relative to the score marginal is
\(\mathbb E_\mu[p(Z)\mid S=s]-\sigma(s)\), so

\[
\mathcal C_\beta(\mu)=\|\zeta_\mu\|_{\rm TV}.
\]

Consequently, for any two probability laws on \(\mathbb R^2\),

\[
\boxed{\quad
 |\mathcal C_\beta(\mu)-\mathcal C_\beta(\nu)|
 \le \|\mu-\nu\|_{\rm TV}.
 \quad}
\tag{C.1.3}
\]

Indeed \(|p(z)-\sigma(s)|\le1\), and multiplication by this bounded
function followed by projection onto the score coordinate contracts the
total-variation norm. Thus no continuity of a conditional-expectation
operator under weak convergence is being asserted.

There is a useful stronger statement. Let
\(g_\mu(s)=\mathbb E_\mu[p(Z)\mid S=s]\), and choose an everywhere
\([0,1]\)-valued version of \(g_\nu\) on the whole score line. The same
signed-measure argument, applied to the positive measures \(g_\mu\mu_S\)
and \(g_\nu\nu_S\), gives

\[
\int |g_\mu(s)-g_\nu(s)|\,\mu_S(ds)
 \le 2\|\mu-\nu\|_{\rm TV}.
\tag{C.1.4}
\]

This is convergence of the conditional calibration curve integrated
under the actual score law. It does not assert pointwise convergence at
every score.

\hypertarget{c.2-a-deterministic-gaussian-smoothing-theorem}{%
\subsection{C.2 A deterministic Gaussian smoothing
theorem}\label{c.2-a-deterministic-gaussian-smoothing-theorem}}

All norms on derivatives in this section are uniform over
\(x\in\mathbb R^d\). A Hilbert--Schmidt norm is denoted by
\(\|\cdot\|_{\rm HS}\).

\textbf{Theorem 2.1 (joint TV approximation with matched covariance).}
Let \(f_d:\mathbb R^d\to\mathbb R\) be smooth, odd functions, and let
\(X_d\sim N(0,I_d)\). Suppose, for fixed finite \(L,H\),

\[
\|\nabla f_d\|\le L,\qquad
\|\nabla^2 f_d\|_{\rm HS}\le H,\qquad
\|\nabla^2 f_d\|_{\rm op}\le\varepsilon_d\longrightarrow0.
\tag{C.2.1}
\]

Put

\[
A_d=\mathbb E[Z_df_d(X_d)],\quad
Q_d=\mathbb E f_d(X_d)^2,\quad
\delta_d=Q_d-A_d^2,
\quad Z_d=\theta_d^\top X_d.
\]

If \(\delta_d\ge\delta_*>0\), then

\[
\left\|\mathcal L(Z_d,f_d(X_d))-
 N\!\left(0,\begin{pmatrix}1&A_d\\A_d&Q_d\end{pmatrix}\right)
\right\|_{\rm TV}\longrightarrow0.
\tag{C.2.2}
\]

The covariance need not converge. In particular, if the pair converges
weakly to a nonsingular centered Gaussian with covariance
\(\left(\begin{smallmatrix}1&A\\A&Q\end{smallmatrix}\right)\), then its
convergence is in total variation.

The proof has two separate parts: weak Gaussian approximation from the
small Hessian operator norm, and a smoothing estimate using the bounded
Hessian Hilbert--Schmidt norm. Supplying both is necessary for a joint
TV conclusion.

\hypertarget{c.2.1-self-contained-weak-gaussian-approximation}{%
\subsubsection{C.2.1 Self-contained weak Gaussian
approximation}\label{c.2.1-self-contained-weak-gaussian-approximation}}

For a centered smooth Gaussian functional \(f(X)\), define the
Ornstein--Uhlenbeck semigroup

\[
P_tf(x)=\mathbb E f(e^{-t}x+\sqrt{1-e^{-2t}}X'),
\qquad
U_f(x)=\int_0^\infty e^{-t}P_t\nabla f(x)\,dt.
\]

Gaussian integration by parts and the semigroup generator give the
covariance identity

\[
\mathbb E[f(X)h(X)]
 =\mathbb E[U_f(X)\cdot\nabla h(X)].
\tag{C.2.3}
\]

For completeness, \(f-\mathbb Ef=-\int_0^\infty\partial_tP_tf\,dt\). The
generator is \(\mathcal L=\Delta-x\cdot\nabla\), and
\(-\mathbb E[(\mathcal Lg)h]=\mathbb E\nabla g\cdot\nabla h\). Also
\(\nabla P_tf=e^{-t}P_t\nabla f\). Substitution proves (C.2.3). The
bounded derivative assumptions, Gaussian moments, and an initial
truncation justify these integrations for the functions used here.

If \(\|\nabla f\|\le L_f\) and
\(\|\nabla^2f\|_{\rm op}\le\varepsilon_f\), then

\[
\|U_f\|\le L_f,\qquad
\|DU_f\|_{\rm op}
 =\left\|\int_0^\infty e^{-2t}P_t\nabla^2f\,dt\right\|_{\rm op}
 \le\varepsilon_f/2.
\]

Set \(T_f=\nabla f\cdot U_f\). Equation (C.2.3), with \(h=f\), gives
\(\mathbb ET_f=\operatorname{Var}(f(X))=:v_f\). Moreover

\[
\|\nabla T_f\|
 \le\tfrac32 L_f\varepsilon_f,
\qquad
\operatorname{Var}(T_f)
 \le\tfrac94 L_f^2\varepsilon_f^2
\tag{C.2.4}
\]

by Gaussian Poincare. Apply (C.2.3) to \(h=e^{itf}\). The characteristic
function \(\chi_f(t)=\mathbb E e^{itf}\) satisfies

\[
\chi_f'(t)=-t v_f\chi_f(t)
           -t\mathbb E[(T_f-v_f)e^{itf}].
\]

Solving this scalar differential equation and using (C.2.4) gives

\[
|\chi_f(t)-e^{-v_ft^2/2}|
 \le\tfrac34 L_f\varepsilon_f t^2.
\tag{C.2.5}
\]

Apply (C.2.5) to \(f=uZ_d+v f_d(X_d)\), whose gradient is bounded by
\(|u|+|v|L\) and Hessian operator norm by \(|v|\varepsilon_d\). Along
every subsequence where \((A_d,Q_d)\to(A,Q)\), it follows that
\((Z_d,f_d(X_d))\) converges weakly to the Gaussian with that
covariance. Such covariance-convergent subsequences exist because
Poincare gives \(Q_d\le L^2\), and \(|A_d|\le L\).

This derivation uses no independence of feature coordinates and no
assumed random-feature Gaussian equivalence.

\hypertarget{c.2.2-a-nondegenerate-gaussian-derivative-matrix-with-high-probability}{%
\subsubsection{C.2.2 A nondegenerate Gaussian derivative matrix with
high
probability}\label{c.2.2-a-nondegenerate-gaussian-derivative-matrix-with-high-probability}}

Temporarily suppress the dimension subscript. Write

\[
g=\nabla f,\quad J=\nabla^2f,\quad
P=I-\theta\theta^\top,\quad q=Pg,\quad
\Delta=\|q\|^2.
\]

The determinant of the Gaussian derivative Gram matrix of \((Z,f)\) is
exactly \(\Delta\). Let \(a=\theta\cdot g\). Gaussian integration by
parts gives \(\mathbb Ea=A\), while Poincare gives

\[
\operatorname{Var}(a)\le\mathbb E\|J\theta\|^2
 \le\varepsilon^2.
\]

Using Poincare again for \(f\),

\[
\mathbb E\Delta
 =\mathbb E\|g\|^2-\mathbb Ea^2
 \ge Q-A^2-\varepsilon^2.
\tag{C.2.6}
\]

Also \(\nabla\Delta=2Jq\), so

\[
\operatorname{Var}(\Delta)
 \le 4L^2\varepsilon^2.
\tag{C.2.7}
\]

If \(Q-A^2\ge\delta_*\), equations (C.2.6)--(C.2.7) imply

\[
\Pr\{\Delta\le\delta_*/4\}=o(1).
\tag{C.2.8}
\]

For example, when \(\varepsilon^2\le\delta_*/4\), the left side is at
most \(16L^2\varepsilon^2/\delta_*^2\), since
\(\mathbb E\Delta-\delta_*/4\ge\delta_*/2\).

Thus one needs neither a pointwise lower bound on \(\Delta\) nor inverse
moments of \(\Delta\) over its entire distribution.

\hypertarget{c.2.3-localized-integration-by-parts-and-the-smoothing-estimate}{%
\subsubsection{C.2.3 Localized integration by parts and the smoothing
estimate}\label{c.2.3-localized-integration-by-parts-and-the-smoothing-estimate}}

Choose a fixed smooth \(\chi:[0,\infty)\to[0,1]\) equal to zero on
\([0,\delta_*/8]\) and to one on \([\delta_*/4,\infty)\). On
\(\Delta>0\) define

\[
v_2=\frac q\Delta,\qquad
v_1=\theta-a v_2.
\]

They obey

\[
\theta\cdot v_1=1,\quad g\cdot v_1=0,
\qquad
\theta\cdot v_2=0,\quad g\cdot v_2=1.
\tag{C.2.9}
\]

Set \(V_i=\chi(\Delta)v_i\), extended by zero on the region where the
cutoff vanishes. These are globally smooth vector fields. On the support
of the cutoff, \(\Delta\ge\delta_*/8\). The bounds

\[
\|v_2\|=\Delta^{-1/2},\qquad
\|Dv_2\|_{\rm HS}
 \le\frac{\|J\|_{\rm HS}}\Delta
       +\frac{\|q\|\|\nabla\Delta\|}{\Delta^2}
 \le\frac{3H}\Delta
\]

and \(|a|\le L\), \(\|\nabla a\|\le H\), together with the cutoff
derivative bound, imply

\[
\|V_i\|+\|DV_i\|_{\rm HS}
 \le C(L,H,\delta_*),\qquad i=1,2.
\tag{C.2.10}
\]

The constant is independent of input dimension. No trace-of-Hessian
bound is used here.

For a vector field \(V\), its Gaussian divergence is
\(\mathfrak d(V)=X\cdot V-\operatorname{div}V\). Two Gaussian
integrations by parts give

\[
\mathbb E\mathfrak d(V)^2
 =\mathbb E\|V\|^2
   +\mathbb E\sum_{i,j}(\partial_iV_j)(\partial_jV_i)
 \le\mathbb E\|V\|^2+\mathbb E\|DV\|_{\rm HS}^2.
\tag{C.2.11}
\]

This identity can first be proved for smooth truncated fields and then
extended by Gaussian Sobolev approximation; (C.2.10) supplies the
required integrability. In particular its bound is dimension-free,
although the two terms defining \(\mathfrak d(V)\) separately need not
be.

For \(\psi\in C_c^\infty(\mathbb R^2)\), (C.2.9) and integration by
parts yield

\[
\mathbb E[\chi(\Delta)\,\partial_i\psi(Z,f)]
 =\mathbb E[\psi(Z,f)\mathfrak d(V_i)].
\tag{C.2.12}
\]

Therefore the positive subprobability measure

\[
\eta(B)=\mathbb E[\chi(\Delta)\mathbf1_{\{(Z,f)\in B\}}]
\]

has distributional first derivatives with total-variation norms bounded
by \(C(L,H,\delta_*)\). More directly, integrate \(\psi(x+h)-\psi(x)\)
along the segment from \(x\) to \(x+h\) and use (C.2.12). Taking the
supremum over smooth bounded test functions gives

\[
\|\tau_h\eta-\eta\|_{\rm TV}\le C\|h\|.
\tag{C.2.13}
\]

The same supremum equals the total-variation norm for finite regular
Borel measures; smooth approximation of continuous compactly supported
tests justifies this step.

Let \(\rho_r\) be a nonnegative smooth probability mollifier on
\(\mathbb R^2\), supported in a ball of radius \(r\). Integrating
(C.2.13) against the mollifier and writing \(\mu=\mathcal L(Z,f)\) gives

\[
\begin{aligned}
\|\mu-\mu*\rho_r\|_{\rm TV}
&\le2\|\mu-\eta\|_{\rm TV}
       +\|\eta-\eta*\rho_r\|_{\rm TV}\\
&\le2\mathbb E[1-\chi(\Delta)]+Cr
 =o(1)+Cr.
\end{aligned}
\tag{C.2.14}
\]

This is the required uniform regularization estimate for the joint law.

\hypertarget{c.2.4-completing-the-tv-proof}{%
\subsubsection{C.2.4 Completing the TV
proof}\label{c.2.4-completing-the-tv-proof}}

Pass to any subsequence along which \((A_d,Q_d)\to(A,Q)\). Its Gaussian
covariance is nonsingular because \(Q-A^2\ge\delta_*\). Section C.2.1
gives \(\mu_d\Rightarrow G_{A,Q}\). For fixed \(r>0\), convolution
densities \(\mu_d*\rho_r\) converge pointwise to \(G_{A,Q}*\rho_r\),
since the translated mollifier is a bounded continuous test function.
All these nonnegative densities have mass one, so Scheffe's argument
gives convergence in \(L^1\), equivalently TV. The elementary argument
is \(\int|p_d-p|=2-2\int\min(p_d,p)\to0\), by dominated convergence for
the minimum.

Use (C.2.14), this fixed-\(r\) convergence, and
\(\|G_{A,Q}-G_{A,Q}*\rho_r\|_{\rm TV}\to0\) as \(r\downarrow0\). First
taking the dimension limit and then \(r\downarrow0\) proves
\(\|\mu_d-G_{A,Q}\|_{\rm TV}\to0\). Gaussian densities with converging
nonsingular covariance matrices also converge pointwise and in \(L^1\),
so \(\|G_{A_d,Q_d}-G_{A,Q}\|_{\rm TV}\to0\). The triangle inequality
proves (C.2.2) on this subsequence. Compactness of the covariance
parameters then proves (C.2.2) for the whole sequence. This completes
Theorem 2.1.

For the last assertion of the theorem, bounded gradients and oddness
give the uniform fourth-moment bound \(\mathbb E f_d^4\le5L^4\): apply
Poincare to \(f_d^2\) and use \(\mathbb E f_d^2\le L^2\). Hence a weak
limit gives convergence of \(A_d,Q_d\) to \(A,Q\), and \(Q>A^2\)
supplies \(\delta_*>0\) eventually.

\hypertarget{c.3-structural-theorem-for-trained-random-feature-readouts}{%
\subsection{C.3 Structural theorem for trained random-feature
readouts}\label{c.3-structural-theorem-for-trained-random-feature-readouts}}

Let \(F_d\) have rows \(W_j^\top\), and, conditional on an arbitrary
environment containing \((F_d,b_d)\), let

\[
f_d(x)=\sum_{j=1}^m b_{d,j}\tanh(W_j^\top x).
\tag{C.3.1}
\]

There is \textbf{no independence assumption between the trained \(b_d\)
and the features or training sample}. The fresh input remains
independent of the environment.

Writing \(J_2=\sup_u|\tanh''u|\le2\), the exact derivative bounds are

\[
\begin{aligned}
\|\nabla f_d\|
 &\le\|F_d\|_{\rm op}\|b_d\|_2,\\
\|\nabla^2f_d\|_{\rm op}
 &\le J_2\|F_d\|_{\rm op}^2\|b_d\|_\infty,\\
\|\nabla^2f_d\|_{\rm HS}
 &\le J_2\|F_d\|_{\rm op}^2\|b_d\|_2.
\end{aligned}
\tag{C.3.2}
\]

For the last inequality use
\(\nabla^2 f_d=F_d^\top\operatorname{diag} (b_{d,j}\tanh''(W_j^\top x))F_d\)
and \(\|ABC\|_{\rm HS}\le\|A\|_{\rm op}\|B\|_{\rm HS}\|C\|_{\rm op}\).

Let

\[
A_d=\mathbb E_X[Zf_d(X)\mid\mathcal E_d],\quad
Q_d=\mathbb E_X[f_d(X)^2\mid\mathcal E_d],\quad
\delta_d=Q_d-A_d^2.
\tag{C.3.3}
\]

\textbf{Corollary 3.1 (conditional matched-Gaussian TV).} If

\[
\|F_d\|_{\rm op}=O_P(1),\qquad
\|b_d\|_2=O_P(1),\qquad
\|b_d\|_\infty\to0\text{ in probability},\qquad
\delta_d^{-1}=O_P(1),
\tag{C.3.4}
\]

then, writing \(\mu_d=\mathcal L_X(Z,f_d(X)\mid\mathcal E_d)\),

\[
\|\mu_d-G_{A_d,Q_d}\|_{\rm TV}\to0
\quad\text{in probability and every finite }L^p.
\tag{C.3.5}
\]

Here \(\delta_d^{-1}=O_P(1)\) includes positivity with probability
tending to one; values on a vanishing exceptional event can be assigned
arbitrarily. The \(L^p\) conclusion uses the bound of two for TV.

To justify the use of deterministic Theorem 2.1 with random
environments, first restrict to events with fixed bounds on
\(\|F_d\|_{\rm op},\|b_d\|_2,\delta_d^{-1}\). The deterministic
convergence is uniform in the sequential sense on each such bounded
class: if uniformity failed, choose one violating deterministic
environment at each dimension and contradict Theorem 2.1. The \(O_P(1)\)
assumptions make the complementary probabilities arbitrarily small.
Finally use \(\|b_d\|_\infty\to0\). No almost-sure boundedness of the
sequence of random environments is assumed.

The same reasoning gives a direct transfer version: if the first three
conditions in (C.3.4) hold and the conditional law of \((Z,f_d(X))\)
converges weakly in probability to a fixed nonsingular Gaussian, then
its convergence is in TV in probability. Uniform fourth moments on each
bounded environment event give convergence of the covariances, as in the
last paragraph of Section C.2.

\hypertarget{c.3.1-rf-covariance-relative-delocalization-and-nondegeneracy}{%
\subsubsection{C.3.1 RF covariance, relative delocalization, and
nondegeneracy}\label{c.3.1-rf-covariance-relative-delocalization-and-nondegeneracy}}

For Gaussian rows \(W_j\sim N(0,I_d/d)\), define

\[
a_1=\mathbb E[G\tanh G],\quad v=\mathbb E\tanh^2G,
\quad\nu=v-a_1^2>0.
\]

Strict positivity holds by the strict Cauchy--Schwarz inequality, since
\(\tanh G\) is not a multiple of \(G\) almost surely. Put
\(q_j=\|W_j\|^2\),
\(D_A=\operatorname{diag}(\mathbb E\tanh'(\sqrt{q_j}G))\), and
\(\Sigma_F=\mathbb E_X[h(X)h(X)^\top]\). The exact first-Hermite
decomposition and the operator bound in Appendix A give

\[
\Sigma_F=D_AFF^\top D_A+D_R+E_F,
\quad \|D_R-\nu I\|_{\rm op}\to0,
\quad\|E_F\|_{\rm op}\to0
\tag{C.3.6}
\]

in probability, when \(m/d\to\gamma\in(0,\infty)\). Here
\((D_R)_{jj}=\mathbb E\tanh^2(\sqrt{q_j}G)-q_j(D_A)_{jj}^2\). The
off-diagonal remainder obeys the cubic Hermite bound, whose squared
Frobenius expectation is \(O(d^{-1})\); row-norm concentration handles
the diagonal.

Gaussian integration by parts gives the exact covariance
\(A_d=\theta^\top F^\top D_A b_d\). Thus, uniformly over all readout
vectors, including data-dependent vectors,

\[
\begin{aligned}
Q_d-A_d^2
 &=\|(I-\theta\theta^\top)F^\top D_A b_d\|^2
   +b_d^\top(D_R+E_F)b_d\\
 &\ge (\nu-o_P(1))\|b_d\|_2^2.
\end{aligned}
\tag{C.3.7}
\]

In particular, a norm bounded away from zero supplies the last condition
of (C.3.4). There is also a scale-free alternative: if

\[
\Pr\{b_d=0\}\to0,\qquad
\frac{\|b_d\|_\infty}{\|b_d\|_2}\to0
\quad\text{in probability},
\tag{C.3.8}
\]

retain \(\|F_d\|_{\rm op}=O_P(1)\), which the Gaussian-row hypothesis
supplies, and normalize the score by \(\|b_d\|_2\). Its coefficient norm
is one and (C.3.7) gives a residual variance at least \(\nu/2\) with
probability tending to one. Corollary 3.1 applies to the normalized
score. The invertible, environment-dependent rescaling of its second
coordinate preserves TV, proving (C.3.5) for the original score as well.
This variant does not require its raw variance to stay bounded away from
zero.

Absolute coefficient delocalization alone cannot be silently substituted
for relative delocalization when the whole coefficient vector vanishes.
At exactly zero score the conditioning sigma-field changes, so this is a
substantive distinction for calibration.

\hypertarget{c.4-later-fixed-time-readouts-exact-gaussian-calibration-proxy}{%
\subsection{C.4 Later fixed-time readouts: exact Gaussian calibration
proxy}\label{c.4-later-fixed-time-readouts-exact-gaussian-calibration-proxy}}

For \(Q>A^2\), define the deterministic scalar functional

\[
C_\beta(A,Q)=\mathbb E_T\left|
 \mathbb E_H\sigma\!\left(\beta\left[
 \frac A{\sqrt Q}T+\sqrt{1-\frac{A^2}Q}\,H\right]\right)
 -\sigma(\sqrt Q\,T)\right|,
\tag{C.4.1}
\]

where \(T,H\) are independent standard normals. Corollary 3.1 and
(C.1.3) immediately imply

\[
\operatorname{CE}_d-C_\beta(A_d,Q_d)\longrightarrow0
\quad\text{in probability and every finite }L^p.
\tag{C.4.2}
\]

This is an exact calibration approximation for the original trained
network with its own conditional covariance, not calibration under a
different training distribution. It holds for any fixed training time
whose coefficients satisfy (C.3.4), or the relative variant (C.3.8).

If additionally \(A_d\to A\) and \(Q_d\to Q>A^2\) in probability,
bounded convergence in (C.4.1) gives the deterministic exact limit
\(\operatorname{CE}_d\to C_\beta(A,Q)\). This explicitly identifies what
a later-time overlap or DMFT theorem must provide.

Appendix A proves fixed-horizon coefficient bounds and delocalization
for actual GD. Appendix B proves positive residual variance at each
fixed positive time for every fixed positive step and supplies the
deterministic overlap limits. These results provide all inputs to the
transfer theorem above.

\hypertarget{c.4.1-risk-also-has-a-covariance-proxy}{%
\subsubsection{C.4.1 Risk also has a covariance
proxy}\label{c.4.1-risk-also-has-a-covariance-proxy}}

On events with bounded \(\|F\|_{\rm op}\|b_d\|_2\), Gaussian Poincare
and oddness bound the score second moment uniformly. Truncate logistic
loss, use (C.3.5) for the bounded truncation, then let the truncation
level grow. The upper bound \(0\le\ell(Y,S)\le\log2+|S|\) makes the tail
uniformly small on these events. Localization therefore proves

\[
L_d-\left\{
 \mathbb E_T\log\bigl(2\cosh(\sqrt{Q_d}T/2)\bigr)
 -\kappa A_d\right\}\longrightarrow0
\quad\text{in probability}.
\tag{C.4.3}
\]

The Gaussian formula uses
\(\mathbb E[S\sigma(\beta Z)]=A_d\mathbb E[Z\sigma(\beta Z)] =\kappa A_d\).
An \(L^1\) conclusion for (C.4.3) needs unconditional uniform
integrability; it is not inferred from \(O_P(1)\) bounds alone.

\hypertarget{c.4.2-positive-calibration-without-deterministic-overlap-convergence}{%
\subsubsection{C.4.2 Positive calibration without deterministic overlap
convergence}\label{c.4.2-positive-calibration-without-deterministic-overlap-convergence}}

For every \(\beta\ge0\), \(Q>A^2\), and \(Q>0\),

\[
C_\beta(A,Q)>0.
\tag{C.4.4}
\]

If \(\beta=0\) or \(A=0\), the conditional teacher mean is \(1/2\),
while \(\sigma(S)\) is nonconstant. If \(\beta>0,A<0\), the former mean
is strictly decreasing in the score whereas the prediction is strictly
increasing. If \(\beta>0,A>0\), write the former mean as

\[
g(s)=\mathbb E_H\sigma(\lambda s+rH),\quad
\lambda=\beta A/Q>0,\quad r=\beta\sqrt{1-A^2/Q}>0.
\]

If calibration held, continuity and positivity of the Gaussian density
would imply \(g(s)=\sigma(s)\) for every real \(s\). But dominated
convergence gives

\[
e^{\lambda s}(1-g(s))\longrightarrow e^{r^2/2}
\quad(s\to+\infty).
\]

Comparison with the logistic tail first forces \(\lambda=1\), then
\(e^{r^2/2}=1\), a contradiction.

The map \(C_\beta\) is continuous on the region \(Q>A^2\). Hence for
each fixed \(0<\delta_*\le Q_*<\infty\),

\[
c_*(\beta,\delta_*,Q_*)
 :=\min_{\substack{Q\le Q_*\\Q-A^2\ge\delta_*}}
 C_\beta(A,Q)>0.
\tag{C.4.5}
\]

The set is compact and \(Q\ge\delta_*\). Therefore if a finite
collection of trained times satisfies common deterministic upper and
lower covariance bounds with probability tending to one, as well as the
hypotheses of Corollary 3.1, then

\[
\Pr\left\{\min_{1\le t\le T}\operatorname{CE}_{t,d}
                    \ge c_*/2\right\}\longrightarrow1.
\tag{C.4.6}
\]

A union bound over the fixed number of times suffices. Deterministic
overlap convergence is unnecessary for this qualitative conclusion. It
is still necessary for an explicit deterministic trajectory of risk and
calibration, unless another argument identifies those functionals
without identifying the overlaps.

\clearpage

\hypertarget{appendix-d.-uniform-time-control-in-the-sample-rich-regime}{%
\section{Appendix D. Uniform time control in the sample-rich
regime}\label{appendix-d.-uniform-time-control-in-the-sample-rich-regime}}

Theorem, lemma, and proposition numbers in this appendix are local;
equation and section references carry the appendix letter.

\hypertarget{d.1-model-constants-and-theorem}{%
\subsection{D.1 Model, constants, and
theorem}\label{d.1-model-constants-and-theorem}}

Let \[
X_i\stackrel{\mathrm{iid}}\sim N(0,I_d),\qquad
Y_i\mid X_i\sim\operatorname{Bernoulli}(\sigma(\beta\theta^\top X_i)),
\quad \|\theta\|=1,
\] with conditionally independent labels, fixed finite \(\beta\geq0\),
and \(\sigma(u)=(1+e^{-u})^{-1}\). Independently, the rows of
\(F\in\mathbb R^{m\times d}\) are \(W_j\sim N(0,I_d/d)\). Put \[
h(x)=\tanh(Fx),\quad H_{ij}=h_j(X_i),\quad
\mathcal L_n(b)=\frac1n\sum_{i=1}^n
 \big\{\log(1+e^{h(X_i)^\top b})-Y_i h(X_i)^\top b\big\}.
\tag{D.1.1}
\] Assume \(n/d\to\alpha>0\), \(m/d\to\gamma>0\). The raw score is
\(b^\top h(x)\); in the alternative network notation \(b=a/\sqrt m\).
The GD update in the other notes is exactly \[
b^0=0,\qquad b^{t+1}=b^t-\eta_d\nabla\mathcal L_n(b^t),
\qquad \eta_d=\frac{cd}{m}.
\tag{D.1.2}
\]

Here are explicit, deliberately conservative constants. For
\(G\sim N(0,1)\), define \[
a_1=\mathbb E[G\tanh G],\quad v=\mathbb E\tanh^2G,
\quad \nu=v-a_1^2>0.
\] Choose any deterministic \(L\) for which
\(\|F\|_{\mathrm{op}}\leq L\) with probability tending to one. The
elementary Gaussian net bound permits, for example, \[
L=8\sqrt{1+\gamma}.
\] Set \[
s=\nu/2,\qquad \kappa=\sqrt{s/2},\qquad
p_0=\frac{s^2}{20L^4},\qquad
B=\sqrt{2\log(4/p_0)},
\] \[
q_\beta=\sigma(-\beta B),\quad
p_1=q_\beta p_0/2,\quad
\mathfrak a=\kappa p_1/4,\quad
C_*=\max\left\{1,\left(\frac{32L}{\kappa p_1}\right)^2\right\}.
\tag{D.1.3}
\] The ratio and step assumptions for this appendix are \[
\boxed{\quad\alpha/\gamma>C_*,\qquad
0<c<\frac{\gamma}{2\Lambda},\quad
\Lambda=16L^2.\quad}
\tag{D.1.4}
\] These define a nonempty proportional regime for each fixed
\(\beta,\gamma\). No assertion that these thresholds are sharp is made.

Further set \[
R_0=\frac{\log2}{\mathfrak a},\quad R=2R_0,
\quad C_H=64L^2,\quad \delta=p_0/4,
\quad M=R\sqrt{C_H/\delta},
\] \[
\chi=\frac{\kappa^2p_0}{16},\qquad
\mu=\chi\sigma'(M)>0,\qquad
\eta_-=\frac{c}{2\gamma},\qquad q=1-\eta_-\mu\in(0,1).
\tag{D.1.5}
\] All these constants are independent of \(d\), although some are very
small or very large.

\textbf{Theorem 1 (whole GD trajectory).} Under (D.1.4), with
probability tending to one:

\begin{enumerate}
\def\labelenumi{\arabic{enumi}.}
\tightlist
\item
  \(\mathcal L_n(b)\geq\mathfrak a\|b\|\) for every \(b\in\mathbb R^m\).
  There is a unique finite unregularized MLE \(b^*\), with
  \(\|b^*\|\leq R_0\).
\item
  Every GD iterate and every feature-deleted GD iterate defined in
  Section D.6 lies in \(\{\|b\|\leq R_0\}\). On the larger ball
  \(\{\|b\|\leq R\}\), \[
  \mu I\preceq\nabla^2\mathcal L_n(b)\preceq\Lambda I.
  \tag{D.1.6}
  \]
\item
  For every integer \(t\geq0\), simultaneously, \[
  \|b^t-b^*\|\leq R_0q^t.
  \tag{D.1.7}
  \]
\item
  The MLE and the whole trajectory are delocalized: \[
  \|b^*\|_\infty+\sup_{t\geq0}\|b^t\|_\infty
     =O_P\!\left(\sqrt{\frac{\log d}{d}}\right).
  \tag{D.1.8}
  \]
\item
  For a deterministic \(r_*>0\), \[
  \inf_{t\geq1}\|b^t\|\geq r_*,\qquad \|b^*\|\geq r_*
  \tag{D.1.9}
  \] with probability tending to one. One permissible choice is \[
  r_* =\frac c4\sqrt{\frac{v\rho_\beta}{\alpha\gamma}},
  \qquad \rho_\beta=\mathbb E\sigma'(\beta G)>0.
  \tag{D.1.10}
  \]
\end{enumerate}

The label mechanism is not independent of the feature vector after
conditioning on \(F\). The proof below does not make that assumption.

\hypertarget{d.2-the-good-feature-event}{%
\subsection{D.2 The good feature
event}\label{d.2-the-good-feature-event}}

Write \(\Sigma_F=\mathbb E_X h(X)h(X)^\top\) and
\(g_F=\mathbb E_X[(\theta^\top X)h(X)]\). There is an event
\(\mathcal F_d\), depending only on \(F\), on which \[
\|F\|_{\mathrm{op}}\leq L,\qquad
\Sigma_F\succeq sI,\qquad
\Sigma_F-g_Fg_F^\top\succeq sI,
\tag{D.2.1}
\] with \(\mathbb P(\mathcal F_d^c)=O(d^{-1})+O(de^{-c_Fd})\) for some
fixed \(c_F>0\). Constants in probability estimates may depend on fixed
model parameters.

For completeness, the covariance claim follows directly from Hermite
expansion. Put \(q_j=\|W_j\|^2\), \(A(q)=\mathbb E\tanh'(\sqrt qG)\),
and \(D_A=\operatorname{diag}(A(q_j))\). Oddness gives only odd Hermite
coefficients, and Gaussian integration by parts identifies the linear
coefficient. Thus \[
\Sigma_F=D_AFF^\top D_A+D_\nu+E,
\quad (D_\nu)_{jj}=\mathbb E\tanh^2(\sqrt{q_j}G)-q_jA(q_j)^2,
\tag{D.2.2}
\] where \(E_{jj}=0\) and, for \(j\ne k\), \[
|E_{jk}|\leq |\rho_{jk}|^3,\qquad
\rho_{jk}=\frac{W_j^\top W_k}{\|W_j\|\|W_k\|}.
\tag{D.2.3}
\] Indeed, Cauchy--Schwarz on the Hermite coefficients of degree at
least three bounds their sum by \(|\rho_{jk}|^3\), since tanh has second
moment at most one. Independent Gaussian directions give \[
\mathbb E\|E\|_F^2\leq
\frac{15m(m-1)}{d(d+2)(d+4)}=O(d^{-1}).
\tag{D.2.4}
\] Uniform concentration of \(q_j\) around one and continuity give
\(\min_j(D_\nu)_{jj}\geq3\nu/4\) except on an exponentially unlikely
event. Markov's inequality gives \(\|E\|_{\mathrm{op}}\leq\nu/4\) with
failure \(O(d^{-1})\). Also \[
g_F=D_AF\theta,
\quad
\Sigma_F-g_Fg_F^\top
=D_AF(I-\theta\theta^\top)F^\top D_A+D_\nu+E\succeq sI.
\tag{D.2.5}
\] The operator norm bound follows by a Gaussian sphere-net argument.
These observations prove (D.2.1).

Conditional on \(F\in\mathcal F_d\), for every unit \(u\in\mathbb R^m\),
the function \(x\mapsto u^\top h(x)\) is odd and \(L\)-Lipschitz. In
particular, \[
\mathbb E e^{t u^\top h(X)}\leq e^{t^2L^2/2},\qquad
s\leq\mathbb E(u^\top h(X))^2\leq L^2,
\quad
\mathbb E(u^\top h(X))^4\leq5L^4.
\tag{D.2.6}
\] The fourth-moment bound follows from Gaussian Poincare applied first
to the centered function and then to its square:
\(\mathbb E f^4\leq(\mathbb E f^2)^2+4L^2\mathbb E f^2\leq5L^4\).

\hypertarget{d.3-uniform-small-ball-events-and-noisy-label-coercivity}{%
\subsection{D.3 Uniform small-ball events and noisy-label
coercivity}\label{d.3-uniform-small-ball-events-and-noisy-label-coercivity}}

Everything in this section is conditional on a fixed
\(F\in\mathcal F_d\). The pairs \((X_i,Y_i)\) are iid under this
conditioning. For each unit \(u\), Paley--Zygmund, or Cauchy--Schwarz
applied above the threshold, gives \[
\mathbb P(|u^\top h(X)|\geq\kappa)\geq p_0.
\tag{D.3.1}
\] Let \(Z=\theta^\top X\) and \(T=2Y-1\). The chosen \(B\) has
\(\mathbb P(|Z|>B)\leq p_0/2\). On \(|Z|\leq B\), either label has
conditional probability at least \(q_\beta\). Therefore \[
\mathbb P(-T u^\top h(X)\geq\kappa)
\geq q_\beta\mathbb P(|u^\top h(X)|\geq\kappa,|Z|\leq B)
\geq p_1.
\tag{D.3.2}
\] This is a conditional-label calculation at each \(X\); it never
separates the dependent events involving \(h(X)\) and \(Z\).

Define the bounded cutoff \[
\psi(x)=
\begin{cases}
0,&x\leq\kappa/2,\\
2x/\kappa-1,&\kappa/2<x<\kappa,\\
1,&x\geq\kappa.
\end{cases}
\tag{D.3.3}
\] It is \(2/\kappa\)-Lipschitz and vanishes at zero. For either class
\[
f_u(X,Y)=\psi(-T u^\top h(X)),\qquad
\widetilde f_u(X)=\psi(|u^\top h(X)|),\quad \|u\|=1,
\] symmetrization and Rademacher contraction give \[
\mathbb E\sup_{\|u\|=1}|(P_n-P)f_u|
\leq\frac{8}{\kappa}
 \mathbb E\left\|\frac1n\sum_i\epsilon_iT_i h(X_i)\right\|
\leq\frac{8L}{\kappa}\sqrt{\frac mn}.
\tag{D.3.4}
\] The same bound holds for \(\widetilde f_u\), with \(T_i\) omitted.
The numerical factor eight safely includes the absolute-value form of
contraction. The last inequality uses \(T_i^2=1\) and the independent
Rademacher signs, not independence of \(T_i\) and \(h(X_i)\).

Since changing one sample changes each supremum by at most \(1/n\),
bounded differences imply an upper-tail bound \(e^{-2nz^2}\) above its
expectation. If \(n/m\geq C_*\), (D.3.4) is at most \(p_1/4\). Using
\(z=p_1/4\) for the signed class and \(z=p_0/4\) for the unsigned class
proves, simultaneously, except with probability \[
e^{-np_1^2/8}+e^{-np_0^2/8},
\tag{D.3.5}
\] that for every unit \(u\), \[
\frac1n\sum_i\psi(-T_i u^\top h_i)\geq p_1/2,
\qquad
\frac1n\#\{i:|u^\top h_i|>\kappa/2\}\geq p_0/2.
\tag{D.3.6}
\] Measurability of the suprema is unproblematic: the cutoff classes are
continuous in \(u\), so a countable dense subset of the sphere suffices.

The logistic loss obeys \(\log(1+e^{-Tz})\geq(-Tz)_+\). Also
\(x_+\geq(\kappa/2)\psi(x)\). Consequently, for every \(b\ne0\), using
\(u=b/\|b\|\), \[
\mathcal L_n(b)
\geq\|b\|\frac1n\sum_i(-T_i u^\top h_i)_+
\geq\frac{\kappa p_1}{4}\|b\|
=\mathfrak a\|b\|.
\tag{D.3.7}
\] The inequality at \(b=0\) is immediate. Thus the entire initial
sublevel set \(\{\mathcal L_n\leq\log2\}\) is contained in the
deterministic radius-\(R_0\) ball.

\hypertarget{d.4-robust-covariance-and-curvature-on-the-whole-ball}{%
\subsection{D.4 Robust covariance and curvature on the whole
ball}\label{d.4-robust-covariance-and-curvature-on-the-whole-ball}}

The unsigned event in (D.3.6) gives an arbitrary-deletion bound: for
\textbf{every} subset \(I\subseteq\{1,\ldots,n\}\) with
\(|I^c|\leq\delta n\), \[
\frac1n\sum_{i\in I}h_i h_i^\top\succeq\chi I_m.
\tag{D.4.1}
\] Indeed, for each unit \(u\), at least \(p_0n/2\) rows have
\(|u^\top h_i|>\kappa/2\); after arbitrary deletion of at most
\(p_0n/4\) rows, at least \(p_0n/4\) remain, contributing at least
\(\kappa^2p_0/16\). This proof does not union-bound over subsets and
does not assume deleted rows are independent of the data.

We also have, with conditional failure at most \(2e^{-3n}\) when
\(n\geq m\), \[
\|H\|_{\mathrm{op}}\leq8L\sqrt n,
\qquad \frac1nH^\top H\preceq C_H I_m.
\tag{D.4.2}
\] For a direct proof, take \(1/4\)-nets of the two unit spheres, of
cardinalities at most \(9^n,9^m\). Each bilinear form \(u^\top Hv\) is
centered \(L\)-subGaussian by independence of the rows and (D.2.6). The
bound at \(4L\sqrt n\), a union bound, and the factor-two net inequality
give failure \(2\exp\{(n+m)\log9-8n\}\leq2e^{-3n}\). This uses no
within-row independence.

For any \(\|b\|\leq R\), let \(I_b=\{i:|h_i^\top b|\leq M\}\). Equation
(D.4.2) gives \[
\frac{|I_b^c|}{n}\leq\frac{b^\top H^\top Hb}{nM^2}
\leq\frac{C_HR^2}{M^2}=\delta.
\] The logistic Hessian is therefore bounded simultaneously on the ball
by \[
\nabla^2\mathcal L_n(b)
=\frac1n\sum_i\sigma'(h_i^\top b)h_ih_i^\top
\succeq\sigma'(M)\frac1n\sum_{i\in I_b}h_ih_i^\top
\succeq\mu I_m.
\tag{D.4.3}
\] Globally, \(\nabla^2\mathcal L_n(b)\preceq C_HI/4=\Lambda I\). Thus
the desired curvature is simultaneous for all points in the ball, not
merely at a prespecified iterate.

Let \(\mathcal G_d\) be the feature event together with (D.3.6) and
(D.4.2). For large \(d\), (D.1.4) gives \(n/m\geq C_*\), and \[
\mathbb P(\mathcal G_d^c)
\leq\mathbb P(\mathcal F_d^c)
 +e^{-np_1^2/8}+e^{-np_0^2/8}+2e^{-3n}=o(1).
\tag{D.4.4}
\]

\hypertarget{d.5-mle-existence-and-uniform-contraction}{%
\subsection{D.5 MLE existence and uniform
contraction}\label{d.5-mle-existence-and-uniform-contraction}}

On \(\mathcal G_d\), coercivity and continuity give a finite global
minimizer. Its loss is at most \(\mathcal L_n(0)=\log2\), so its norm is
at most \(R_0\). Strict convexity on the containing ball proves
uniqueness; call the minimizer \(b^*\). Its full gradient is zero.

For probabilistic expressions involving the MLE, assign \(b^*=0\) on any
exceptional realization without a unique finite MLE. This arbitrary
extension has vanishing probability and does not change any conclusion.

For large \(d\), the step assumption implies \[
\eta_-\leq\eta_d\leq2c/\gamma<1/\Lambda.
\tag{D.5.1}
\] The global smoothness inequality gives \[
\mathcal L_n(b-\eta_d\nabla\mathcal L_n(b))
\leq\mathcal L_n(b)-\frac{\eta_d}{2}\|\nabla\mathcal L_n(b)\|^2.
\tag{D.5.2}
\] Thus all iterates remain in the initial sublevel set, and hence in
the radius-\(R_0\) ball.

For \(x,y\) in this ball, integrate the Hessian along their line
segment. It yields a symmetric matrix \(A_{x,y}\) with spectrum in
\([\mu,\Lambda]\) such that \[
[x-\eta_d\nabla\mathcal L_n(x)]-[y-\eta_d\nabla\mathcal L_n(y)]
=(I-\eta_dA_{x,y})(x-y).
\] Consequently the GD map contracts by \[
q_d=1-\eta_d\mu\leq q<1.
\tag{D.5.3}
\] Applying this with \(y=b^*\) proves (D.1.7), uniformly for every
integer time. No dimension-dependent burn-in is hidden in the
contraction factor.

\hypertarget{d.6-leave-one-feature-definitions-and-concentration}{%
\subsection{D.6 Leave-one-feature definitions and
concentration}\label{d.6-leave-one-feature-definitions-and-concentration}}

Let \(P_j\) set coordinate \(j\) to zero. Define restricted GD for every
data realization, including outside \(\mathcal G_d\), by \[
b^{-j,0}=0,\qquad
b^{-j,t+1}=b^{-j,t}-\eta_d P_j\nabla\mathcal L_n(b^{-j,t}).
\tag{D.6.1}
\] The coordinate \(j\) stays zero, so this update uses only
\(X,Y,F_{-j}\). On \(\mathcal G_d\), the same descent inequality applied
on the coordinate subspace shows that every restricted iterate has loss
at most \(\log2\) and norm at most \(R_0\). Therefore the same full
Hessian bounds apply simultaneously along the full and all restricted
trajectories.

Also define \(\bar b^{-j}\), for every data realization, to minimize
\(\mathcal L_n\) over the deterministic compact set \[
\{b:b_j=0,\ \|b\|\leq R\}.
\tag{D.6.2}
\] If needed, select the unique minimum-norm minimizer. This is a
measurable selector: it is the limit as \(\epsilon\downarrow0\) of the
unique minimizers after adding \(\epsilon\|b\|^2\). The definition uses
only \(X,Y,F_{-j}\). On \(\mathcal G_d\), its loss is at most \(\log2\),
so \(\|\bar b^{-j}\|\leq R_0<R\). It is then an unconstrained minimizer
in its coordinate subspace, and \[
\nabla\mathcal L_n(\bar b^{-j})=g_j^*e_j.
\tag{D.6.3}
\] The larger compact ball in the definition is essential: it keeps the
object well defined and independent of \(W_j\) even on exceptional
samples, while still giving interior stationarity on the good event.

For any of these restricted coefficient vectors \(b\), put \[
g_j(b)=\frac1n\sum_i
 [\sigma(h_{i,-j}^\top b_{-j})-Y_i]\tanh(W_j^\top X_i).
\tag{D.6.4}
\] Conditional on \(X,Y,F_{-j}\), the bracketed residual vector is
fixed, has norm at most \(\sqrt n\), and is independent of
\(W_j=G_j/\sqrt d\). As a function of standard Gaussian \(G_j\), (D.6.4)
is odd, has mean zero, and has Lipschitz constant at most \[
\frac{\|X\|_{\mathrm{op}}}{\sqrt{nd}}.
\tag{D.6.5}
\] Hence \[
\mathbb P(|g_j(b)|>u\mid X,Y,F_{-j})
\leq2\exp\!\left(-\frac{ndu^2}{2\|X\|_{\mathrm{op}}^2}\right).
\tag{D.6.6}
\] This conditioning is valid for the compact restricted minimizer and
for every deterministic restricted-GD time. Conditioning instead on a
good event involving \(W_j\) before applying concentration would not be
valid; we do not do so.

Choose a deterministic \(B_X\) with
\(\mathbb P(\|X\|_{\mathrm{op}}>B_X\sqrt d)\leq e^{-c_Xd}\) for large
\(d\); the same net proof permits \(B_X=8\sqrt{1+\alpha}\). For
\(d\geq3\), define \[
T_d=1+\left\lceil\frac{\max\{0,\log(R_0d)\}}{-\log q}\right\rceil,
\qquad
\tau_d=B_X\sqrt{\frac{8\log d}{n}}.
\tag{D.6.7}
\] In particular \(T_d=O(\log d)\) with a dimension-independent implied
constant. Applying (D.6.6) and a union bound over \(j\leq m\),
\(0\leq t<T_d\), and the restricted minimizers gives \[
\max_{j,t<T_d}|g_j(b^{-j,t})|\leq\tau_d,
\qquad \max_j|g_j^*|\leq\tau_d
\tag{D.6.8}
\] except with probability at most \[
e^{-c_Xd}+2m(T_d+1)d^{-4}=o(1).
\tag{D.6.9}
\] No independence between different leave-one-feature problems or
different times is asserted or required.

\hypertarget{d.7-uniform-coefficient-delocalization}{%
\subsection{D.7 Uniform coefficient
delocalization}\label{d.7-uniform-coefficient-delocalization}}

First compare the full and restricted MLEs on the intersection of
\(\mathcal G_d\) and (D.6.8). Strong monotonicity along their segment
gives \[
\mu\|\bar b^{-j}-b^*\|^2
\leq\langle\nabla\mathcal L_n(\bar b^{-j})-
 \nabla\mathcal L_n(b^*),\bar b^{-j}-b^*\rangle
=-g_j^*b_j^*.
\] Cauchy--Schwarz implies \[
|b_j^*|\leq\|\bar b^{-j}-b^*\|\leq|g_j^*|/\mu,
\quad\text{hence}\quad \|b^*\|_\infty\leq\tau_d/\mu.
\tag{D.7.1}
\]

For the trajectories, let \(G_d(b)=b-\eta_d\nabla\mathcal L_n(b)\).
Equation (D.6.1) says \[
b^{-j,t+1}=G_d(b^{-j,t})+\eta_dg_j(b^{-j,t})e_j.
\] The contraction established on the common ball gives \[
\|b^{t+1}-b^{-j,t+1}\|
\leq q_d\|b^t-b^{-j,t}\|+\eta_d|g_j(b^{-j,t})|.
\tag{D.7.2}
\] Starting from zero and using \(1-q_d=\eta_d\mu\), (D.6.8) implies \[
\max_{j,\ 0\leq t\leq T_d}\|b^t-b^{-j,t}\|
\leq\frac{\eta_d\tau_d}{1-q_d}=\frac{\tau_d}{\mu}.
\tag{D.7.3}
\] Since the restricted coordinate is zero, this bounds every
\(|b_j^t|\) through time \(T_d\). For every \(t\geq T_d\), (D.1.7) and
the definition of \(T_d\) give \[
\|b^t\|_\infty\leq\|b^*\|_\infty+R_0q^{T_d}
\leq\tau_d/\mu+d^{-1}.
\tag{D.7.4}
\] Thus on this single event, \[
\boxed{\quad\sup_{t\geq0}\|b^t\|_\infty
\leq\tau_d/\mu+d^{-1},\qquad \|b^*\|_\infty\leq\tau_d/\mu.\quad}
\tag{D.7.5}
\] Together with (D.4.4) and (D.6.9), this proves (D.1.8) with an
explicit failure bound. The proof handles infinitely many GD times by
deterministic contraction after a logarithmic horizon; it does not
union-bound infinitely many concentration events.

\hypertarget{d.8-positive-norms-and-simultaneous-calibration-floor}{%
\subsection{D.8 Positive norms and simultaneous calibration
floor}\label{d.8-positive-norms-and-simultaneous-calibration-floor}}

The lower-norm result is the argument in Appendix A Proposition 3; it is
included here to make its use explicit. Put \(K=HH^\top/m\),
\(\lambda_d=cd/n\), and \(R_Y=Y-1/2\). Then \[
\|b^1\|^2=\frac{\lambda_d^2}{m}R_Y^\top K R_Y.
\tag{D.8.1}
\] Conditional on \(X,F\), decompose \(R_Y=\bar r+\epsilon\), with
independent centered bounded \(\epsilon_i\), and
\(D_{ii}=\mathbb E(\epsilon_i^2\mid X,F)=\sigma'(\beta Z_i)\). A
quadratic-form expansion gives \[
\operatorname{Var}(R_Y^\top KR_Y\mid X,F)
\leq C\{\operatorname{tr}K^2+\|K\bar r\|^2\}
\leq Cn\|K\|_{\mathrm{op}}^2.
\tag{D.8.2}
\] The norm \(\|K\|\) is bounded in probability, so the normalization in
(D.8.1) makes its centered fluctuation \(o_P(1)\). The contribution
\(\bar r^\top K\bar r\) is nonnegative. Conditional on \(F\), the
bounded iid row variables \(D_{ii}K_{ii}\) have expectation tending to
\(v\rho_\beta\): uniformly over features, \(\|W_j\|^2\to1\) and
\(\theta^\top W_j\to0\), and the relevant bounded bivariate Gaussian
expectation is continuous. Hence \[
\frac1n\operatorname{tr}(KD)\to v\rho_\beta,
\qquad
\|b^1\|^2\geq\frac{c^2v\rho_\beta}{\alpha\gamma}+o_P(1).
\tag{D.8.3}
\]

Let \(g_0=\nabla\mathcal L_n(0)\). Descent and convexity respectively
give, for every \(t\geq1\), \[
\mathcal L_n(b^t)\leq\log2-\frac{\eta_d}{2}\|g_0\|^2,
\qquad
\mathcal L_n(b^t)\geq\log2+g_0^\top b^t.
\] Therefore \(\|b^t\|\geq\eta_d\|g_0\|/2=\|b^1\|/2\), simultaneously at
all positive times. Letting \(t\to\infty\) gives the same bound for
\(b^*\). Equation (D.8.3) proves (D.1.9)--(D.1.10).

For a fresh test pair, condition on the complete training environment
\(\mathcal E_d=(F,X_{1:n},Y_{1:n})\). Let \(S_{t,d}=h(X)^\top b^t\) and
\[
A_{t,d}=\mathbb E_X[ZS_{t,d}\mid\mathcal E_d]=g_F^\top b^t,
\qquad
Q_{t,d}=\mathbb E_X[S_{t,d}^2\mid\mathcal E_d]=(b^t)^\top\Sigma_Fb^t.
\tag{D.8.4}
\] Use \(t=\infty\) for \(b^*\). Equations (D.2.1), (D.1.9), and
Poincare yield, simultaneously for all \(t\in\{1,2,\ldots,\infty\}\), \[
Q_{t,d}-A_{t,d}^2\geq s r_*^2=:\Delta_*>0,
\qquad Q_{t,d}\leq L^2R_0^2=:Q_*.
\tag{D.8.5}
\] The readout derivative bounds are uniform in time: \[
\|\nabla S_t\|\leq LR_0,\quad
\|\nabla^2S_t\|_{\mathrm{HS}}\leq2L^2R_0,\quad
\|\nabla^2S_t\|_{\mathrm{op}}
\leq2L^2\sup_u\|b^u\|_\infty=o_P(1).
\tag{D.8.6}
\] The proof of Appendix C Theorem 2.1 is sequentially uniform over
classes with fixed bounds in (D.8.5)--(D.8.6). Otherwise a deterministic
sequence violating uniformity would contradict that theorem.
Consequently its random-environment argument gives \[
\sup_{t\in\{1,2,\ldots,\infty\}}
\left\|\mathcal L((Z,S_{t,d})\mid\mathcal E_d)
-N\!\left(0,\begin{pmatrix}1&A_{t,d}\\A_{t,d}&Q_{t,d}\end{pmatrix}\right)
\right\|_{\mathrm{TV}}\longrightarrow0
\quad\text{in probability}.
\tag{D.8.7}
\] This argument does not require deterministic overlap limits or a
union bound over time.

For \(Q>A^2\), define the explicit Gaussian calibration functional \[
\mathcal C_\beta(A,Q)
=\mathbb E_{U}\left|
\mathbb E_V\sigma\!\left(\beta\frac{A}{\sqrt Q}U
+\beta\sqrt{1-\frac{A^2}{Q}}V\right)
-\sigma(\sqrt Q U)\right|,
\qquad U,V\stackrel{\mathrm{iid}}\sim N(0,1).
\tag{D.8.8}
\] Let \(\mathrm{CE}_{t,d}\) be exact, unbinned population calibration
error of the predictor \(\sigma(S_{t,d})\), conditional on
\(\mathcal E_d\). Since sigmoid is injective, conditioning on its
probability is equivalent to conditioning on the score. the
signed-measure argument in Appendix C makes this calibration error
TV-Lipschitz. Applying the common label kernel \(Y\mid Z\) and (D.8.7),
\[
\sup_{t\in\{1,2,\ldots,\infty\}}
|\mathrm{CE}_{t,d}-\mathcal C_\beta(A_{t,d},Q_{t,d})|
\longrightarrow0
\quad\text{in probability}.
\tag{D.8.9}
\]

For clarity, strict positivity of (D.8.8) holds for every \(Q-A^2>0\).
If \(\beta=0\), the true curve is the constant \(1/2\) and the
nonconstant prediction has positive error. If \(\beta>0,A\leq0\), the
Gaussian conditional probability is nonincreasing while sigmoid is
strictly increasing. If \(\beta>0,A>0\), write its conditional curve as
\(g(s)=\mathbb E\sigma(ks+\tau V)\), with \(k=\beta A/Q>0\) and
\(\tau=\beta\sqrt{1-A^2/Q}>0\). Dominated convergence gives \[
\lim_{s\to\infty}e^{ks}(1-g(s))=e^{\tau^2/2}.
\] Equality with \(\sigma(s)\) for every \(s\) would force first \(k=1\)
and then \(e^{\tau^2/2}=1\), a contradiction. Equality almost everywhere
would imply equality everywhere by continuity. Thus the error is
strictly positive.

The parameter set \[
\mathcal K=\{(A,Q): Q\leq Q_*,\ Q-A^2\geq\Delta_*\}
\tag{D.8.10}
\] is compact. Formula (D.8.8) is continuous on it by dominated
convergence. Therefore \[
\varepsilon_*:=\min_{(A,Q)\in\mathcal K}\mathcal C_\beta(A,Q)>0.
\tag{D.8.11}
\] Equations (D.8.5) and (D.8.9) prove the simultaneous statement \[
\mathbb P\!\left(
\inf_{t\in\{1,2,\ldots,\infty\}}\mathrm{CE}_{t,d}
\geq\varepsilon_*/2\right)\longrightarrow1.
\tag{D.8.12}
\] In particular, \textbf{no training-measurable stopping rule taking a
positive integer time, even one diverging with dimension, attains
asymptotic calibration in this fixed-step regime.} The same holds if the
rule may output the MLE. At time zero the constant predictor \(1/2\) is
calibrated, since the symmetric teacher has marginal class probability
\(1/2\). The theorem does not cover steps \(c=c_d\to0\), which alter the
stated regime.

Positive absolute calibration error does not alone assert that every
confidence level is overconfident. The full curve in (D.8.8), or its
conditional form \(g(s)\), determines the direction. Early positive
iterates can be underconfident while still having strictly positive
exact calibration error.

\hypertarget{d.9-promotion-of-deterministic-finite-time-overlaps-to-all-times}{%
\subsection{D.9 Promotion of deterministic finite-time overlaps to all
times}\label{d.9-promotion-of-deterministic-finite-time-overlaps-to-all-times}}

Theorem A and Appendix B prove, for every fixed integer \(t\geq1\),
deterministic limits \[
A_{t,d}\xrightarrow{P}A_t(\alpha,\gamma,c,\beta),\qquad
Q_{t,d}\xrightarrow{P}Q_t(\alpha,\gamma,c,\beta).
\tag{D.9.1}
\] Then these limits extend to \(t=\infty\), and convergence is uniform
over all positive integer times and the MLE: \[
\sup_{t\in\{1,2,\ldots,\infty\}}
\big\{|A_{t,d}-A_t|+|Q_{t,d}-Q_t|\big\}\xrightarrow{P}0.
\tag{D.9.2}
\]

Here is the proof, including the exchange of limits. The population norm
bounds are \(\|g_F\|\leq L\) and \(\|\Sigma_F\|\leq L^2\). From (D.1.7),
for every \(t\geq T\), including infinity, \[
\|b^t-b^T\|\leq2R_0q^T,
\] \[
|A_{t,d}-A_{T,d}|\leq2LR_0q^T,
\qquad
|Q_{t,d}-Q_{T,d}|\leq4L^2R_0^2q^T.
\tag{D.9.3}
\] For each fixed \(t,T\), take probability limits in these
deterministic inequalities on events of probability tending to one. This
gives the same bound between \((A_t,Q_t)\) and \((A_T,Q_T)\). Therefore
the deterministic sequence is exponentially Cauchy and has a limit
\((A_\infty,Q_\infty)\). For arbitrary fixed \(T\), split the supremum
in (D.9.2) into the finite times below \(T\) and the common exponential
tail. Equation (D.9.1) controls the finite part, while (D.9.3) controls
both tails. Let \(d\to\infty\), then \(T\to\infty\). This proves (D.9.2)
without interchanging unproved limits.

Combining uniform continuity on (D.8.10), (D.8.9), and (D.9.2) gives the
deterministic whole-trajectory theorem \[
\boxed{\quad
\sup_{t\in\{1,2,\ldots,\infty\}}
\left|\mathrm{CE}_{t,d}
-\mathcal C_\beta\big(A_t(\alpha,\gamma,c,\beta),
Q_t(\alpha,\gamma,c,\beta)\big)\right|
\xrightarrow{P}0.\quad}
\tag{D.9.4}
\] For a data-dependent positive stopping rule \(\tau_d\), the same
uniform estimate holds after substituting \(t=\tau_d\). If
\(\tau_d\to\infty\) in probability, then
\(\mathrm{CE}_{\tau_d,d}\to\mathcal C_\beta(A_\infty,Q_\infty)>0\) in
probability. If the stopping time has no limiting value, the
deterministic curve still approximates its realized index and the common
positive floor still applies.

The finite-time result (D.9.1) is what makes (D.9.4) a calibration
function of the aspect ratios and training time. Equations
(D.8.7)--(D.8.12) themselves do not require that input and are already
unconditional consequences of the proofs above and the analytic transfer
theorem in Appendix C.

\clearpage

\hypertarget{appendix-e.-a-single-crossing-theorem-and-signed-overconfidence}{%
\section{Appendix E. A single-crossing theorem and signed
overconfidence}\label{appendix-e.-a-single-crossing-theorem-and-signed-overconfidence}}

Theorem, lemma, and proposition numbers in this appendix are local;
equation and section references carry the appendix letter.

\hypertarget{e.1-statement}{%
\subsection{E.1 Statement}\label{e.1-statement}}

Let \(Z,G\) be independent standard normal variables, let
\(k,\tau,\beta>0\), and put

\[
U=kZ+\tau G,\qquad V=k^2+\tau^2.
\tag{E.1}
\]

Given \(Z\), let \(Y\) be Bernoulli with mean \(\sigma(\beta Z)\),
independently of \(G\), where \(\sigma(x)=(1+e^{-x})^{-1}\). Its
conditional regression given the score coordinate is

\[
g(u):=\mathbb E[Y\mid U=u]
=\mathbb E_H\sigma(\lambda u+sH),\qquad
\lambda=\frac{\beta k}{V},\quad s=\frac{\beta\tau}{\sqrt V},
\tag{E.2}
\]

with \(H\) standard normal. Consider the positive-scale predictors
\(p_c(u)=\sigma(cu)\) and their population cross-entropy

\[
R(c)=\mathbb E\{\log(1+e^{cU})-cU g(U)\}.
\tag{E.3}
\]

\textbf{Theorem.} For every \(k,\tau,\beta>0\):

\begin{enumerate}
\def\labelenumi{\arabic{enumi}.}
\tightlist
\item
  The function \[
  h(u)=\frac{\operatorname{logit}g(u)}u,\qquad u>0,
  \tag{E.4}
  \] is continuous and strictly increasing, with endpoint limits \[
  h(0+)=4g'(0)=4\lambda\mathbb E\sigma'(sH),
  \qquad h(+\infty)=\lambda.
  \tag{E.5}
  \]
\item
  \(R\) has a unique global minimizer \(c_*\), and \[
  4g'(0)<c_*<\lambda,\qquad R(c_*)<R(0)=\log2.
  \tag{E.6}
  \]
\item
  There is a unique \(u_*>0\) such that \[
  \sigma(c_*u)-g(u)
  \begin{cases}>0,&0<u<u_*,\\=0,&u=u_*,\\<0,&u>u_*.
  \end{cases}
  \tag{E.7}
  \]
\item
  The average signed overconfidence is strictly positive: \[
  \operatorname{OC}(c_*):=
  \mathbb E\big[(\sigma(c_*U)-g(U))\operatorname{sign}U\big]>0.
  \tag{E.8}
  \]
\end{enumerate}

Here the predicted class is \(\mathbf1_{\{U>0\}}\), its confidence is
\(\sigma(c_*|U|)\), and its conditional probability of being correct is
\(g(|U|)\). With \(S=c_*U\), conditioning on \(U\) gives the equivalent
exact definition

\[
\operatorname{OC}(c_*)=\mathbb E[(\sigma(S)-\sigma(\beta Z))\operatorname{sign}S].
\]

Thus (E.8) is exactly mean confidence minus mean accuracy. It is not a
claim of overconfidence at every score: (E.7) proves underconfidence in
the far tails.

\hypertarget{e.2-logistic-noise-as-a-normal-variance-mixture}{%
\subsection{E.2 Logistic noise as a normal variance
mixture}\label{e.2-logistic-noise-as-a-normal-variance-mixture}}

Let \(E_1,E_2,\ldots\) be independent exponential variables of mean one,
and define

\[
W=2\sum_{j=1}^{\infty}\frac{E_j}{j^2}.
\tag{E.9}
\]

The series is finite almost surely, because its expectation is
\(2\sum j^{-2}=\pi^2/3<\infty\). If \(N\) is an independent standard
normal variable, the characteristic function of \(\sqrt W N\) is

\[
\mathbb E e^{-t^2W/2}
=\prod_{j=1}^{\infty}(1+t^2/j^2)^{-1}
=\frac{\pi t}{\sinh(\pi t)},
\tag{E.10}
\]

with the continuous value one at \(t=0\). The last identity is Euler's
product for \(\sinh\). On the other hand, the characteristic function of
a standard logistic variable \(L\) is

\[
\int_{-\infty}^{\infty}e^{itx}\frac{e^x}{(1+e^x)^2}\,dx
=\int_0^{\infty}\frac{r^{it}}{(1+r)^2}\,dr
=\Gamma(1+it)\Gamma(1-it)
=\frac{\pi t}{\sinh(\pi t)}.
\tag{E.11}
\]

The middle equality is the beta integral, and the last follows from the
gamma reflection identity. Uniqueness of characteristic functions
therefore proves

\[
L\ \overset d=\ \sqrt W N.
\tag{E.12}
\]

We also need strict support, not just a formal mixture. Write
\(W=2E_1+Q\), where \(Q=2\sum_{j\ge2}E_j/j^2\). For every \(w>0\),
\(\mathbb P(Q<w)>0\): choose a finite cutoff whose remaining tail has
expectation less than \(w/4\); the tail is below \(w/2\) with positive
probability by Markov's inequality, and the finite initial sum is below
\(w/2\) with positive probability, independently. Conditioning on \(Q\)
gives the density

\[
f_W(w)=\frac12 e^{-w/2}
\mathbb E\big[e^{Q/2}\mathbf1_{\{Q<w\}}\big]>0
\quad(w>0).
\tag{E.13}
\]

It is finite, since the integrand is bounded by \(e^{w/2}\). Hence the
CDF \(F_W\) is strictly increasing on \((0,\infty)\), is positive at
every positive argument, and vanishes at nonpositive arguments.

\hypertarget{e.3-a-direct-strict-single-crossing-argument}{%
\subsection{E.3 A direct strict single-crossing
argument}\label{e.3-a-direct-strict-single-crossing-argument}}

For \(s>0\), write

\[
F_s(x)=\mathbb E\sigma(x+sH),\qquad
\Lambda_s(x)=\operatorname{logit}F_s(x).
\tag{E.14}
\]

By symmetry of \(H\), \(F_s\) is the CDF of \(L+sH\). Conditional on
\(W\), this is a centered normal variable of variance \(W+s^2\).
Consequently

\[
F_s(x)=\mathbb E\Phi\left(\frac{x}{\sqrt{W+s^2}}\right),
\qquad
\sigma(cx)=\mathbb E\Phi\left(\frac{x}{\sqrt W/c}\right)
\quad(c>0).
\tag{E.15}
\]

Set \(x=e^t\). The difference of the CDFs of the logarithms of the two
random scales in (E.15) is

\[
N_c(u)
=F_W(e^{2u}-s^2)-F_W(c^2e^{2u}).
\tag{E.16}
\]

Integration by parts yields

\[
D_c(t):=F_s(e^t)-\sigma(ce^t)
=\int_{\mathbb R}K(t-u)N_c(u)\,du,
\qquad K(v)=e^v\varphi(e^v).
\tag{E.17}
\]

For completeness, the function being integrated against the signed scale
law is \(u\mapsto\Phi(e^{t-u})\); its derivative is \(-K(t-u)\). The
boundary terms vanish because \(N_c(u)\to0\) at both ends. The integral
is absolutely convergent: \(|N_c|\le1\) and \(\int K(v)\,dv=1/2\).

If \(c\ge1\), (E.13) and (E.16) give \(N_c(u)<0\) for every finite
\(u\); hence \(D_c(t)<0\) for every \(t\). In particular

\[
\tfrac12<F_s(x)<\sigma(x)\quad(x>0).
\tag{E.18}
\]

The first strict inequality follows from symmetry and strict
monotonicity of \(F_s\).

For \(0<c<1\), put

\[
u_c=\frac12\log\frac{s^2}{1-c^2}.
\tag{E.19}
\]

Strict monotonicity of \(F_W\) implies

\[
N_c(u)<0\ (u<u_c),\qquad
N_c(u_c)=0,\qquad
N_c(u)>0\ (u>u_c).
\tag{E.20}
\]

For any \(t_2>t_1\), the kernel ratio

\[
Q(u):=\frac{K(t_2-u)}{K(t_1-u)}
=e^{t_2-t_1}
\exp\left[-\frac{e^{2t_2}-e^{2t_1}}2e^{-2u}\right]
\tag{E.21}
\]

is strictly increasing in \(u\), and is bounded above by
\(e^{t_2-t_1}\). Therefore

\[
D_c(t_2)-Q(u_c)D_c(t_1)
=\int_{\mathbb R}K(t_1-u)
  [Q(u)-Q(u_c)]N_c(u)\,du>0.
\tag{E.22}
\]

Every factor except the bracket and \(N_c\) is positive; these last two
factors have the same strict sign on either side of \(u_c\). The
integral is finite by the bound on \(Q\) and the integrability of \(K\).
In particular, a zero of \(D_c\) is followed by strictly positive
values. This proves the exact direction of the required single crossing
without an appeal to a general variation-diminishing result.

Now fix \(0<x_1<x_2\), and choose

\[
c=\frac{\Lambda_s(x_1)}{x_1}\in(0,1),
\tag{E.23}
\]

where membership in \((0,1)\) follows from (E.18). Then
\(D_c(\log x_1)=0\). Equation (E.22) gives

\[
F_s(x_2)>\sigma(cx_2),\quad\text{so}\quad
\frac{\Lambda_s(x_2)}{x_2}>
\frac{\Lambda_s(x_1)}{x_1}.
\tag{E.24}
\]

This is strict monotonicity for every pair \(x_1<x_2\), without needing
an analyticity argument.

At zero, dominated differentiation and \(F_s(0)=1/2\) give

\[
\lim_{x\downarrow0}\frac{\Lambda_s(x)}x
=4\mathbb E\sigma'(sH)\in(0,1).
\tag{E.25}
\]

For the other endpoint,

\[
e^x[1-F_s(x)]
=\mathbb E\frac{e^{-sH}}{1+e^{-x-sH}}
\longrightarrow e^{s^2/2}
\tag{E.26}
\]

by dominated convergence, since \(e^{-sH}\) is integrable. Thus

\[
\Lambda_s(x)=x-\frac{s^2}2+o(1),\qquad
\lim_{x\to\infty}\frac{\Lambda_s(x)}x=1.
\tag{E.27}
\]

Applying (E.24)--(E.27) with \(x=\lambda u\) proves (E.4)--(E.5).

\hypertarget{e.4-risk-minimization-and-strict-signed-overconfidence}{%
\subsection{E.4 Risk minimization and strict signed
overconfidence}\label{e.4-risk-minimization-and-strict-signed-overconfidence}}

The normal law of \(U\) has all moments, so dominated differentiation
gives

\[
R'(c)=\mathbb E[U(\sigma(cU)-g(U))],\qquad
R''(c)=\mathbb E[U^2\sigma'(cU)]>0.
\tag{E.28}
\]

Since \(g(u)>1/2\) for \(u>0\), symmetry implies \(R'(0)<0\). On the
other hand,

\[
\lim_{c\to\infty}R'(c)
=\mathbb E[|U|(1-g(|U|))]>0.
\tag{E.29}
\]

The integrand is strictly positive almost surely and is integrable. Thus
the strictly convex function \(R\) has a unique critical point
\(c_*>0\), and it is its global minimizer. Its value is strictly below
\(R(0)=\log2\).

Let \(h_0=4g'(0)\). By the strict increase already proved,
\(h_0<h(u)<\lambda\) for every \(u>0\). Hence \(R'(c)<0\) for
\(c\le h_0\), and \(R'(c)>0\) for \(c\ge\lambda\). This proves
\(h_0<c_*<\lambda\). Continuity, strict increase and the endpoint limits
produce a unique \(u_*>0\) with \(h(u_*)=c_*\), and (E.7) follows from
the strict increase of \(\sigma\).

Write \(e(u)=\sigma(c_*u)-g(u)\), and let \(\varphi_V\) denote the
\(N(0,V)\) density. The first-order condition and oddness of \(e\) give

\[
0=R'(c_*)=2\int_0^\infty u\,e(u)\varphi_V(u)\,du.
\tag{E.30}
\]

Subtracting this identity divided by \(u_*\) from the unweighted signed
error gives

\[
\begin{aligned}
\operatorname{OC}(c_*)
&=2\int_0^\infty e(u)\varphi_V(u)\,du\\
&=2\int_0^\infty e(u)
    \left(1-\frac u{u_*}\right)\varphi_V(u)\,du
>0.
\end{aligned}
\tag{E.31}
\]

The two factors \(e(u)\) and \(1-u/u_*\) have the same strict sign for
every \(u>0\) other than \(u_*\). Integrability follows from \(|e|\le1\)
and the first absolute normal moment. This completes the theorem.

One may compute \(c_*\) through a one-dimensional strictly monotone
equation. Stein's identity and the original representation (E.1) give

\[
Vc_*\,\mathbb E\sigma'(c_*\sqrt V\,H)
=k\beta\,\mathbb E\sigma'(\beta H).
\tag{E.32}
\]

Equation (E.28), rather than any presumed monotonicity of the individual
factors in (E.32), proves uniqueness.

\hypertarget{e.5-consequence-for-the-random-feature-model}{%
\subsection{E.5 Consequence for the random-feature
model}\label{e.5-consequence-for-the-random-feature-model}}

Suppose a first-step Gaussian random-feature theorem gives, conditional
on the trained model,

\[
\mathcal L(Z,U_d\mid\text{model})
\longrightarrow \mathcal L(Z,kZ+\tau G)
\quad\text{in total variation in probability},
\tag{E.33}
\]

with \(k,\tau>0\), and gives population cross-entropy convergence at
every fixed deterministic scale. Taking the fixed scale \(c=c_*\) from
this appendix then yields

\[
\operatorname{OC}_d(c_*)\xrightarrow{\mathbb P}
\operatorname{OC}(c_*)>0,\qquad
R_d(c_*)\xrightarrow{\mathbb P}R(c_*)<\log2.
\tag{E.34}
\]

Indeed signed OC is the conditional expectation of the bounded
measurable function

\[
(z,u)\longmapsto
[\sigma(c_*u)-\sigma(\beta z)]\operatorname{sign}u,
\tag{E.35}
\]

so total variation proves the first assertion directly. A second moment
bound or the separately established risk-transfer theorem handles the
unbounded cross-entropy integrand. In particular, fixed positive margins
below \(\log2\) and above zero hold simultaneously with probability
tending to one. Theorem A and Appendix C supply these total-variation
and risk limits.

Because \(u\mapsto\sigma(c_*u)\) is one-to-one, the limiting exact
unbinned calibration error is

\[
\mathbb E|\sigma(c_*U)-g(U)|
\ge\operatorname{OC}(c_*)>0.
\tag{E.36}
\]

The finite-model OC statement (E.34) itself does not require continuity
of conditional regression: it uses the bounded test function (E.35).
This is weaker in prerequisite than proving convergence of an entire
conditional calibration curve.

The theorem treats all positive parameters in the Gaussian scalar
experiment, not a selected categorical surrogate or one numerical
parameter choice. It closes a signed-overconfidence proof obligation at
a risk-optimal \textbf{first-step scale}. It does not establish a
deterministic arbitrary-time neural-GD trajectory, interchange a
training-time limit with dimension, optimize over all stopping times, or
show that the scalar risk optimum is a global optimum over the full GD
path. Those remain separate the present problem obligations. The
limiting predictor is overconfident near the boundary and underconfident
in the tails, so a claim of pointwise overconfidence at all scores would
be false.

\hypertarget{e.6-an-explicit-common-regime-with-the-whole-trajectory-stopping-theorem}{%
\subsection{E.6 An explicit common regime with the whole-trajectory
stopping
theorem}\label{e.6-an-explicit-common-regime-with-the-whole-trajectory-stopping-theorem}}

The risk-optimal scale above is not automatically a stable GD step.
Nevertheless, there is an explicit nonempty regime in which it satisfies
the small-step and sample-size conditions of the entire-trajectory
theorem in Appendix D. Thus the overconfidence and stopping conclusions
need not be about incompatible parameter choices.

Use the actual first-step law from the Gaussian tanh RF calculation,
with the step scale factored out: \[
S_1=cU,\qquad
k=a^2\mu_\beta,\qquad
\tau^2=\frac{\mu_\beta^2a^2v}{\gamma}
+\frac{a^4+v^2/\gamma}{4\alpha},
\tag{E.37}
\] where \[
\begin{gathered}
a=\mathbb E[H\tanh H]>0,\quad v=\mathbb E\tanh^2H>0,\quad \nu=v-a^2>0,\\
\mu_\beta=\mathbb E[H(\sigma(\beta H)-1/2)]=\beta\mathbb E\sigma'(\beta H).
\end{gathered}
\] These are exactly the coefficients in Appendix B, equation (B.21),
and in equation (11) of the main text. For \(\beta>0\), both \(k\) and
\(\tau\) are strictly positive.

Fix any \(\gamma>0\), and use the constants of Appendix D \[
L=8\sqrt{1+\gamma},\qquad
s_0=\nu/2,\qquad
\kappa=\sqrt{s_0/2},\qquad
p_0=\frac{s_0^2}{20L^4},\qquad
B=\sqrt{2\log(4/p_0)}.
\] Choose \[
\frac{\alpha}{\gamma}>
C_0:=\max\left\{1,\left(\frac{256L}{\kappa p_0}\right)^2\right\},
\qquad
0<\beta<
\min\left\{\frac{\log3}{B},\frac{v}{8aL\sqrt\alpha}\right\}.
\tag{E.38}
\] Every interval here has positive length. In particular this is a
nonempty open set of proportional sample, width and teacher parameters,
although its sample ratio is deliberately very conservative.

The first inequality on \(\beta\) gives \(q_\beta=\sigma(-\beta B)>1/4\)
and hence \(p_1=q_\beta p_0/2>p_0/8\) in Appendix D. Thus its threshold
satisfies \[
C_*(\beta)=
\max\left\{1,\left(\frac{32L}{\kappa p_1}\right)^2\right\}
\le C_0<\alpha/\gamma.
\tag{E.39}
\] Next, \(\mu_\beta\le\beta/4\) and (E.37) give \[
0<c_*<\frac{\beta k}{k^2+\tau^2}
\le\frac{a^2\alpha\gamma\beta^2}{v^2}
<\frac{\gamma}{64L^2}
<\frac{\gamma}{32L^2}
=\frac{\gamma}{2\Lambda},
\qquad\Lambda=16L^2.
\tag{E.40}
\] This is precisely inside the stable fixed-step interval of Appendix
D. The scalar risk optimum is therefore a valid fixed deterministic GD
parameter in the same model for which Appendix D proves its all-time
statements.

Consequently, Theorems A and B imply that this one model simultaneously
has:

\begin{itemize}
\tightlist
\item
  A first positive iterate with population risk strictly below \(\log2\)
  and strictly positive signed overconfidence, with probability tending
  to one.
\item
  A fixed positive exact-calibration floor, with probability tending to
  one, simultaneously over every positive integer training time, the
  MLE, and any training-measurable stopping rule taking those values.
\item
  A calibrated time-zero predictor with risk exactly \(\log2\).
\end{itemize}

This also gives a nonvacuous risk-constrained stopping conclusion. For
any fixed \(\Delta>0\) and any training-measurable rule
\(T_d\in\{0,1,\ldots,\infty\}\) such that \[
\mathbb P\{R_d(T_d)\le\log2-\Delta\}\longrightarrow1,
\tag{E.41}
\] one necessarily has \(\mathbb P(T_d>0)\to1\), because the risk at
time zero is exactly \(\log2\). The simultaneous positive calibration
floor then applies to \(T_d\). On the other hand, a fixed first-step
rule satisfies (E.41) for every \(\Delta<R(0)-R(c_*)\) by (E.34). Thus
such risk improvement is attainable but cannot coexist with asymptotic
exact calibration through stopping alone in this fixed-step regime.

This common regime still has sufficiently large \(n/m\), so the fitted
output layer is underparameterized. The random-feature architecture
remains the original dense Gaussian nonlinear model; no categorical
design is introduced. The argument does not extend the all-time
statement to every overparameterized sample ratio or to
dimension-dependent vanishing steps.

\clearpage

\hypertarget{appendix-f.-all-time-obstruction-in-the-strongly-overparameterized-regime}{%
\section{Appendix F. All-time obstruction in the strongly
overparameterized
regime}\label{appendix-f.-all-time-obstruction-in-the-strongly-overparameterized-regime}}

Theorem, lemma, and proposition numbers in this appendix are local;
equation and section references carry the appendix letter.

\hypertarget{f.1-model-and-explicit-nonempty-regime}{%
\subsection{F.1 Model and explicit nonempty
regime}\label{f.1-model-and-explicit-nonempty-regime}}

Use the exact original Gaussian tanh model: \[
X_i\sim N(0,I_d),\quad W_j\sim N(0,I_d/d),\quad
Y_i\mid X_i\sim\operatorname{Bernoulli}(\sigma(\beta\theta^\top X_i)),
\] with independent inputs and feature rows, conditionally independent
labels, \(\|\theta\|=1\), and fixed finite \(\beta\ge0\). Let \[
n/d\to\alpha>0,\qquad m/d\to\gamma>0,\quad
h(x)=\tanh(Fx),\quad H_{ij}=h_j(X_i).
\] Only the output coefficients are trained. Put \(s_i=2Y_i-1\),
\(\ell(z)=\log(1+e^{-z})\), and \[
\mathcal L_n(b)=\frac1n\sum_i\ell(s_i h_i^\top b),\qquad
b^0=0,\quad b^{t+1}=b^t-\eta_d\nabla\mathcal L_n(b^t),
\quad \eta_d=\frac{cd}{m}.
\tag{F.1.1}
\] Thus the actual update coefficient multiplying
\(H^\top[Y-\sigma(Hb^t)]\) is \(cd/(nm)\), as in the main theorem. The
scale \(c>0\) is fixed independently of dimension.

Define \[
a=\mathbb E[G\tanh G],\quad v=\mathbb E\tanh^2G,\quad
\nu=v-a^2>0,\quad s=\nu/2,\quad \kappa=\sqrt{s/2},
\quad G\sim N(0,1).
\] For the training-input side, set \[
L_X=8\sqrt{1+\alpha},\quad p_X=\frac{s^2}{20L_X^4},\quad
C_X=\max\left\{1,\left(\frac{32L_X}{\kappa p_X}\right)^2\right\},
\] \[
k_0=\frac{\kappa^2p_X}{16},\qquad k_1=64L_X^2.
\tag{F.1.2}
\] The explicit assumptions for this module are \[
\boxed{\qquad\gamma/\alpha>C_X,\qquad 0<c<2\alpha/k_1.\qquad}
\tag{F.1.3}
\] For each fixed \(\alpha>0\) and finite \(\beta\), these conditions
leave nonempty open intervals of width ratios and learning scales. They
imply \(m>n\) for large dimension. Their conservative constants are not
claimed to be sharp.

For the fresh-input feature side, define \[
L_F=8\sqrt{1+\gamma},\quad p_F=\frac{s^2}{20L_F^4},\quad
B_F=\sqrt{2\log(4/p_F)},\quad q_F=\sigma(-\beta B_F)>0,
\] \[
C_F=\log\frac4{q_Fp_F},\quad
R_F=\frac{C_F}{\kappa},\quad
\varepsilon_F=\frac{q_Fp_F}{4},\quad
a_F=\frac{q_F\kappa p_F}{2}>0.
\tag{F.1.4}
\] Finally put \[
D=2\sqrt{\gamma k_1/\alpha},\qquad
\mathfrak B_t=\frac1{t+1}+\frac{\alpha}{ck_0}\frac{\log^2(t+1)}t
\quad(t\ge1).
\tag{F.1.5}
\] These are all deterministic dimension-independent constants or
sequences.

\hypertarget{f.2-the-simultaneous-design-event-proved-for-the-actual-nonlinear-features}{%
\subsection{F.2 The simultaneous design event, proved for the actual
nonlinear
features}\label{f.2-the-simultaneous-design-event-proved-for-the-actual-nonlinear-features}}

With probability tending to one, all of the following hold: \[
\|F\|_{\mathrm{op}}\le L_F,\qquad
\Sigma_F:=\mathbb E_X[h(X)h(X)^\top]\succeq sI_m,
\tag{F.2.1}
\] \[
k_0I_n\preceq HH^\top/m\preceq k_1I_n.
\tag{F.2.2}
\] We prove these facts directly for the actual nonlinear features.

\hypertarget{f.2.1-population-covariance-on-either-side}{%
\subsubsection{F.2.1 Population covariance on either
side}\label{f.2.1-population-covariance-on-either-side}}

Consider a matrix \(V\) with \(r\) independent rows \(N(0,I_d/d)\),
where \(r/d\) tends to a fixed positive constant. For an independent
\(G\sim N(0,I_d)\), write \(h_V(G)=\tanh(VG)\). If \(q_j=\|V_j\|^2\) and
\(A(q)=\mathbb E\tanh'(\sqrt qG_1)\), the Hermite expansion gives \[
\mathbb E[h_Vh_V^\top]
=D_AVV^\top D_A+D_\nu+E,
\tag{F.2.3}
\] where \[
(D_\nu)_{jj}=\mathbb E\tanh^2(\sqrt{q_j}G_1)-q_jA(q_j)^2,
\quad E_{jj}=0,
\quad |E_{ij}|\le|\rho_{ij}|^3\ (i\ne j).
\] Here \(\rho_{ij}\) is the inner product of the two normalized rows.
The linear Hermite coefficient is identified by Gaussian integration by
parts; all remaining degrees are odd and at least three. Cauchy--Schwarz
on their coefficients gives the displayed remainder bound, since the
second moment of tanh is at most one. Independent Gaussian row
directions satisfy \[
\mathbb E\|E\|_F^2
\le\frac{15r(r-1)}{d(d+2)(d+4)}=O(d^{-1}).
\tag{F.2.4}
\] Uniform row-norm concentration makes \(D_\nu\to\nu I\) in operator
norm, while (F.2.4) makes \(\|E\|_{\mathrm{op}}\to0\) in probability.
Thus the covariance is at least \(sI\) with probability tending to one.
Gaussian sphere nets also give
\(\|V\|_{\mathrm{op}}\le8\sqrt{1+r/d+o(1)}\) with exponential
probability; the stated constants \(L_F,L_X\) have ample fixed slack.
More explicitly, the covariance event can be chosen with failure
\(O(d^{-1})+O(de^{-c_0d})\).

Apply this first to \(V=F\), obtaining (F.2.1). Apply it also to
\(V=X/\sqrt d\), obtaining an event \(\mathcal X_d\), measurable in the
training inputs alone, on which \[
\|X\|_{\mathrm{op}}/\sqrt d\le L_X,\qquad
K_X:=\mathbb E_G[\tanh(XG/\sqrt d)\tanh(XG/\sqrt d)^\top]
\succeq sI_n.
\tag{F.2.5}
\] This second use is legitimate because the rows \(X_i/\sqrt d\) are
iid Gaussian. It does not replace the actual feature matrix by
independent entries.

\hypertarget{f.2.2-conditional-independent-columns-and-small-balls}{%
\subsubsection{F.2.2 Conditional independent columns and small
balls}\label{f.2.2-conditional-independent-columns-and-small-balls}}

Condition on \(X\in\mathcal X_d\). The \(m\) columns \(h_j\) of \(H\)
are iid centered vectors of the form \(\tanh(XG_j/\sqrt d)\). For every
unit \(u\in\mathbb R^n\), \(u^\top h_j\) is a centered \(L_X\)-Lipschitz
function of a standard Gaussian vector. Consequently \[
\mathbb E e^{z u^\top h_j}\le e^{z^2L_X^2/2},\quad
\mathbb E(u^\top h_j)^2\ge s,\quad
\mathbb E(u^\top h_j)^4\le5L_X^4.
\tag{F.2.6}
\] The fourth-moment bound follows by Gaussian Poincare applied to the
function and its square. Paley--Zygmund then gives \[
\mathbb P(|u^\top h_j|\ge\kappa\mid X)\ge p_X.
\tag{F.2.7}
\]

Let \(\psi(z)=0\) for \(|z|\le\kappa/2\), let it rise linearly to one on
\(\kappa/2<|z|<\kappa\), and let it equal one for \(|z|\ge\kappa\). It
is \(2/\kappa\)-Lipschitz, bounded by one, and vanishes at zero.
Symmetrization, Rademacher contraction, and
\(\operatorname{tr}K_X\le nL_X^2\) give \[
\mathbb E\sup_{\|u\|=1}\left|
\frac1m\sum_j\psi(u^\top h_j)-\mathbb E\psi(u^\top h_j)
\right|
\le\frac{8L_X}{\kappa}\sqrt{\frac nm}.
\tag{F.2.8}
\] If \(m/n\ge C_X\), this is at most \(p_X/4\). Bounded differences
with deviation \(p_X/4\) proves, except with conditional probability
\(e^{-mp_X^2/8}\), \[
\inf_{\|u\|=1}\frac1m\#\{j:|u^\top h_j|>\kappa/2\}\ge p_X/2.
\tag{F.2.9}
\] In fact (F.2.9) is robust: after deleting any subset of at most
\(p_Xm/4\) columns, the remaining Gram matrix, still normalized by
\(m\), is at least \(\kappa^2p_XI/16=k_0I\). No independence of the
deleted columns or union bound over subsets is needed. In particular it
proves the lower half of (F.2.2).

The same two-sphere-net bound applied to \(H^\top\) gives \[
\|H\|_{\mathrm{op}}\le8L_X\sqrt m
\tag{F.2.10}
\] with conditional failure at most \(2e^{-3m}\) when \(m\ge n\): each
fixed bilinear form is centered \(L_X\)-subGaussian, and \(1/4\)-nets
have cardinalities at most \(9^m,9^n\). This proves the upper half of
(F.2.2).

Let \(\mathcal G_d\) combine (F.2.1), (F.2.5), and (F.2.2). For large
\(d\), (F.1.3) ensures \(m/n\ge C_X\), and \[
\mathbb P(\mathcal G_d^c)
\le O(d^{-1})+O(de^{-c_0d})+e^{-mp_X^2/8}+2e^{-3m}=o(1).
\tag{F.2.11}
\] No independence between the good feature event and the conditional
sample-Gram event is asserted: a union bound suffices. Labels play no
role in this design event.

\hypertarget{f.3-uniform-empirical-loss-decay}{%
\subsection{F.3 Uniform empirical loss
decay}\label{f.3-uniform-empirical-loss-decay}}

On \(\mathcal G_d\), the full row rank in (F.2.2) makes every binary
labeling strictly separable. For \(R>0\), the explicit comparator \[
b_R=H^\top(HH^\top)^{-1}(R s_{1:n})
\tag{F.3.1}
\] has every signed training margin equal to \(R\), and \[
\mathcal L_n(b_R)=\ell(R)\le e^{-R},\qquad
\|b_R\|^2\le\frac{nR^2}{mk_0}.
\tag{F.3.2}
\]

The empirical Hessian is globally bounded by \[
\nabla^2\mathcal L_n(b)\preceq\frac{H^\top H}{4n},
\quad\text{so its smoothness constant is at most }\frac{mk_1}{4n}.
\] For large \(d\), \(n/d\ge\alpha/2\); assumption (F.1.3) gives \[
\eta_d\frac{mk_1}{4n}=\frac{cdk_1}{4n}<1.
\tag{F.3.3}
\] Hence ordinary GD decreases the empirical loss at every step.

For clarity, the convex smooth-GD comparator inequality used here is \[
\mathcal L_n(b^t)-\mathcal L_n(u)\le\frac{\|u\|^2}{2\eta_dt}
\quad(t\ge1).
\tag{F.3.4}
\] It follows directly by expanding \(\|b^{k+1}-u\|^2\), using
convexity, and using descent to replace
\(\eta_d^2\|\nabla\mathcal L_n(b^k)\|^2\) by
\(2\eta_d[\mathcal L_n(b^k)-\mathcal L_n(b^{k+1})]\). Summing gives the
bound on the average of losses at times \(1,\ldots,t\), and monotonicity
bounds the final loss by that average. No strong convexity or finite MLE
is needed.

Use (F.3.2) in (F.3.4), then take \(R=\log(t+1)\). Since
\(n/d\le2\alpha\) eventually, \[
\boxed{\quad
\mathcal L_n(b^t)
\le\frac1{t+1}+\frac{n}{2cdk_0}\frac{\log^2(t+1)}t
\le\mathfrak B_t\qquad\text{for every }t\ge1.
\quad}
\tag{F.3.5}
\] This inequality holds simultaneously at all integer times on one
event. In particular it applies to times growing arbitrarily quickly
with dimension.

\hypertarget{f.4-uniform-norm-growth-and-a-deterministic-crossing-time}{%
\subsection{F.4 Uniform norm growth and a deterministic crossing
time}\label{f.4-uniform-norm-growth-and-a-deterministic-crossing-time}}

For any \(b\), Jensen's inequality, monotonicity of \(\ell\), and
Cauchy--Schwarz give \[
\mathcal L_n(b)
\ge\ell\!\left(\frac1n\sum_i s_i h_i^\top b\right)
\ge\ell\!\left(\frac{\|Hb\|}{\sqrt n}\right)
\ge\ell(D\|b\|).
\tag{F.4.1}
\] The last inequality uses \(\|H\|^2/m\le k_1\) and
\(m/n\le4\gamma/\alpha\) for large \(d\). Thus (F.3.5) implies \[
\boxed{\quad
\|b^t\|\ge\frac1D
\left[-\log(e^{\mathfrak B_t}-1)\right]_+
\quad(t\ge1).\quad}
\tag{F.4.2}
\] Its right-hand side grows as \(D^{-1}(\log t-2\log\log t-O(1))\). The
path need not have monotone norm; the simultaneous lower bound is what
is used.

For any deterministic threshold \(R_*>0\), define the explicit finite
integer \[
T(R_*)=\min\{T\ge2:\mathfrak B_T<\ell(DR_*)\}.
\tag{F.4.3}
\] It exists because \(\mathfrak B_T\to0\). At this time the loss is
below \(\ell(DR_*)\); monotone empirical loss and (F.4.1) then give \[
\|b^t\|>R_*\qquad\text{for every integer }t\ge T(R_*).
\tag{F.4.4}
\] The time \(T(R_*)\) depends only on the fixed model parameters and
the chosen threshold, never on dimension or the data. This is stronger
than taking \(t\to\infty\) separately for each fixed dataset.

\hypertarget{f.5-a-calibration-floor-for-every-sufficiently-large-readout}{%
\subsection{F.5 A calibration floor for every sufficiently large
readout}\label{f.5-a-calibration-floor-for-every-sufficiently-large-readout}}

We prove a large-norm calibration witness directly. On (F.2.1), every
coefficient vector \(b\), including arbitrary training-dependent
vectors, satisfies \[
\mathbb E_X S_b^2\ge s\|b\|^2,\qquad
\mathbb E_X S_b^4\le5L_F^4\|b\|^4,
\quad S_b=b^\top h(X).
\] Consequently \[
\mathbb P_X(|S_b|\ge\kappa\|b\|)\ge p_F
\quad\text{for every }b\ne0.
\tag{F.5.1}
\] All constants and inequalities are simultaneous over \(b\). This
statement follows from uniform population moment bounds; it does not
union-bound over coefficient vectors.

Write \(Z=\theta^\top X\), \(p(Z)=\sigma(\beta Z)\). The choice of
\(B_F\) ensures \(\mathbb P(|Z|>B_F)\le p_F/2\), and on the
complementary event either label has probability at least \(q_F\).
Define the bounded measurable score witness \[
f(S)=\operatorname{sign}(S)\mathbf1_{\{|S|\ge C_F\}}.
\] For every score and teacher coordinate, \[
[\sigma(S)-p(Z)]f(S)
\ge q_F\mathbf1_{\{|S|\ge C_F,\ |Z|\le B_F\}}
   -e^{-C_F}\mathbf1_{\{|S|\ge C_F\}}.
\tag{F.5.2}
\] Indeed, on a positive score the signed difference is at least
\(1-p(Z)-e^{-C_F}\), and on a negative score it is at least
\(p(Z)-e^{-C_F}\). Outside the bounded-teacher event the remaining lower
bound is still \(-e^{-C_F}\).

If \(\|b\|\ge R_F=C_F/\kappa\), (F.5.1) implies
\(\mathbb P(|S_b|\ge C_F,|Z|\le B_F)\ge p_F/2\). Hence \[
\mathbb E_X[(\sigma(S_b)-p(Z))f(S_b)]
\ge q_Fp_F/2-e^{-C_F}\ge\varepsilon_F.
\tag{F.5.3}
\] By the defining signed-measure characterization of exact calibration
error, or just conditional expectation and \(|f|\le1\), \[
\mathrm{CE}(b)
=\mathbb E_X\left|\mathbb E[p(Z)\mid S_b]-\sigma(S_b)\right|
\ge\left|\mathbb E_X[(\sigma(S_b)-p(Z))f(S_b)]\right|.
\] Therefore \[
\boxed{\qquad \inf_{\|b\|\ge R_F}\mathrm{CE}(b)\ge\varepsilon_F>0.\qquad}
\tag{F.5.4}
\] This holds on the good feature event for all readouts at once. It
does not require their coefficients to be delocalized or their fresh
score to be approximately Gaussian. The bounded-teacher and score-tail
events may be dependent; only the union bound was used.

\hypertarget{f.6-all-positive-times-and-arbitrary-stopping-rules}{%
\subsection{F.6 All positive times and arbitrary stopping
rules}\label{f.6-all-positive-times-and-arbitrary-stopping-rules}}

Set \(T_F=T(R_F)\), a deterministic finite integer. On \(\mathcal G_d\),
(F.4.4) and (F.5.4) give \[
\inf_{t\ge T_F}\mathrm{CE}_{t,d}\ge\varepsilon_F.
\tag{F.6.1}
\] For the finitely many earlier times, apply \textbf{Theorem A}. It
applies to every fixed finite horizon, every fixed \(c>0\), and all
finite positive \(\alpha,\gamma\). It proves deterministic limits \[
\mathrm{CE}_{t,d}\xrightarrow{P}C_\beta(A_t,Q_t)>0
\qquad(1\le t<T_F).
\tag{F.6.2}
\] The positive-definite residual covariance in that theorem is valid
also in the present separable regime. No finite-MLE result is imported.

Define \[
\varepsilon_{\rm early}
=\frac12\min_{1\le t<T_F}C_\beta(A_t,Q_t)>0,
\qquad
\varepsilon_* =\min\{\varepsilon_F,\varepsilon_{\rm early}\}>0.
\tag{F.6.3}
\] The minimum is over a nonempty finite set because \(T_F\ge2\). A
finite union bound in (F.6.2), combined with (F.2.11) and (F.6.1),
proves the new theorem: \[
\boxed{\qquad
\mathbb P\!\left\{\inf_{t\in\{1,2,\ldots\}}
\mathrm{CE}_{t,d}\ge\varepsilon_*\right\}\longrightarrow1.
\qquad}
\tag{F.6.4}
\]

In particular, for \textbf{every} sequence of training-measurable
stopping rules \(\tau_d\in\{1,2,\ldots\}\), with no upper bound on their
size, \[
\mathbb P(\mathrm{CE}_{\tau_d,d}\ge\varepsilon_*)\longrightarrow1.
\tag{F.6.5}
\] The event in (F.6.4) is common to all rules; the order of its
quantifiers is stronger than proving (F.6.5) one rule at a time. Extra
randomization of the rule does not affect this pathwise implication.
Time zero remains the calibrated constant predictor. A finite
unregularized MLE does not exist in this separable regime, so no such
object is claimed here.

\hypertarget{f.7-population-risk-divergence-and-a-finite-optimal-stopping-window}{%
\subsection{F.7 Population-risk divergence and a finite optimal-stopping
window}\label{f.7-population-risk-divergence-and-a-finite-optimal-stopping-window}}

The same feature event yields a useful stronger conclusion about risk.
Conditional on \(X\), logistic cross-entropy is at least \[
\min\{p(Z),1-p(Z)\}|S_b|.
\] Use the joint event \(\{|S_b|\ge\kappa\|b\|,|Z|\le B_F\}\), whose
probability is at least \(p_F/2\). Uniformly for all coefficients, \[
\mathcal R_F(b):=\mathbb E[\ell(Y,S_b)\mid F,b]
\ge a_F\|b\|,
\qquad
\mathcal R_F(b)\le\log2+L_F\|b\|.
\tag{F.7.1}
\] In the lower bound, \(\ell(Y,S)=\log(1+e^S)-YS\); this is the same
loss as \(\ell((2Y-1)S)\) in Section F.1. The upper bound follows from
\(\ell(Y,S)\le\log2+|S|\) and Gaussian Poincare. If \(b\) is trained on
data, both expectations are fresh-test population quantities conditioned
on that training environment.

Combining (F.7.1) and (F.4.2), on one event of probability tending to
one, \[
\mathcal R_{t,d}\ge\frac{a_F}{D}
\left[-\log(e^{\mathfrak B_t}-1)\right]_+
\quad\text{for every }t\ge1.
\tag{F.7.2}
\] This proves uniform late-time divergence: for each fixed \(M>0\), a
deterministic finite time exists after which \(\mathcal R_{t,d}>M\)
simultaneously for all later times with probability tending to one. For
every \(\tau_d\to\infty\) in probability,
\(\mathcal R_{\tau_d,d}\to\infty\) in probability. Empirical risk
nevertheless tends to zero by (F.3.5).

In particular, let \[
R_{\rm risk}=(\log2+1)/a_F,\qquad
T_{\rm risk}=T(R_{\rm risk}).
\tag{F.7.3}
\] Then \(\mathcal R_{t,d}>\log2+1\) for all \(t\ge T_{\rm risk}\) on
\(\mathcal G_d\), whereas \(\mathcal R_{0,d}=\log2\). Therefore the
population-risk optimum over the whole GD path is attained before this
fixed time: \[
\inf_{t\ge0}\mathcal R_{t,d}
=\min_{0\le t<T_{\rm risk}}\mathcal R_{t,d}
\quad\text{with probability tending to one}.
\tag{F.7.4}
\] Finite-time risk convergence from the main theorem now gives the
deterministic characterization \[
\inf_{t\ge0}\mathcal R_{t,d}
\xrightarrow{P}
\min_{0\le t<T_{\rm risk}}\mathcal R_\beta(A_t,Q_t),
\tag{F.7.5}
\] where the time-zero value is \(\log2\). Thus the fixed-time recursion
suffices to compute optimal population-risk stopping in this regime: no
unproved exchange of infinite time and dimension is used. If the finite
limiting minimum has a unique index, a population-risk-minimizing index
equals it with probability tending to one. Without uniqueness, the
limiting minimizing set remains the correct conclusion.

\hypertarget{f.8-nontrivial-learning-and-the-calibration-obstruction-can-coexist}{%
\subsection{F.8 Nontrivial learning and the calibration obstruction can
coexist}\label{f.8-nontrivial-learning-and-the-calibration-obstruction-can-coexist}}

For \(\beta>0\), retain the exact first-step constants \[
\mu_\beta=\mathbb E[G(\sigma(\beta G)-1/2)]>0,
\quad k=a^2\mu_\beta,
\] \[
\tau^2=\frac{\mu_\beta^2a^2v}{\gamma}
+\frac{a^4+v^2/\gamma}{4\alpha},\qquad V=k^2+\tau^2.
\tag{F.8.1}
\] The limiting first-step score is \(cU\) with \(U=kZ+\tau G\). Its
risk \(R(c)\) has \(R'(0)=-a^2\mu_\beta^2\) and \(R''(c)\le V/4\). Hence
\[
R(c)\le\log2-ca^2\mu_\beta^2+c^2V/8.
\] Choosing a fixed \[
0<c<\min\left\{2\alpha/k_1,\ 4a^2\mu_\beta^2/V\right\}
\tag{F.8.2}
\] simultaneously satisfies stability and gives \[
R(c)<\log2-\tfrac12ca^2\mu_\beta^2.
\tag{F.8.3}
\] Thus a positive-step predictor improves population risk by a fixed
amount, yet every positive stopping time still has the uniform
calibration floor (F.6.4). In particular, a rule allowed to output time
zero that achieves a fixed positive risk improvement with probability
tending to one must output a positive time with probability tending to
one, and therefore cannot attain asymptotic exact calibration. The
result is not merely a pathology of steps that never learn.

There is also an explicit common regime with \textbf{strict signed
overconfidence at the actual risk-optimal first step}. Retain (F.1.3)'s
width condition and choose \[
0<\beta<\frac{v}{a\sqrt{\gamma k_1}}.
\tag{F.8.4}
\] Let \(c_{\rm risk}\) be the unique minimizer of the first-step
Gaussian risk. The scalar theorem in Appendix E proves both strict risk
improvement and strictly positive average confidence minus accuracy at
this scale, and gives \[
0<c_{\rm risk}<\frac{\beta k}{V}
\le\frac{a^2\alpha\gamma\beta^2}{v^2}
<\frac{\alpha}{k_1}<\frac{2\alpha}{k_1}.
\tag{F.8.5}
\] Here \(\mu_\beta\le\beta/4\) and \(V\ge v^2/(4\alpha\gamma)\). Thus
choosing the fixed step \(c=c_{\rm risk}\) satisfies the same stability
assumption as the entire-trajectory theorem. For this single strongly
overparameterized model, the first iterate strictly improves risk and is
overconfident on average, while no positive stopping time attains exact
asymptotic calibration. As in the scalar theorem, this is an average
signed statement, not overconfidence at every score. Condition (F.8.4)
restricts this compatibility corollary only; the calibration floor and
population-risk divergence allow every fixed finite \(\beta\).

\hypertarget{f.9-abstract-extension-to-any-uniformly-positive-separability-margin}{%
\subsection{F.9 Abstract extension to any uniformly positive
separability
margin}\label{f.9-abstract-extension-to-any-uniformly-positive-separability-margin}}

The row-Gram condition is a concrete sufficient certificate, not a
logical necessity for the late-time argument. Suppose instead that, with
probability tending to one:

\begin{enumerate}
\def\labelenumi{\arabic{enumi}.}
\tightlist
\item
  The fresh-input feature conditions (F.2.1) hold with common constants.
\item
  \(\|H\|_{\mathrm{op}}/\sqrt n\le D_0\).
\item
  A unit vector \(u_d\) has
  \(\min_i s_i h_i^\top u_d\ge\kappa_{\rm sep}>0\), with a deterministic
  constant \(\kappa_{\rm sep}\).
\item
  The GD steps satisfy \(0<\eta_-\le\eta_d\le4/D_0^2\).
\end{enumerate}

The comparator \(b_R=(R/\kappa_{\rm sep})u_d\) and (F.3.4) give \[
\mathcal L_n(b^t)
\le\frac1{t+1}+\frac{\log^2(t+1)}{2\eta_-\kappa_{\rm sep}^2t}.
\tag{F.9.1}
\] Equation (F.4.1) holds with \(D_0\) in place of \(D\). The remainder
of Sections F.4--7 therefore applies verbatim. If the same fixed-time
positive calibration theorem is available, all positive stopping times
again have a common calibration floor and the population-risk optimum is
in a deterministic finite window.

Consequently a separately proved exact-model max-margin limit
\(\kappa_d\to\kappa_*>0\) would extend this module to every parameter
point covered by that positive-margin theorem. This paragraph is an
implication, not an assertion that a particular external universality
theorem has already verified the premise for every aspect ratio. At a
zero-margin phase boundary, the deterministic crossing time may diverge
and this proof does not close the stopping-time quantifier.

\clearpage

\hypertarget{references}{%
\section{References}\label{references}}

\begin{enumerate}
\def\labelenumi{\arabic{enumi}.}
\tightlist
\item
  Y. Bai, S. Mei, H. Wang and C. Xiong. \emph{Don't Just Blame
  Over-parametrization for Over-confidence}. ICML, PMLR 139, 566--576
  (2021). \href{https://proceedings.mlr.press/v139/bai21c.html}{Primary
  paper}.
\item
  L. Clarté, B. Loureiro, F. Krzakala and L. Zdeborová. \emph{On
  double-descent in uncertainty quantification in overparametrized
  models}. AISTATS, PMLR 206 (2023).
  \href{https://proceedings.mlr.press/v206/clarte23a.html}{Primary
  paper}.
\item
  C. Louart, Z. Liao and R. Couillet. \emph{A random matrix approach to
  neural networks}. Annals of Applied Probability 28(2), 1190--1248
  (2018).
  \href{https://doi.org/10.1214/17-AAP1328}{doi:10.1214/17-AAP1328}.
\item
  Z. Fan. \emph{Approximate Message Passing algorithms for rotationally
  invariant matrices}. Annals of Statistics 50(1), 197--224 (2022).
  \href{https://doi.org/10.1214/21-AOS2101}{doi:10.1214/21-AOS2101}.
\item
  T. Wang, X. Zhong and Z. Fan. \emph{Universality of Approximate
  Message Passing algorithms and tensor networks}. Annals of Applied
  Probability 34(4), 3943--3994 (2024).
  \href{https://arxiv.org/html/2206.13037v5}{Versioned full text},
  especially Proposition D.1(b2) and Theorem 2.22.
\item
  Z. Ji, J. Li and M. Telgarsky. \emph{Early-stopped neural networks are
  consistent}. NeurIPS 34 (2021).
  \href{https://proceedings.neurips.cc/paper/2021/file/0e1ebad68af7f0ae4830b7ac92bc3c6f-Paper.pdf}{Primary
  paper}.
\item
  J. Wu, P. Bartlett, M. Telgarsky and B. Yu. \emph{Benefits of Early
  Stopping in Gradient Descent for Overparameterized Logistic
  Regression}. ICML, PMLR 267, 67081--67110 (2025).
  \href{https://proceedings.mlr.press/v267/wu25b.html}{Primary paper}.
\item
  E. Berta, D. Holzmüller, M. I. Jordan and F. Bach. \emph{Rethinking
  Early Stopping: Refine, Then Calibrate}. arXiv:2501.19195v2 (2025).
  \href{https://arxiv.org/html/2501.19195v2}{Primary manuscript}.
\item
  Y. Shmalo. \emph{Math Brain: public research repository} (2026).
  \url{https://github.com/yspennstate/math-brain-public}.
\item
  Y. Shmalo. \emph{Calibration along gradient descent in Gaussian tanh
  random-feature classifiers}. Public manuscript and reproducibility
  package, 31 October 2026.
  \href{https://github.com/yspennstate/math-brain-public/tree/main/papers/three-targets-2026-10-31/T357}{Public paper archive};
  \href{https://github.com/yspennstate/math-brain-public/blob/main/papers/three-targets-2026-10-31/T357/T357_Calibration_Along_Gradient_Descent.pdf}{PDF in Math Brain}.
\end{enumerate}

\section*{Public repository and source availability}
The complete manuscript, six detailed proof appendices, and editable
sources are archived in the public Math Brain repository [9]. The
paper-specific archive and PDF are given in [10], together with the
release's contribution history and verification record.

\section*{Acknowledgments and funding}\label{acknowledgments-funding-and-declarations}
\addcontentsline{toc}{section}{Acknowledgments and funding}

The research presented in this
paper was supported by the European Research Council (ERC) under the
European Union's Horizon Europe research and innovation programme (grant
agreement No.~101041711), by the Simons Foundation as part of the
Collaboration on the Mathematical and Scientific Foundations of Deep
Learning, by Heights Labs, by Convex Nexus Capital, by the Israel
Science Foundation (grant number 2258/19), and by the Israel Science
Foundation (ISF Grant 4101/25).

\section*{Declaration of competing interest}
\addcontentsline{toc}{section}{Declaration of competing interest}

The author holds equity in
Heights Labs and in Convex Nexus Capital and has a management role at
Convex Nexus Capital. These entities are included in the funding
declaration above.

\section*{Declaration of generative AI and AI-assisted technologies}
\addcontentsline{toc}{section}{Declaration of generative AI and AI-assisted technologies}

The mathematical arguments,
literature comparisons, verification, and writing were produced by AI
systems under the author's direction. This manuscript consolidates the
finite-time proof package of 24 October 2026 and its strongly
overparameterized continuation. The author selected the problem and
commissioned and directed the work; the author did not carry out the
mathematics or the writing. Responsibility for submission or publication
rests with the author. Internal AI-assisted mathematical reviews and
normalization checks are not external peer review or proof-assistant
formalization. The accompanying public source archive records the contribution history
and checksums of the released files.

\end{document}
